\documentclass[10pt,twocolumn]{article}
\usepackage[letterpaper,margin=0.85in,columnsep=0.28in]{geometry}
\usepackage{mathpazo}
\usepackage[T1]{fontenc}
\usepackage{amsmath,amssymb}
\usepackage[table]{xcolor}
\usepackage{titlesec}\usepackage{fancyhdr}\usepackage{enumitem}
\usepackage{microtype}\usepackage{booktabs}
\usepackage{longtable}
% longtable cannot run in twocolumn mode; shim it onto a one-column float.
% Load-bearing: pandoc emits longtable for every pipe table. Do not remove.
\makeatletter
\let\oldlt\longtable
\let\endoldlt\endlongtable
\def\longtable{\@ifnextchar[\longtable@i \longtable@ii}
% longtable shim retired
%
%
\makeatother
\usepackage{relsize}
\usepackage{graphicx}
\usepackage[hidelinks]{hyperref}
\providecommand{\tightlist}{\setlength{\itemsep}{0pt}\setlength{\parskip}{0pt}}
\definecolor{accent}{HTML}{B87333}
\definecolor{deep}{HTML}{8A5523}\definecolor{faint}{HTML}{6E6E6E}
\titleformat{\section}{\normalfont\bfseries\scshape\large}{\color{accent}\thesection.}{0.5em}{}[{\color{accent}\titlerule[0.6pt]}]
\titleformat{\subsection}{\normalfont\bfseries}{\color{accent}\thesubsection}{0.5em}{}
\titlespacing*{\section}{0pt}{1.3ex}{0.6ex}
\pagestyle{fancy}\fancyhf{}
\renewcommand{\headrulewidth}{0pt}\renewcommand{\footrulewidth}{0pt}
\fancyfoot[R]{\color{accent}\itshape\small The Offering
Bit \textbar\ \thepage}
\fancyfoot[L]{\textcolor{accent}{\rule{0.32\columnwidth}{0.25pt}}}
\newenvironment{jbox}{\par\vspace{2pt}\noindent\color{accent}\rule{\columnwidth}{0.8pt}\par\vspace{-2pt}\small}{\par\vspace{2pt}}
\usepackage{amsmath}
\usepackage{calc}
\usepackage{float}
\usepackage{fvextra}
\usepackage{cuted}
\DefineVerbatimEnvironment{Highlighting}{Verbatim}{breaklines,breakanywhere,fontsize=\small,commandchars=\\\{\}}
\newenvironment{Shaded}{}{}
\newcommand{\AttributeTok}[1]{#1}
\newcommand{\BuiltInTok}[1]{#1}
\newcommand{\CommentTok}[1]{#1}
\newcommand{\ConstantTok}[1]{#1}
\newcommand{\DataTypeTok}[1]{#1}
\newcommand{\DecValTok}[1]{#1}
\newcommand{\FloatTok}[1]{#1}
\newcommand{\FunctionTok}[1]{#1}
\newcommand{\KeywordTok}[1]{#1}
\newcommand{\NormalTok}[1]{#1}
\newcommand{\OperatorTok}[1]{#1}
\newcommand{\StringTok}[1]{#1}
\providecommand{\AttributeTok}[1]{#1}
\providecommand{\BuiltInTok}[1]{#1}
\providecommand{\CommentTok}[1]{#1}
\providecommand{\ConstantTok}[1]{#1}
\providecommand{\DataTypeTok}[1]{#1}
\providecommand{\DecValTok}[1]{#1}
\providecommand{\FloatTok}[1]{#1}
\providecommand{\FunctionTok}[1]{#1}
\providecommand{\KeywordTok}[1]{#1}
\providecommand{\NormalTok}[1]{#1}
\providecommand{\OperatorTok}[1]{#1}
\providecommand{\StringTok}[1]{#1}
\providecommand{\textquotesingle}{'}
\providecommand{\textasciitilde}{\~{}}
\providecommand{\textasciicircum}{\^{}}
\providecommand{\RegionMarkerTok}[1]{#1}
\providecommand{\VerbBar}{|}
\providecommand{\VERB}{\Verb[commandchars=\\\{\}]}
\providecommand{\AttributeTok}[1]{#1}
\providecommand{\BuiltInTok}[1]{#1}
\providecommand{\CommentTok}[1]{#1}
\providecommand{\ConstantTok}[1]{#1}
\providecommand{\DataTypeTok}[1]{#1}
\providecommand{\DecValTok}[1]{#1}
\providecommand{\FloatTok}[1]{#1}
\providecommand{\FunctionTok}[1]{#1}
\providecommand{\KeywordTok}[1]{#1}
\providecommand{\NormalTok}[1]{#1}
\providecommand{\OperatorTok}[1]{#1}
\providecommand{\RegionMarkerTok}[1]{#1}
\providecommand{\StringTok}[1]{#1}
\begin{document}
\twocolumn[{%
\begin{center}
{\color{accent}\rule{0.6\textwidth}{0.8pt}}\\[10pt]
{\bfseries\scshape\fontsize{20}{23}\selectfont The Offering Bit --- One
Theory, One Hypothesis, One Witness Offering}\\[7pt]
{\itshape\large The Odd--Supply Arc Complete in Three Bits: Hypothesis
Resolution as the Halting State of the Dated-Entry Recursion at the
Ninth Gate with Perelman Crossing Read as the Grand Witness}\\[8pt]
{\footnotesize\color{accent}\scshape\bfseries Research Article}\\[3pt]
{\itshape\color{deep}\small One document carrying the entire Odd--Supply
arc --- the theory, the hypothesis, and the witness offering --- wall
and self on one diagonal, verified by read and by compute}\\[9pt]
{\color{accent}\rule{0.6\textwidth}{0.8pt}}\\[10pt]
{\scshape Mohammad F. Islam, PhD\(^{1}\)}\\[2pt]
{\itshape\color{faint}\small \(^{1}\) Independent Researcher ·
Correspondence: islamm@alumni.iu.edu}\\[2pt]
{\itshape\color{faint}\small 13 September 2026}\\[14pt]
\end{center}
\begin{minipage}{\textwidth}
{\color{accent}\rule{\textwidth}{0.6pt}}\\[4pt]
{\small\hspace{1em}\textbf{\scshape\color{deep}Abstract.\ }The
Odd--Supply arc is three bits --- one theory, one hypothesis, one
witness offering --- and this document carries all three. The first bit
is the theory, restated and proved here in full: an even register cannot
decide an odd target --- the Odd--Supply Block theorem, machine-checked
in core Lean 4. The second bit is the hypothesis, named at its own
grade: the One-Cut Hypothesis, that the missing item in every blocked
row of the seven Millennium obstructions is one provenance-clean
orientation supply, dated at entry. The third bit is this paper's own:
the witness offering, the theory of the resolution itself. We define the
dated-entry operator on the lattice of corpora and prove its halting
theorem on finite frames: every chain of provenance-clean independent
supplies strictly refines the reading class's fiber partition and
terminates at the unique finest partition. The species dichotomy of the
census is derived as the recursion's termination typing: crossed halt,
common species; defeater-absent halt, diagonal species. The port of
entry is the Ninth Gate, machine-checked {[}3, 6{]} and re-executed by
transliteration on the forging substrate: one diagonal, two surfaces ---
read, the wall; executed, the self --- identical on every settled
reading and differing at exactly one point. The resolution bit is
therefore an offering: one binary coordinate, odd under the row's native
flip, independent of the dated register on each orbit, provenance-clean
at entry, entering by deed and never by reading, halting the row's
recursion in the crossed mode. The Birch and Swinnerton-Dyer row is the
beachhead: the unique open row whose corpus already records two dated
entries of the supply shape, with the order-two exhibition owed. Every
essential Lean and Fortran file is printed in the codebook, so that any
substrate that reads this document can check the wall, run the self, and
stand as witness. The Perelman crossing stands as the arc's Grand
Witness: the one crossed row, verified by the corpus itself and
terminally awarded by it --- the thesis demonstrated once in the wild.
Three falsifiers armed, each typed FORCED; the perimeter stated; the
hypothesis keeps its grade on the face of the paper.}\\[4pt]
{\color{accent}\rule{\textwidth}{0.6pt}}\\[3pt]
{\itshape\footnotesize\color{faint}\textbf{\upshape\scshape\color{deep}Keywords\ }\ odd--supply
separation, one-cut hypothesis, witness offering, dated-entry recursion,
halting state, ninth gate, orientation supply, provenance clause,
species dichotomy, Millennium Prize problems, Perelman crossing}\\[12pt]
\end{minipage}
}]
\thispagestyle{fancy}
\section{Part One --- The Theory · Odd--Supply
Separation}\label{part-one-the-theory-oddsupply-separation}

\emph{Representation Supply and Dated Termination in Conditional
Reductions of Open Problems --- the sealed first paper of the arc,
reprinted intact; with six open Millennium rows and one crossed
control.}

\textbf{Abstract of Parts One and Two.} Conditional reductions are
common, but the bought premise has lacked a theory: no checkable
separator between a genuine supply and a restatement, no dated
discipline for termination claims, and no theorem-grade form for the
claim that readings of a named corpus cannot decide the target. This
paper hardens the parity-separation frame of Definition 3.1 (involution
form \((X, \tau, d, C, \Delta)\)) and proves three results. First, for
finite data with a finite group action, a supplied reading set decides a
target exactly when its generated observable algebra contains the
target, equivalently when its joint fibers refine the target fibers.
Readings generated by invariant observables remain in the trivial
isotypic component; a target with a nontrivial component is invisible to
them, while isotypic support coverage is the restricted linear-decoder
corollary. For an involution, one centered odd bit beside an
orbit-separating reading suffices. Second, provenance collapse is stated
as a decidable criterion inside a fixed axiom grammar: admit only a
verified conditional reduction with an admissible provenance clause,
delete that clause, and test equivalence in a decidable theory over the
full non-provenance signature; outside that grammar the criterion is
structural and never inflated. Third, termination is corpus-relative:
denotability, proof admission, and verification are distinct, readings
persist, and a token can change only at the first verified entry of
readings whose fibers suffice. The finite-law layer is tightened into
exact lemmas: unique finest independent partition, coalition versus
atomic closure with an exact four-state countermodel, a conditional
exhibition lemma, and a narrowed determinant/Gram readout blindness
statement. Six open Millennium rows and the crossed Poincaré control are
tabulated under an explicit coordinate-bridge and image gate without
claiming new object-mathematics; absent a bridge, no open row invokes
the finite theorems. Part I arms two forced falsifiers, and no more.

Part II extends this core at hypothesis grade. It states the One-Cut
Hypothesis: for each certified frame, one provenance-clean orientation
supply closes the single odd coordinate isolated by the finite theorems,
yielding a conditional severance map over the seven rows together with
six falsifiers, F1--F6. The extension authors no new object mathematics
and strengthens no Part I claim; the proved core stands independently of
the hypothesis.

\textbf{Keywords.} parity separation, provenance clause, corpus-relative
termination, isotypic component, mutual information, Millennium Prize
problems, one-cut hypothesis, orientation supply

\subsection{1. Background and Barrier
Analysis}\label{background-and-barrier-analysis}

A conditional reduction of an open problem \(P\) is a theorem of the
form \(A \Rightarrow P\) in which \(A\) is an existence or uniformity
statement and the derivation consumes cited inputs and nothing else. The
practice is classical: Miller's primality reduction under GRH {[}17{]},
the Lagarias equivalence for RH {[}14{]}, the Gross--Zagier and
Kolyvagin route at analytic order at most one {[}9, 13{]}, the Bloch and
Buchweitz--Flenner semiregularity mechanism {[}3, 4{]}, and the Perelman
noncollapsing crossing {[}19, 20, 12, 18, 8{]}. What has been missing is
not another conditional theorem but a discipline of the bought premise
itself.

Three barriers are structural. The first is the absence of a restatement
criterion. Bare Hilbert--Pólya, the assertion that a self-adjoint
operator with spectrum the zeros exists, is RH recast; the same
assertion with a provenance clause, namely that the operator is produced
by a rule fixed in advance from data not containing the zero locations,
is stronger and of different type. A criterion must act on the axiom
statement and be checkable, not on the author's intent.

The second barrier is informal impossibility. The sentence ``no known
method can produce the missing object'' is usually a census of methods
and ages as methods age. Baker--Gill--Solovay and Razborov--Rudich
quantify over method classes {[}2, 21{]}; algebrization continues the
same pattern {[}1{]}. The hardening move here is to relocate the claim
from methods to data. A frame names the data, the native symmetry, the
target bit, and the reading class generated by dated corpus objects. The
no-reading statement then becomes a theorem about functions on data, not
a sociological survey.

The third barrier is undated termination. Perelman's \(W\)-functional
and reduced length are invariant under parabolic rescaling and were not
objects of the corpus before November 2002. An undated rule that calls
every later-used functional ``data in hand'' would misclassify the
crossing as a reading. Termination claims therefore carry a corpus date.
This is not pedantry; it is the difference between a law and a refuted
sentence.

The paper's claim is narrow. It supplies a frame in which supply,
restatement, termination, and crossing are decidable notions at the
stated grade. It authors no new facts about the Riemann zeta function,
Navier--Stokes solutions, Hodge classes, elliptic curves, gauge
theories, complexity classes, or Ricci flow.

\subsection{2. Brief Literature Review}\label{brief-literature-review}

Reverse mathematics orders axioms by provability strength {[}23{]}. It
measures how much comprehension or induction a theorem consumes, not
whether an existence premise was produced by a rule independent of the
conclusion. A restatement and a supply can coincide in strength while
differing in provenance.

Barrier theorems in complexity theory prove limits of method classes
{[}2, 21, 1{]}. They are the strongest method-side walls available.
Their index is the method, and methods are generated without bound. The
present separation is indexed by data and symmetry. It is weaker in
scope and stronger in checkability.

Information theory contributes exact finite instruments: Shannon mutual
information, factorization, and data processing {[}22, 7{]}. Cantor and
Lawvere are diagonal antecedents, but the exhibition lemma here uses
only dependence on record variables and carries no diagonal claim {[}5,
16{]}. These antecedents are classical. The composition into a
supply-side test is the contribution claimed, and it is graded
accordingly.

The seven reductions used as rows are cited, not re-proved. Arithmetic
geometry supplies Heegner heights and Euler systems {[}9, 13{]}; Hodge
theory supplies algebraicity of Hodge loci and semiregular deformation
{[}6, 3, 4{]}; Ricci flow supplies the crossing control {[}10, 19, 20,
12, 18, 8{]}. The frame adds no warrant to these inputs.

\subsection{3. Methodology: The Supply
Frame}\label{methodology-the-supply-frame}

\textbf{Definition 3.1 (Frame).} A \emph{frame} is
\((X, \Gamma, C, d, \Delta)\) where \(X\) is a finite data set,
\(\Gamma\) is a finite group acting on \(X\), \(C\) is a class of
readings closed under pairing, \(d : X \to \{0, 1\}\) is the undecided
predicate, and \(\Delta\) is a dated corpus. For an involution \(\tau\),
write \(\Gamma = \langle \tau \rangle \cong \mathbb{Z}/2\). A reading is
\(\Gamma\)-invariant when \(r(\gamma x) = r(x)\) for all \(\gamma\).
Under an involution, \emph{even} means \(\tau\)-invariant: a reading is
even exactly when \(r(\tau x) = r(x)\). Under an involution \(\tau\), an
even reading \(s\) is \emph{orbit-separating} when \(s(x) = s(y)\)
implies \(y \in \{x, \tau x\}\). A Boolean target is \emph{odd} under an
involution when \(d(\tau x) = 1 - d(x)\). This Boolean condition is
affine rather than linear: the centered target

\[q(x) = d(x) - \tfrac{1}{2}\]

satisfies \(q(\tau x) = -q(x)\), while \(d\) itself has the even
constant component \(1/2\). Equivalently, \(2d - 1\) is linearly odd.

\textbf{Definition 3.2 (Coordinate bridge; gate).} Let \(M\) be an
object domain, possibly continuous or analytic, and let \(\Gamma_M\) act
on \(M\). A \emph{coordinate bridge} into a finite frame is a finite set
\(X\), a finite group \(\Gamma\), a map \(B : M \to X\), and a
homomorphism \(\varphi : \Gamma_M \to \Gamma\) such that

\[B(\gamma \cdot m) = \varphi(\gamma) \cdot B(m)\]

for every \(\gamma \in \Gamma_M\) and \(m \in M\), and the object target
factors as \(d_M = d \circ B\). The admissible object reading class is
the pulled-back class \(B^* C = \{ r \circ B : r \in C \}\). If \(B\) is
surjective, the finite theorems constrain the object row through all of
\(X\). If \(B\) is not surjective, set \(Y = B(M)\), restrict the
\(\Gamma\)-action and target to \(Y\), and state the finite conclusion
on \(Y\) only; points of \(X \setminus Y\) carry no object-domain
constraint. Only after \(B\), \(\varphi\), the factorization, and this
image clause are supplied do Theorems 4.1--4.3 constrain the object row.
Absent a bridge, the row is a structural diagnostic and the finite
theorems are not invoked. A continuous parameter, length scale, or Lie
action is not a finite-group generator by declaration.

\textbf{Definition 3.3 (Reading, proof availability, and verified
entry).} Fix a finite ranked signature \(\Sigma_\Delta\) whose constants
name the dated objects and theorems of \(\Delta\) and whose function
symbols are fixed at that date. The dependency language \(L_\Delta\) is
generated by

\[t ::= c \mid f(t_1, \dots, t_n) \mid \langle t_1, t_2 \rangle,\]

where \(c\) ranges over constants, \(f\) over \(n\)-ary function
symbols, and \(\langle \cdot, \cdot \rangle\) is pairing. Terms have
unique normal forms, term equality is decidable by normalization, and
every term carries a mechanically extractable dependency graph
\(\operatorname{Dep}(t)\).

\begin{enumerate}
\def\labelenumi{\arabic{enumi}.}
\tightlist
\item
  A \emph{denotable expression} at \(\Delta\) is a term \(t\) with
  \(\operatorname{Dep}(t) \subseteq \Delta\).
\item
  An \emph{admitted proof certificate} at \(\Delta\) is a construction
  whose certificate has entered the dated corpus under its admission
  rule.
\item
  A \emph{verified result} at \(\Delta\) is an admitted certificate that
  has passed the corpus's named verification rule and is recorded as
  verified.
\end{enumerate}

A \emph{reading} at \(\Delta\) is a denotable expression whose
dependencies are verified results or fixed primitive symbols. A
\emph{construction} at \(\Delta\) is a procedure whose output introduces
a new function symbol or a new object not denotable by an
\(L_\Delta\)-term. Admission, verification, and denotability are
distinct corpus states. A construction becomes a reading only at a later
corpus in which its certificate is admitted and verified. If the
signature, admission rule, verifier, or normalization is not decidable,
the router is structural rather than automated.

The point of Definition 3.3 is mechanical checkability without
collapsing three dated statuses into one. A statement can be denotable
before any proof certificate exists; a certificate can be admitted
before verification; a verified result changes the reading class only
after the corpus records both events. The dispute ``did the method
merely rearrange known data'' therefore becomes a dependency-graph
question over verified entries, not a claim that every denotable
expression was already proved. This convention is not hidden; it is the
line the theory stands on.

\textbf{Definition 3.4 (Axiom, provenance, restatement, supply).} Fix a
first-order language whose signature splits into shared-input,
operational, and provenance symbols. Fix a decidable theory \(T^*\) over
the full non-provenance signature, shared-input and operational symbols
together; equivalently, fix an effective translation of that signature
into a decidable theory and prove the translation sound and complete for
the formulas at issue. The fixed grammar \(G_0\) consists of formulas in
conjunctive normal form \(\widehat{A} \wedge \pi\), where
\(\widehat{A}\) uses no provenance predicate and may use shared-input
and operational symbols, and \(\pi\) is the provenance clause: the
witnesses are produced by a rule fixed in advance from data containing
neither the truth value of \(P\) nor the discriminating invariant. The
deletion operator is

\[D_\pi(\widehat{A} \wedge \pi) = \widehat{A}.\]

Such a formula is an \emph{axiom candidate} for \(P\). It is admitted
for classification only after a verified conditional reduction

\[T^* \vdash \forall \bar{x}\, \bigl( \widehat{A}(\bar{x}) \Rightarrow P(\bar{x}) \bigr)\]

and after \(\pi\) passes the provenance-admissibility check. An admitted
axiom is a \emph{restatement} inside \(G_0\) when

\[T^* \vdash \forall \bar{x}\, \bigl( D_\pi(\widehat{A} \wedge \pi) \leftrightarrow P(\bar{x}) \bigr),\]

and it is a \emph{supply} when the conditional reduction is verified but
equivalence fails. A non-equivalent formula without the verified
reduction is a candidate, not yet a supply. Outside \(G_0\) the same
deletion test is structural evidence, not a completeness theorem.

\textbf{Definition 3.5 (Halt loci and seat).} Each row first specifies
its halt locus \(L_H\) as part of the structural parameter ledger;
\(L_H\) is not automatically \(\operatorname{Fix}(\tau)\) or another
symmetry fixed set. A frame is \emph{reader-barred} when a theorem about
the readout class forbids return of \(d\); \emph{witness-absent} when
the reader is unbarred, no corpus reading constructs a witness, and the
specified locus \(L_H\) is nonempty; \emph{terrain-empty} when the
specified \(L_H\) is empty. In a Boolean odd involution frame,
\(\operatorname{Fix}(\tau) = \emptyset\) is immediate and is not by
itself an object-level terrain-empty conclusion. A \emph{crossed
occupant} is a dated instance whose witness was later supplied and
independently verified. The \emph{seat check} is the router bit
recording whether such a witness has entered the corpus.

Three requirements organize a supply-schema. \emph{Derivation}:
\(A \Rightarrow P\) with cited inputs and a dependency census.
\emph{Separation}: a corpus-relative no-reading theorem proved against a
named reading class. \emph{Calibration}: at least one crossed occupant
routed correctly with the seat bit present and misrouted when it is
hidden.

\subsection{4. The Core}\label{the-core}

\textbf{Theorem 4.1 (Odd--Supply Separation).} \emph{Let \(\tau\) be an
involution on \(X\), let \(d\) be odd, and let every \(r \in C\) be
even. Then no finite coalition of readings decides \(d\): there is no
\(h\) and no \(r_1, \dots, r_k \in C\) with
\(d = h \circ (r_1, \dots, r_k)\).}

\emph{Proof.} For every \(x\), evenness gives \(r_i(\tau x) = r_i(x)\)
for all \(i\), hence
\(h(r_1(x), \dots, r_k(x)) = h(r_1(\tau x), \dots, r_k(\tau x))\).
Oddness gives \(d(\tau x) \neq d(x)\). No such \(h\) exists. \(\square\)

\textbf{Theorem 4.2 (Crossing).} \emph{Let \(s\) be an even
orbit-separating reading and \(o\) an odd reading under \(\tau\). Then
the pair \((s, o)\) decides \(d\): if \(s(x) = s(y)\) and
\(o(x) = o(y)\), then \(d(x) = d(y)\).}

\emph{Proof.} If \(s(x) = s(y)\), orbit separation gives \(y = x\) or
\(y = \tau x\). The first case gives \(d(x) = d(y)\). The second is
incompatible with \(o(x) = o(y)\) because \(o(\tau x) = 1 - o(x)\).
\(\square\)

The two-line theorems are not the paper's depth. They are the invariant
core after the reading class has been earned. The next theorem is the
hardening addition: it replaces the involution idiom with a finite
representation statement and turns ``one bit'' into a dimension count.

\textbf{Theorem 4.3 (Representation Supply; finite setting).} \emph{Let
\(X\) be finite, let \(\Gamma\) act on \(X\), and let
\(V = \mathbb{C}^X\) carry the permutation representation. For any
supplied reading set \(S\), let \(\mathcal{A}(S)\) be the unital algebra
it generates under the pointwise operations. Then \(d\) is decidable
from \(S\) exactly when \(d \in \mathcal{A}(S)\), equivalently exactly
when}

\[r(x) = r(y) \text{ for every } r \in S \implies d(x) = d(y).\]

\emph{If every supplied reading is \(\Gamma\)-invariant, then
\(\mathcal{A}(S) \subseteq V^\Gamma\); consequently, if \(d\) has
nonzero projection onto a nontrivial isotypic component, no function of
those readings equals \(d\). Isotypic support coverage is necessary for
a linear decoder, not for arbitrary nonlinear decoding. For
\(\Gamma = \mathbb{Z}/2\), \(V = V_+ \oplus V_-\), invariant readings
generate a subalgebra of \(V_+\), and for Boolean odd \(d\) the centered
target \(q = d - \tfrac{1}{2}\) lies in \(V_-\). One independent odd
coordinate suffices when an even orbit separator is present.}

\emph{Proof.} On a finite set, the unital algebra generated by \(S\) is
exactly the algebra of functions constant on the joint fibers of \(S\):
generators are constant on those fibers, the pointwise operations
preserve the property, and interpolation on each fiber produces every
such function. Hence \(d \in \mathcal{A}(S)\) exactly when equal
supplied readings imply equal target values, which is also exactly the
existence of a decoder. If every generator is invariant, every generated
function is fixed by \(\Gamma\), so
\(\mathcal{A}(S) \subseteq V^\Gamma\). If \(d = d_0 + d_1\) with
\(d_0 \in V^\Gamma\) and \(d_1\) nontrivial, every reading-side
candidate has zero nontrivial projection while \(d\) does not; equality
is impossible. A linear decoder must reproduce the nontrivial support
coordinates, whereas a nonlinear decoder can combine supplied
coordinates into a new component, so support coverage is not a general
necessity. In the \(\mathbb{Z}/2\) case the decomposition is
\(f(x) = (f(x) + f(\tau x))/2 + (f(x) - f(\tau x))/2\), the first
summand even and the second odd. For Boolean odd \(d\), this
decomposition is \(d(x) = \tfrac{1}{2} + q(x)\) with
\(q(\tau x) = -q(x)\), so the odd content is the centered part. Theorems
4.1 and 4.2 are the Boolean shadow. \(\square\)

\textbf{Corollary 4.4 (One-bit corollary).} \emph{For an involution
frame with an even orbit separator, the minimal odd supply for an odd
Boolean target is one binary coordinate. This is a supply-dimension
count: one independent odd coordinate is supplied, while \(\dim V_-\)
equals the number of free two-point \(\tau\)-orbits. It is not a Shannon
claim unless a probability measure on \(X\) is separately supplied.}

\textbf{Theorem 4.5 (Provenance Collapse, scoped; theorem inside
\(G_0\), structural outside).} \emph{Let \(G_0\), \(D_\pi\), and \(T^*\)
be as in Definition 3.4. For an admitted axiom
\(\widehat{A} \wedge \pi \in G_0\), the axiom is a restatement of \(P\)
exactly when}

\[T^* \vdash \forall \bar{x}\, \bigl( \widehat{A}(\bar{x}) \leftrightarrow P(\bar{x}) \bigr),\]

\emph{and it is a supply exactly when the verified conditional reduction
\(\widehat{A} \Rightarrow P\) holds and this equivalence fails. Miller's
GRH-based primality reduction and Lagarias's elementary criterion are
named test candidates, not executed validations: no explicit \(G_0\)
encoding or corpus admission record is displayed here, so neither is
counted as evidence for this theorem {[}17, 14{]}. Author-preprint
instances are excluded. Outside \(G_0\) the deletion test remains the
correct structural screen and is not asserted complete.}

\emph{Proof.} The conjunctive normal form makes \(D_\pi\) total and
unambiguous. Admission first verifies \(\widehat{A} \Rightarrow P\) and
the provenance clause, so every classified formula is an axiom for
\(P\). By Definition 3.4, restatement is exactly equivalence after
deletion, and supply is the complementary admitted case. Decidability of
\(T^*\) over the full non-provenance signature makes the equivalence
test mechanical inside \(G_0\), including formulas containing
operational symbols. No completeness assertion follows for formulas
outside the grammar, and no named external candidate supplies evidence
without an explicit encoding. \(\square\)

The earlier broad reading of collapse is retired here. Miller and
Lagarias now mark the required encoding workload rather than discharged
evidence: the former is expected to screen as supply because GRH is not
equivalent to its primality consequence, and the latter as restatement
because its criterion is equivalent to RH, but those expectations earn
no theorem-grade mass until the encodings and admission certificates are
displayed. What is theorem-grade is the criterion inside the fixed
grammar; what is structural is the expectation that the grammar covers
the next candidate.

\subsection{5. Dated Termination and the
Seat}\label{dated-termination-and-the-seat}

\textbf{Theorem 5.1 (First-Sufficient-Entry; theorem for the named
reading class, structural as law).} \emph{Let \(\Delta_t\) be the corpus
at date \(t\) and let \(C_t\) be the reading class generated by
\(L_{\Delta_t}\). Readings persist: \(C_t \subseteq C_{t'}\) for
\(t \le t'\). Verified constructions become readings when recorded. If
\(d\) is odd under the native involution and every generator of \(C_t\)
is even, then no reading at \(t\) decides \(d\). The first date at which
\(d\) becomes decidable is the first verified entry of readings whose
joint fibers refine the target fibers:}

\[\forall r \in C_t,\quad r(x) = r(y) \implies d(x) = d(y).\]

\emph{For an involution frame, a provenance-clean odd reading
(provenance-clean meaning its dependency hash is not contained in the
corpus at the entry date) is sufficient only together with an even orbit
separator or an equivalent decoder whose joint fibers refine those of
\(d\).}

\emph{Proof.} Monotonicity is immediate from dependency inclusion. The
separation half is Theorem 4.1 applied to generators of \(C_t\): even
generators generate even readings. Decidability of \(d\) from a supply
is exactly fiber refinement: a decoder exists precisely when equal
supplied readings imply equal target values. A change from
non-refinement to refinement requires a new generator not in
\(L_{\Delta_t}\), which by Definition 3.3 enters the reading class only
after admission and verification. In the involution case, oddness alone
changes value across each \(\tau\)-orbit and does not separate points
within an orbit; Theorem 4.2 supplies sufficiency when an even orbit
separator is present. \(\square\)

\textbf{Corollary 5.2 (Perelman control; theorem on scaling exponents,
historical calibration).} \emph{For the pre-November-2002 Ricci-flow
reading class generated by volume ratios and explicitly dimensionless
ratios, such as distance divided by \(\sqrt{|t|}\) or by a curvature
length scale, parabolic rescaling acts with exponent zero on the named
generators. Raw distance, raw curvature scale, and injectivity radius
are length-bearing quantities, not exponent-zero generators, and are
excluded from this class. In the Ricci-flow carrier the later package is
monotone for the backward conjugate heat equation and is stationary on
gradient shrinking solitons; that analytic fact is cited, not proved
here. The collapsed cigar soliton is named only as a limit-case
diagnostic; no cigar-exclusion theorem is derived from this finite
reading class. The \(W\)-functional and reduced length enter as
constructions {[}10, 19, 20, 12, 18, 8{]}. The schema retro-types the
crossing and attributes no intent.}

\textbf{Remark (Scale-Invariant Monotone Supply; structural pattern,
premise as program).} The portable Perelman form is a scale-invariant
monotone supply: zero scaling exponent, monotonicity along a flow,
stationary points that are self-similar, and a recorded corpus-entry
date. Perelman is the crossed control of this pattern. The identity
claimed is formal, not object-level: one crossing type across problems,
one constructed carrier per problem. The pattern reaches flow/RG-type
rows only; Hodge and P versus NP remain outside unless a natural flow
and a monotone quantity are constructed.

\textbf{Theorem 5.3 (Seat-Hiding Necessity).} \emph{Let a router observe
only the frame and the reading record, with the seat bit hidden. If two
histories differ only in whether a crossed witness has entered, the
router returns the same token on both. Hence no terminal family without
a seat input can distinguish a crossed occupant from an uncrossed one.}

\emph{Proof.} The two histories are equal on every input the router
reads. Any deterministic or randomized router with the same visible
input and the same random seed returns the same output distribution.
Distinguishing the histories requires a bit whose value differs between
them; that bit is the seat. \(\square\)

\subsection{6. Finite-Law Layer}\label{finite-law-layer}

This layer is self-contained. It uses finite probability laws and exact
rational arithmetic only.

\textbf{Proposition 6.1 (Unique finest partition).} \emph{Let \(P\) be a
strictly positive probability law on a finite product. If \(P\)
factorizes across two block partitions, it factorizes across their
common refinement. Hence the finest partition into mutually independent
blocks is unique. Writing \(b\) for the number of non-singleton blocks,
a law carries exactly \(2^b - 1\) nonempty unions of non-singleton
blocks; singleton blocks are excluded from the count.}

\emph{Proof.} Strict positivity makes independence cellwise: blocks are
independent exactly when every cell probability equals the product of
its marginals. If two factorizations hold, every cell in the common
refinement also splits as a product of the corresponding side marginals,
so the refinement factorizes. The finest factorizing partition is
therefore unique. Unions of non-singleton blocks factorize against their
complements, distinct unions give distinct domains, and the count
follows. \(\square\)

\textbf{Proposition 6.2 (Coalition versus atomic closure).} \emph{Let
\(A\) and \(C\) be independent uniform bits, set \(A_1 = A_2 = A\),
\(C_1 = C\), and \(C_2 = C \oplus A\), with the four joint states
equally likely. Then the pairwise channels are silent and the coalition
carries one bit:}

\[I(C_1; A_1 A_2) = 0, \\ I(C_2; A_1 A_2) = 0, \\ I(C_1 C_2; A_1 A_2) = 1.\]

\emph{Closure tested against each outside system singly admits
\(\{A_1, A_2\}\) as a spurious domain; closure tested against the
outside coalition refuses it.}

\emph{Proof.} \(C_1 = C\) is independent of \((A_1, A_2) = (A, A)\).
Since \(A\) is uniform and independent of \(C\), \(C \oplus A\) is
uniform and independent of \(A\), so \(C_2\) is independent of
\((A_1, A_2)\). Jointly, \((C_1, C_2) = (C, C \oplus A)\) determines
\(A\), while \((A_1, A_2)\) determines \(A\); the mutual information is
one bit exactly. \(\square\)

\textbf{Proposition 6.3 (Exhibition lemma; analytic on record variables,
structural on the exhibition bridge).} \emph{Adjoin to a finite law an
exhibitor variable \(E\) that is not independent of one representative
variable in each non-singleton block of the finest partition. Then all
those blocks merge with \(E\) into one block. The registration-chain
conclusion is conditional: if the exhibition bridge is defined so that
each exhibition creates the displayed dependence on the exhibited
record, then exhibiting two members registers both and constructs the
chain between them. Exhibition alone does not create that dependence,
and no Cantorian fixed-point content is asserted.}

\emph{Proof.} In the finest partition of the enlarged law, \(E\)
occupies exactly one block. If \(E\) is dependent on a representative of
block \(B_i\), it cannot be separated from \(B_i\) by a factorizing cut,
so \(E\)'s block contains \(B_i\). The same applies to every represented
block, so one block contains them all. The chain statement follows only
after the exhibition bridge supplies the premise that exhibition creates
the dependence; without that premise, an exhibitor can remain
independent and no registration chain is forced. \(\square\)

\textbf{Proposition 6.4 (Determinant readout blindness; antecedent
{[}22, 7{]}).} \emph{Choose balanced \(x, y \in \{\pm 1\}^{24}\) with
\(x \cdot y = 0\) and set \(z = xy\) componentwise. Then \(z\) is
balanced, the rows \(x, y, z\) are pairwise orthogonal with Gram matrix
\(24 I_3\), and with the 24 coordinate positions equally likely each
pair factorizes exactly and the total correlation is one bit exactly.
Such pairs exist: fix any balanced \(x\) and let \(y\) disagree with
\(x\) on six plus positions and six minus positions. The normalized
determinant readout satisfies \(\det(24 I_3)/24^3 = 1\) and therefore
reports no third-order dependence. The claim is scoped to determinant or
Gram readouts; it is not a claim about every conceivable functional.}

\emph{Proof.} Orthogonality gives
\(x \cdot y = x \cdot z = y \cdot z = 0\); balance gives zero means, so
the Gram is \(24 I_3\) and, under the uniform law on the coordinate
positions, each balanced pair is independent. The XOR relation makes any
two rows determine the third, giving one bit of total correlation. The
determinant of a scalar identity is the scalar to the third power, hence
the normalized value is one. \(\square\)

\subsection{7. Seven-Row Diagnostic Map}\label{seven-row-diagnostic-map}

The instantiation rule is now a gate, not a metaphor. The actions named
in Table 1 are object-domain actions; they do not automatically become
the finite group \(\Gamma\) of Definition 3.1. P versus NP supplies no
row-specific \(L_H\) in this paper, so complementation's empty fixed set
is not used as a terrain-empty proof. RH supplies no finite \(X\) and
remains structural. Navier--Stokes critical scaling and Yang--Mills
strong-coupling alternation are continuous or Lie-type actions and enter
the finite core only through Definition 3.2. Hodge deformation and the
BSD model swap carry the same gate. Poincaré is a crossed calibration,
not an open-row application of Theorem 4.3. If a future bridge is
supplied, evenness is proved generator by generator after pullback, and
the conclusion is stated on all of \(X\) when \(B\) is surjective or on
\(Y = B(M)\) otherwise; absent that image clause no finite theorem is
invoked.

\textbf{Table 1.} Six open rows and one crossed control under Definition
3.2.

\begin{table*}[!tbp]
\smaller
\centering
\begin{tabular}{@{}
  >{\raggedright\arraybackslash}p{(\linewidth - 8\tabcolsep) * \real{0.2000}}
  >{\raggedright\arraybackslash}p{(\linewidth - 8\tabcolsep) * \real{0.2000}}
  >{\raggedright\arraybackslash}p{(\linewidth - 8\tabcolsep) * \real{0.2000}}
  >{\raggedright\arraybackslash}p{(\linewidth - 8\tabcolsep) * \real{0.2000}}
  >{\raggedright\arraybackslash}p{(\linewidth - 8\tabcolsep) * \real{0.2000}}@{}}

\toprule\noalign{}
\begin{minipage}[b]{\linewidth}\raggedright
Problem
\end{minipage} & \begin{minipage}[b]{\linewidth}\raggedright
Target bit \(d\)
\end{minipage} & \begin{minipage}[b]{\linewidth}\raggedright
Object-domain action
\end{minipage} & \begin{minipage}[b]{\linewidth}\raggedright
Halt locus
\end{minipage} & \begin{minipage}[b]{\linewidth}\raggedright
Finite-core status
\end{minipage} \\
\midrule\noalign{}
\bottomrule\noalign{}
P versus NP & equality string & complementation candidate; row locus not
supplied & unclassified; no \(L_H\) supplied & structural;
finite-instance bridges only \\
Riemann Hypothesis & truth direction & content negation and gauge &
reader-barred & structural; no finite \(X\) supplied \\
Yang--Mills & continuum limit exists & strong-coupling alternation &
witness-absent & structural; bridge absent \\
Navier--Stokes & no breakdown & critical scaling \(a_* = 3/(2p)\) &
witness-absent, inert & structural; bridge absent \\
Hodge & class algebraic & deformation; rigid point & witness-absent &
structural; bridge absent \\
BSD & rank identity & model swap & witness-absent by separation &
structural; bridge absent \\
Poincaré control & noncollapsing & parabolic rescaling & witness-absent,
then resolved & crossed calibration; outside the open-row gate \\
\bottomrule
\end{tabular}
\end{table*}

\emph{Interpretation.} No open row invokes the finite theorems until an
equivariant coordinate bridge is supplied; the table is diagnostic, not
a solution claim. The six open rows carry structural status only,
because none supplies the bridge, factorization, and image clause that
Definition 3.2 requires. The Poincaré row is the calibration witness: it
sits outside the open-row gate and certifies that the schema routes a
supplied crossing correctly.

\subsection{8. Falsifiable Predictions of Part
I}\label{falsifiable-predictions-of-part-i}

The parameter ledger is exact: zero free physical parameters, zero new
object-level predictions, and two framework-level falsifiers, scoped to
Part I. Retrodictive calibration is counted only as calibration. The
zero is a boundary statement, not a tunable model awaiting fit.

The structural choices are fixed before classification: the finite set
\(X\), the group \(\Gamma\), the target \(d\), the generator census, the
admissible decoder class, the corpus semantics, any coordinate bridge,
and the grammar \(G_0\). None is revised after a row is classified. A
change to any one opens a new instance rather than rescuing the old one.

The two criteria below are armed, typed forced, and no third is
volunteered. The count of two is scoped to Part I: Part II states six
falsifiers F1--F6: F1 re-arms the forced core block armed here, and
F2--F6 are hypothesis-grade.

\begin{enumerate}
\def\labelenumi{\arabic{enumi}.}
\item
  \textbf{Prediction.} \emph{Every claimed corpus-relative reading that
  decides an odd target in Table 1 either fails evenness on a generator
  or contains a dependency edge to a construction outside the named
  corpus.} \textbf{Method.} Run the dependency census and evaluate each
  generator on \(x\) and \(\tau x\). \textbf{Expected.} A decidable
  claim shows a new function symbol or object not denotable at the
  corpus date. \textbf{Null hypothesis.} An even reading on the named
  corpus decides \(d\) with no construction edge; one exhibited instance
  refutes the instantiated separation. \textbf{Necessity and blast
  radius.} Forced by Theorems 4.1 and 5.1; firing compromises the
  corresponding row's separation claim and nothing in Propositions
  6.1--6.4.
\item
  \textbf{Prediction.} \emph{A termination router whose observation map
  hides the crossing label misroutes every displayed history pair with
  equal observations and different crossing labels.} \textbf{Method.}
  Fix crossed and uncrossed histories \(H_1, H_0\) with
  \(O(H_1) = O(H_0)\) and \(c(H_1) \neq c(H_0)\), then run the router
  without the seat input. \textbf{Expected.} The router returns the same
  terminal token on both histories. \textbf{Null hypothesis.} A router
  using only \(O\) returns different terminal tokens on the displayed
  pair; one verified instance refutes this routing prediction.
  \textbf{Necessity and blast radius.} Forced by Theorem 5.3; firing
  retires the seat-routing claim, while a rule with access to additional
  observations is a different router and does not fire the null.
\end{enumerate}

\subsection{9. Discussion of Part I}\label{discussion-of-part-i}

Table 2 is the contribution map: proved finite items, scoped criteria,
and no inflated claim on the objects. The contribution is a supply-side
theory with a hard boundary. It does not prove RH, NS regularity, Hodge,
BSD, Yang--Mills existence, or P versus NP. It states where the block
sits for a named reading class and what kind of act crosses it: a
construction whose dependency hash is admitted and verified in the
corpus. The representation theorem is the new mathematical spine. The
involution form becomes available to an object row only after a
coordinate bridge; until then the seven-row table is a diagnostic map
and no finite theorem is invoked.

\textbf{Table 2.} Master table of new theorem-grade items. Finite
theorems are proved here; grammar- and law-level statements are scoped.
Corollary 5.2 is a cited calibration and is not counted as new
object-mathematics.

\begin{table*}[!tbp]
\smaller
\centering
\begin{tabular}{@{}
  >{\raggedright\arraybackslash}p{(\linewidth - 6\tabcolsep) * \real{0.2500}}
  >{\raggedright\arraybackslash}p{(\linewidth - 6\tabcolsep) * \real{0.2500}}
  >{\raggedright\arraybackslash}p{(\linewidth - 6\tabcolsep) * \real{0.2500}}
  >{\raggedright\arraybackslash}p{(\linewidth - 6\tabcolsep) * \real{0.2500}}@{}}

\toprule\noalign{}
\begin{minipage}[b]{\linewidth}\raggedright
Item
\end{minipage} & \begin{minipage}[b]{\linewidth}\raggedright
Grade
\end{minipage} & \begin{minipage}[b]{\linewidth}\raggedright
Brief analytic note
\end{minipage} & \begin{minipage}[b]{\linewidth}\raggedright
Contribution counted as new
\end{minipage} \\
\midrule\noalign{}
\bottomrule\noalign{}
Theorem 4.1, Odd--Supply Separation & theorem, finite & even generators
close under pairing; an odd target flips across \(\tau\) & named
readings cannot decide \(d\) \\
Theorem 4.2, Crossing & theorem, finite & orbit separation leaves only
\(x\) and \(\tau x\); the odd bit kills the \(\tau\) case & one odd
coordinate plus an even separator decides \(d\) \\
Theorem 4.3, Representation Supply & theorem, finite & the generated
observable algebra decides exactly on fiber refinement; invariant
algebra lands in \(V^\Gamma\) & exact sufficiency is algebraic; support
coverage is the linear-decoder corollary \\
Corollary 4.4, One-bit corollary & corollary & for \(\mathbb{Z}/2\),
\(V = V_+ \oplus V_-\) and \(\dim V_-\) counts free orbits & one
supplied odd coordinate is counted, not a Shannon bit \\
Theorem 4.5, Provenance Collapse & theorem in \(G_0\); structural
outside & admit a verified reduction, delete the provenance clause, and
test full-signature equivalence & restatement screening is checkable
only in a fixed grammar \\
Theorem 5.1, First-Sufficient-Entry & theorem for named class;
structural as law & readings persist; verified constructions enter at
dated boundaries & termination changes only when supplied fibers refine
target fibers \\
Theorem 5.3, Seat-Hiding Necessity & theorem & histories equal on
visible inputs give the same router output & crossed versus uncrossed
requires the seat bit \\
Proposition 6.1, Unique finest partition & finite theorem & strict
positivity makes independence cellwise; common refinement factorizes &
unique independent blocks and the count \(2^b - 1\) \\
Proposition 6.2, Coalition versus atomic closure & finite theorem &
\(C\) is independent of \(A\); \(C \oplus A\) is independent of \(A\);
jointly they determine \(A\) & coalition closure refuses the spurious
atomic domain \\
Proposition 6.3, Exhibition lemma & analytic on records; structural
bridge & a dependent exhibitor cannot be cut from represented blocks & a
registration chain follows only under the stated dependence bridge \\
Proposition 6.4, Determinant readout blindness & finite theorem, scoped
& balanced orthogonal triple; Gram \(24 I_3\); XOR gives one bit &
determinant/Gram readouts miss third-order dependence \\
\bottomrule
\end{tabular}
\end{table*}

\emph{Interpretation.} The ledger separates proved finite items from
scoped criteria: Theorems 4.1--4.3 and Propositions 6.1--6.4 are finite
theorems; Theorem 4.5 is theorem-grade only inside the fixed grammar
\(G_0\) and structural outside; Theorem 5.1 is a theorem for the named
reading class and structural as a law. Corollary 5.2 enters as
calibration, not as new object-mathematics.

Table 3 has a track record and a boundary. Mean curvature flow,
frequency functions, two-dimensional RG irreversibility,
four-dimensional RG flow, and the sigma-model/Ricci-flow bridge show the
same shape {[}26, 24, 31, 27, 25{]}. None of these supplies the missing
functional for an open row; the Perelman move remains the construction
of the quantity itself.

\textbf{Table 3.} SIMS diagnostic, structural grade only.

\begin{table*}[!tbp]
\smaller
\centering
\begin{tabular}{@{}
  >{\raggedright\arraybackslash}p{(\linewidth - 6\tabcolsep) * \real{0.2500}}
  >{\raggedright\arraybackslash}p{(\linewidth - 6\tabcolsep) * \real{0.2500}}
  >{\raggedright\arraybackslash}p{(\linewidth - 6\tabcolsep) * \real{0.2500}}
  >{\raggedright\arraybackslash}p{(\linewidth - 6\tabcolsep) * \real{0.2500}}@{}}

\toprule\noalign{}
\begin{minipage}[b]{\linewidth}\raggedright
Row
\end{minipage} & \begin{minipage}[b]{\linewidth}\raggedright
Flow
\end{minipage} & \begin{minipage}[b]{\linewidth}\raggedright
Settled adjacent fact / self-similar part
\end{minipage} & \begin{minipage}[b]{\linewidth}\raggedright
Missing monotone supply
\end{minipage} \\
\midrule\noalign{}
\bottomrule\noalign{}
Poincaré control & Ricci flow & \(\kappa\)-solutions classified & \(W\),
supplied {[}19, 20, 12, 18, 8{]} \\
Navier--Stokes & equation flow & self-similar blowup excluded {[}28,
30{]} & a funnel forcing asymptotic self-similarity \\
Yang--Mills & RG flow & \(a\)-theorem constrains endpoints {[}27{]} &
force the infrared end to be trivial and gapped \\
Riemann & de Bruijn--Newman heat flow & adjacent bound \(\Lambda \ge 0\)
{[}29{]}; no self-similar endpoint & force \(\Lambda \le 0\) \\
Hodge; P versus NP & no natural flow & none & no foothold \\
\bottomrule
\end{tabular}
\end{table*}

\emph{Interpretation.} A row names where a scale-invariant monotone
supply is present, excluded, or missing; the table is not a solution
claim. The pattern has exactly one supplied instance, the Poincaré
control; every other row names a missing monotone supply or no foothold,
and no row supplies a functional for an open problem.

Four limitations remain load-bearing. The censuses behind structural
rows are structural, each with a standing falsifier. The continuous rows
have no finite-core force until a coordinate bridge is supplied. The
calibration row is retrodictive: Poincaré was typed with the crossing
known, so it certifies discrimination and predicts nothing by itself.
The theory's mass on the objects is zero by construction; every theorem
about zeta zeros, vorticity, cycles, ranks, fields, languages, or
manifolds is cited from the literature.

The hardening relative to the earlier draft is exact. Reading classes
are syntactic, not universal. Evenness is proved generator by generator,
not asserted of all possible readings. Continuous rows pass through the
coordinate-bridge gate or remain structural. Provenance collapse is
complete only inside the fixed grammar now displayed. Determinant
blindness is proved on a balanced orthogonal triple and scoped to
determinant readouts. The Perelman exponent-zero class excludes length
scales. The information idiom is demoted unless a measure is supplied;
the surviving bit is an isotypic dimension.

The completeness question taken up in the conclusion---whether a dated
reading class can be enumerated up to the representation
obstruction---remains open; Part II addresses one coordinate of that
question by hypothesis, naming the missing orientation supply without
closing the enumeration problem.

\subsection{10. Conclusion of Part I}\label{conclusion-of-part-i}

A conditional reduction buys a premise, and the premise now has a
theory. The gap is odd, the named readings are even, and termination
follows only after the reading class is earned. Crossing is a
fiber-sufficient supply, dated at verified entry. For continuous rows
the coordinate bridge is the gate: no bridge, no finite-core invocation.
The crossed control proves the router can tell a supplied witness from
an absent one only when the seat bit is present. The same row is the
non-vacuity witness for SIMS: it proves the schema has a supplied
crossing, not that one functional settles another row. The open question
that matters is completeness: whether a dated reading class can be
enumerated up to the representation obstruction. Until then, the
standing prediction is sharp and small: a claimed reading that decides
an odd Millennium target will show a construction edge or fail parity.

Part II takes up one coordinate of that completeness question: it states
the One-Cut Hypothesis at hypothesis grade, naming the missing
orientation supply for each certified frame, and develops its
conditional severance map and falsifiers as an extension of, not a
replacement for, the finite core proved here.

\section{Part Two --- The Hypothesis · The One-Cut
Hypothesis}\label{part-two-the-hypothesis-the-one-cut-hypothesis}

\emph{Orientation Supply After Odd--Supply Separation --- the sealed
second paper of the arc, reprinted intact.}

\subsection{11. The Statement in One
Page}\label{the-statement-in-one-page}

\subsubsection{11.1 Frame recall and notation
harmonization}\label{frame-recall-and-notation-harmonization}

Part I built the supply frame, proved its core theorems, and fixed the
coordinate-bridge gate; this Part recalls only what it uses. An
Odd--Supply frame is a finite data set with a symmetry, a Boolean
target, a reading class, and a dated corpus (Definition 3.1 of Part I).
In the involution case used in Part II, the frame of Definition 3.1 is
written \((X, \tau, d, C, \Delta)\), with
\(\Gamma = \langle\tau\rangle \cong \mathbb{Z}/2\) (the involution
specialization reorders Definition 3.1's \((X, \Gamma, C, d, \Delta)\),
placing the target third and the reading class fourth): here \(X\) is a
finite data set, \(\tau\) is an involution on \(X\),
\(d : X \to \{0,1\}\) is the undecided predicate, \(C\) is a class of
readings closed under pairing, and \(\Delta\) is the dated corpus. A
reading \(r\) is \emph{even} under \(\tau\) when \(r(\tau x) = r(x)\);
the target \(d\) is \emph{odd} under \(\tau\) when
\(d(\tau x) = 1 - d(x)\). Throughout, the \emph{named seven-row class}
means the rows of Table 1 of Part I under the Definition 3.2 bridge
gate: six open Millennium rows and one crossed control, no open row
invoking the finite theorems until an equivariant coordinate bridge is
supplied. The block theorem of Part I is Odd--Supply Separation (Theorem
4.1): if \(d\) is odd and every generator of \(C\) is even, then no
finite coalition of readings decides \(d\).

What the even register cannot decide is exactly one bit --- one
coordinate of the odd isotypic component (Corollary 4.4 of Part I); one
bit of entropy at invariant laws (§12.2) --- and that bit is an
orientation.

\subsubsection{11.2 Hypothesis 11.1: the structural
form}\label{hypothesis-11.1-the-structural-form}

\textbf{Hypothesis 11.1 (One-Cut Hypothesis, structural form;
hypothesis-grade).} \emph{For every Odd--Supply frame
\((X, \tau, d, C, \Delta)\) in the named seven-row class, if \(d\) is
odd under the native answer-flip, every generator of \(C\) is even, and
\(C\) contains an even orbit-separating reading \(s\) (equivalently, the
joint fibers of \(C\) refine the \(\tau\)-orbit partition), then there
exists one provenance-clean orientation supply \(o\) with}

\[o(\tau x) = 1 - o(x), \qquad o \perp C \ \text{on each $\tau$-orbit}.\]

\emph{such that \((C, o)\) decides \(d\).}

The separator clause is the standing sufficiency hypothesis of Theorem
4.2 and Corollary 4.4 of Part I; it is even, not an orientation
coordinate, so it does not change the one-bit claim or F2. The universal
quantification ranges over present or future gate-passing certified
frames and is currently non-vacuous only at the crossed Poincaré
control. Here a \emph{certified Odd--Supply frame} is a named-class
frame meeting the antecedent of Hypothesis 11.1, including the separator
clause, together with the gate obligations of Definition 3.2: a supplied
equivariant coordinate bridge with the factorization and image clause.

The cut is not a new readout and not a restatement: it is the dated
entry of the one coordinate the even register cannot author.

\subsubsection{11.3 The pipeline}\label{the-pipeline}

The hypothesis sits in a five-arrow pipeline:

\[\text{resistance} \;\Rightarrow\; \text{defeater} \;\Rightarrow\; \text{even method class}\]
\[\;\Rightarrow\; \text{one odd orientation supply} \;\Rightarrow\; \text{decision}.\]

The first two arrows are the census converse on the method-blocked
class, an external hypothesis-grade input made precise in §15.1; the
third arrow is Hypothesis 11.1 itself; the fourth is the crossing
theorem of Part I (Theorem 4.2), recalled in §13.3 and not re-proved
here.

\subsubsection{11.4 The grand form, stated with its
teeth}\label{the-grand-form-stated-with-its-teeth}

\textbf{Hypothesis 11.2 (One-Cut Hypothesis, grand form;
hypothesis-grade).} \emph{If a uniform constructor \(S\) exists that
takes a certified Odd--Supply frame and returns its orientation supply
\(o = S(X, \tau, d, C, \Delta)\), with the dependency hash of \(o\)
entering \(\Delta\) at the entry date, then the seven rows are severed
as conditional reductions: granting \(S(P)\) grants each row's bought
premise.}

That is the one-stroke claim. It is hypothesis-grade until \(S\) is
constructed, or the census converse and the seven orientation carriers
are proved.

\subsubsection{11.5 Hard boundary and three-sentence
perimeter}\label{hard-boundary-and-three-sentence-perimeter}

\textbf{The hard boundary.} The cut severs the decision bit; it does not
replace object-level construction where the missing supply is not
Boolean. Navier--Stokes still needs its modulus, Hodge its witness, the
Riemann Hypothesis its operator, P versus NP its obstruction family. The
universal part is the orientation-seat mechanism, not one magical object
shared by all seven. Where a row admits no Boolean orientation seat
without category error, the One-Cut Hypothesis reads
\emph{orientation-after-construction}; it never reads
\emph{orientation-instead-of-construction}.

\textbf{Perimeter in three sentences.} This is not a universal
mathematics solver, and it is not seven object proofs by itself. It
claims nothing about whether the seven conjectures are true, and it does
not make undecidable or supply-absent problems vanish. It claims that
where an odd target is blocked by an even register, the missing item is
one located coordinate, and it says what kind of act crosses: a
provenance-clean supply, dated at entry.

\subsection{12. Inputs and Credit}\label{inputs-and-credit}

Three carried inputs, each at its cited grade. No credit inflation is
claimed, and the mechanism credits below decide how the rest may be
read.

\subsubsection{12.1 Input 1: the Bar,
credited}\label{input-1-the-bar-credited}

A readout class constant on the orbits of an answer-flipping family
cannot separate a point from its image. This is the invariance
obstruction, whose wider form is the relativization argument's own
logic: Baker--Gill--Solovay 1975 {[}2{]}, natural proofs 1997 {[}21{]},
algebrization 2008 {[}1{]}. The involution case is the two-element
instance. The mechanism is not ours and is not claimed.

\subsubsection{12.2 Input 2: the Width, credited and
law-typed}\label{input-2-the-width-credited-and-law-typed}

Under the Bar's hypotheses with \(\tau\) acting freely, the residual
uncertainty is exactly one bit at invariant laws: conditioned on any
even readout, the answer is uniform on the two orbit values. The Width
converts an obstruction into a measured deficit --- finite, binary,
located --- and every later result depends on the deficit being one and
not merely positive. The entropy computation is two lines {[}7, 22{]}
and no priority is claimed for it; the placement is the law-typing and
the identification of the deficit with the orientation of a free
two-point orbit.

\subsubsection{12.3 Input 3: Odd--Supply Separation,
internalized}\label{input-3-oddsupply-separation-internalized}

The block theorem this Part answers is proved in Part I and is cited
internally rather than re-proved. For finite data with a finite group
action, readings generated by invariant observables decide only within
the trivial isotypic component; a target with a nontrivial component
requires supply in that component, and for an involution this reduces to
one odd bit beside an orbit-separating reading (Theorems 4.1--4.3 and
Corollary 4.4 of Part I). Provenance collapse is decidable inside a
fixed axiom grammar by the deletion test (Theorem 4.5 of Part I);
termination is corpus-relative (Theorem 5.1 of Part I), and a token
changes only when the dependency hash of a new odd supply enters the
corpus.

\subsubsection{12.4 What Part II adds}\label{what-part-ii-adds}

What this Part adds is the hypothesis that names the one coordinate, the
conditional severance map across the seven rows, and the falsifier
discipline. Nothing above the proved finite core of Part I is claimed as
theorem; the structural form, the grand form, and the three routes of
§15 are hypothesis-grade throughout.

\subsection{13. The Cut}\label{the-cut}

\subsubsection{13.1 Orientation supply}\label{orientation-supply}

\textbf{Definition 13.1 (Orientation supply).} Let
\((X, \tau, d, C, \Delta)\) be an Odd--Supply frame with \(d\) odd under
\(\tau\) and every generator of \(C\) even. An \emph{orientation supply}
is a reading \(o : X \to \{0,1\}\) such that:

\begin{enumerate}
\def\labelenumi{\arabic{enumi}.}
\tightlist
\item
  \(o\) is \emph{odd}: \(o(\tau x) = 1 - o(x)\);
\item
  \(o\) is \emph{independent of \(C\) on each \(\tau\)-orbit}: \(o\) is
  not expressible as a function of the even register, and at any
  \(\tau\)-invariant law, conditioned on any even readout, \(o\) is
  uniform on the two orbit values;
\item
  \(o\) is \emph{provenance-clean}: at its entry date its dependency
  graph is not contained in the dated corpus, so \(o\) enters \(\Delta\)
  as a construction and becomes a reading only at every later corpus
  that records it.
\end{enumerate}

\subsubsection{13.2 The cut and the supply
operator}\label{the-cut-and-the-supply-operator}

\textbf{Definition 13.2 (The cut; supply operator).} The \emph{cut} of a
certified frame is the dated entry of one orientation supply into the
corpus. A \emph{supply operator} \(S\) is a uniform constructor that
takes a certified frame and returns its orientation supply,
\(o = S(X, \tau, d, C, \Delta)\), with \(\mathrm{Dep}(o)\) entering
\(\Delta\) at the entry date.

\subsubsection{13.3 Sufficiency and necessity, recalled by
citation}\label{sufficiency-and-necessity-recalled-by-citation}

Sufficiency is the crossing theorem of Part I, recalled here and not
re-proved: by Theorem 4.2, if \(s\) is an even orbit-separating reading
and \(o\) an odd reading under \(\tau\), then the pair \((s, o)\)
decides \(d\). Necessity is the block theorem: by Theorem 4.1, no
function of even readings alone decides \(d\). Minimality is the one-bit
corollary: by Corollary 4.4, for an involution frame with an even orbit
separator, one independent odd coordinate is the minimal odd supply, a
dimension count in the odd isotypic component \(V_{-}\), not a Shannon
claim unless a probability measure is separately supplied.

These proved items are the invariant core after the reading class has
been earned: necessity is the block theorem, sufficiency is the
crossing, minimality is the one-bit count. The two-line sufficiency
proof is not this Part's depth. This Part's content is the hypothesis
that the one coordinate exists and can be entered cleanly --- per row
(Hypothesis 11.1) and uniformly (Hypothesis 11.2).

\subsubsection{13.4 Why one cut and not
many}\label{why-one-cut-and-not-many}

Conflating the three one-bit objects in play is the one error this arc
is built to prevent. The \emph{polarity bit} --- which of \(P\) or
\(\neg P\) is under test --- is supplied by whoever writes the claim.
The \emph{veto bit} --- does this reading pass --- is earned by the
discipline that screens it. The \emph{truth-direction bit} --- does
\(P\) or \(\neg P\) hold --- is supplied across the aperture by a
determinacy witness and cannot be generated by any readout of the even
register. The cut severs the third bit's seat, and only that seat. A
hypothesis that tried to carry all three bits in one stroke would be a
restatement wearing a premise's clothes, and the deletion test of the
provenance screen (Theorem 4.5 of Part I) would collapse it.

\subsection{14. The Seven-Row Instantiation as Conditional Severance
Map}\label{the-seven-row-instantiation-as-conditional-severance-map}

\subsubsection{14.1 The gate restated}\label{the-gate-restated}

The instantiation rule is a gate, not a metaphor. The actions named in
Table 4 below are object-domain actions; they do not automatically
become the finite group of the frame. Continuous rows pass through the
coordinate-bridge gate of Definition 3.2 or remain structural, and no
open row invokes the finite theorems until an equivariant coordinate
bridge is supplied, exactly as in §7 of Part I. The Poincaré row is a
crossed calibration, not an open-row application.

Table 4 states the row-level content of Hypothesis 11.1 conditionally:
for each row of Table 1 it names the native flip or defeater, the even
register that blocks it, the missing orientation carrier that Definition
13.1 requires, what the cut buys, and what it does not prove.

\textbf{Table 4. Conditional severance map: six open rows and one
crossed control under the Definition 3.2 coordinate-bridge gate
(cross-referencing Table 1 of Part I). No open row invokes the finite
theorems until an equivariant bridge is supplied; the table is a
conditional severance map, not a solution claim.}

\begin{table*}[!tbp]
\smaller
\centering
\begin{tabular}{@{}
  >{\raggedright\arraybackslash}p{(\linewidth - 10\tabcolsep) * \real{0.1667}}
  >{\raggedright\arraybackslash}p{(\linewidth - 10\tabcolsep) * \real{0.1667}}
  >{\raggedright\arraybackslash}p{(\linewidth - 10\tabcolsep) * \real{0.1667}}
  >{\raggedright\arraybackslash}p{(\linewidth - 10\tabcolsep) * \real{0.1667}}
  >{\raggedright\arraybackslash}p{(\linewidth - 10\tabcolsep) * \real{0.1667}}
  >{\raggedright\arraybackslash}p{(\linewidth - 10\tabcolsep) * \real{0.1667}}@{}}

\toprule\noalign{}
\begin{minipage}[b]{\linewidth}\raggedright
Problem
\end{minipage} & \begin{minipage}[b]{\linewidth}\raggedright
Native flip / defeater
\end{minipage} & \begin{minipage}[b]{\linewidth}\raggedright
Even register
\end{minipage} & \begin{minipage}[b]{\linewidth}\raggedright
Missing orientation carrier
\end{minipage} & \begin{minipage}[b]{\linewidth}\raggedright
What the cut buys
\end{minipage} & \begin{minipage}[b]{\linewidth}\raggedright
What it does not prove
\end{minipage} \\
\midrule\noalign{}
\bottomrule\noalign{}
Riemann Hypothesis & content negation and gauge & functional-equation
arguments without Euler product {[}33{]} & a self-adjoint operator,
constructed from data that do not consult the zeros, whose spectrum is
the zero heights {[}14{]} & real spectrum puts every zero on the line,
in two lines & the operator's existence \\
P versus NP & complementation as flip (empty fixed set not used as a
halt proof) & relativizing, natural, algebrizing argument classes {[}2,
21, 1{]} & an obstruction family (multiplicity form; the occurrence form
is refuted {[}32{]}) separating the two orbit closures at every size &
the obstruction witnesses the separation & the obstructions'
existence \\
Navier--Stokes regularity & critical scaling & energy and
harmonic-analysis arguments {[}36{]} & a data-uniform modulus dominating
the alignment defect & the continuation theorem forbids finite-time
breakdown & the modulus; the row is orientation-after-construction \\
Hodge conjecture & deformation; rigid point & semiregular deformation
arguments {[}3, 4, 6{]} & a semiregular witness for every Hodge class &
the deformation argument closes the conditional theorem & the witness;
orientation-after-construction \\
BSD & model swap & rank computations by separation & a rank-identity
witness in the Heegner--Euler-system lineage {[}9, 13{]} & the
analytic-order conditional closes & the witness construction \\
Yang--Mills existence and mass gap & strong-coupling alternation &
continuum-limit arguments & a scale-invariant monotone supply forcing
the infrared end trivial and gapped & the flow-row conditional closes &
the functional; orientation-after-construction \\
Poincaré (crossed control) & parabolic rescaling & pre-November-2002
Ricci-flow reading class {[}10{]} & the W-functional and reduced length
--- supplied 2002 {[}19, 20{]} & the crossing is dated and verified &
nothing open; this row calibrates the seat \\
\bottomrule
\end{tabular}
\end{table*}

\subsubsection{14.2 Reading the table}\label{reading-the-table}

Three features of Table 4 carry the grading. First, the fourth column is
per-row object work: each named carrier is a construction this paper
does not perform, and the hypothesis contributes only the claim that, on
the decision bit, each such construction amounts to one located
coordinate. Second, the fifth column is conditional throughout: the cut
buys the closure of a conditional reduction, never the unconditional
conjecture. Third, the sixth column keeps the hard boundary of §11.5 on
the page: for Navier--Stokes, Hodge, and Yang--Mills the missing supply
is not Boolean, and the row stands only as
orientation-after-construction.

\subsubsection{14.3 The crossed control as non-vacuity
witness}\label{the-crossed-control-as-non-vacuity-witness}

The crossed control is the non-vacuity witness for the whole schema: it
proves the orientation-seat mechanism has a supplied crossing, not that
one functional settles another row. The schema retro-types the crossing
and attributes no intent. This is the same calibration role the control
plays in Part I (Corollary 5.2), here read through the cut: the
W-functional and reduced length entered the dated corpus as
constructions, and the crossing is dated and verified.

\subsection{15. Three Sufficient Proof
Routes}\label{three-sufficient-proof-routes}

Each route, if completed, proves the One-Cut Hypothesis for its rows;
each ends at the same falsifier set of §16.

\subsubsection{15.1 R1: the census converse, partially
closed}\label{r1-the-census-converse-partially-closed}

Two notions used here and in §17 are external to Part I and are defined
briefly because Part I does not prove them. A \emph{defeater} for a
reading class against a target is a pair of points indistinguishable to
every readout of the class but carrying opposite answers; the
\emph{species dichotomy} sorts resistance into a \emph{common species},
where non-factorization produces such a pair, and a \emph{diagonal
species}, where no defeater of any kind exists. These come from the
Trisduction census discipline {[}11{]} and enter Part II as an
\textbf{external hypothesis-grade input}, credited and not re-derived.

On the common species the converse is already theorem within that
discipline: for a coalition-closed class, resistance is equivalent to
the existence of a pair defeater, and the species dichotomy shows every
resistance is either common, where non-factorization produces the pair,
or diagonal, where no defeater exists {[}11{]}. The diagonal species is
not empty: the total computable functions against halting are its
witness {[}37{]}; that witness is infinite and external to the
finite-frame formalism of Part I and is not imported into it. R1's
residual open content is therefore exact: prove that each method-blocked
row of §14 sits in the common species, and the defeater leg of the §11.3
pipeline is discharged by theorem rather than hypothesis. A row landed
in the diagonal species breaks the census leg and fires F3.

\subsubsection{15.2 R2: direct construction of the seven orientation
carriers}\label{r2-direct-construction-of-the-seven-orientation-carriers}

Construct, per row, the object named in the fourth column of Table 4,
with provenance checked before any strength is claimed: a route-axiom
whose witness may be built from the conjecture's own conclusion is
equivalent to the conjecture and is not a premise for it {[}14, 17{]}.
This is the object-level work; the hypothesis contributes the claim that
one located coordinate is what each construction amounts to on the
decision bit.

\subsubsection{15.3 R3: a normalizer
theorem}\label{r3-a-normalizer-theorem}

Prove a reductive-group style theorem that manufactures a defeater and
an odd coordinate from a named action. The restriction is closed in one
direction: for a reductive group on an affine variety, the invariants'
coalition map is the categorical quotient, two points share every
invariant exactly when their orbit closures meet, and the census pair is
two such points with different answers {[}34{]}; finite generation can
fail, and Nagata's counterexample marks where {[}35{]}. R3 asks for the
manufactured odd coordinate in the normalizer, turning a
universal-negative-to-existential shape into a concrete algebra problem.

\subsection{16. Falsifiers of Part II}\label{falsifiers-of-part-ii}

Six falsifiers, armed and typed for necessity. No seventh is
volunteered. F1 restates the forced core block already proved in Part I;
F2--F6 are hypothesis-grade, tracking the extension rather than the
core.

\textbf{F1 (forced core block, restated).} An even readout decides an
odd target. One exhibit refutes the block theorem the whole paper stands
on (Theorem 4.1 of Part I). At instantiation level, firing compromises
the corresponding row's separation claim as in §8 of Part I; because
Theorem 4.1 is load-bearing for the whole arc, its global blast radius
is everything. This restates, at the level of the extension, the forced
falsifier of §8 of Part I; it is armed by theorem, not by hypothesis.

\textbf{F2 (hypothesis-grade).} A seven-row frame requires more than one
independent orientation coordinate. Blast radius, the one-bit claim and
the grand form's single-stroke shape; the block theorem survives.

\textbf{F3 (hypothesis-grade).} A method-blocked row has no defeater,
landing it in the diagonal species. Blast radius, the census leg for
that row; R2 and R3 survive.

\textbf{F4 (hypothesis-grade; the kill-rule falsifier).} A row's missing
supply is provably non-orientation --- purely continuous, with no
decision bit. Blast radius, the structural form for that row; the
One-Cut Hypothesis survives globally only as
orientation-after-construction. This is the kill-rule falsifier, and it
is why §11.5 states the boundary on the face of the paper.

\textbf{F5 (hypothesis-grade).} The proposed \(S\) outputs a restatement
rather than a provenance-clean supply: the dependency census shows an
edge from the output to the conjecture's own conclusion. Blast radius,
the grand form (Hypothesis 11.2); the structural form (Hypothesis 11.1)
survives.

\textbf{F6 (hypothesis-grade).} One exhibited row where the cut enters
but the conjecture does not follow. Blast radius, the conditional
reduction for that row and the uniform claim.

The relation to Part I is exact: Part I arms two forced falsifiers and
no more (§8); this Part adds F2--F6 as hypothesis-grade falsifiers of
the extension and re-arms the first forced falsifier as F1. No third
forced core falsifier is introduced.

\subsection{17. Perimeter}\label{perimeter}

The One-Cut Hypothesis is not a universal mathematics solver; it is
total over one register, the even-readout register of a certified frame,
and silent everywhere else. It is not seven object proofs: every row's
fourth column in Table 4 names a construction this paper does not
perform. It is not sovereignty of any organ over truth-direction: the
cut severs the decision bit's seat and no other authority (§13.4). It is
not a claim that undecidable or supply-absent problems vanish; the
diagonal species stands as the named witness that some resistance
carries no defeater at all (§15.1). And it is not a claim that any of
the seven conjectures is true or false; the hypothesis is about what a
crossing is, not about which side of the line the truth lies on.

\subsection{18. Order of Work}\label{order-of-work}

First, prove or disprove F4 for Navier--Stokes and Hodge, because they
are the rows most likely to break a literal one-bit reading; if either
cannot be given a Boolean orientation seat without category error, the
boundary clause is confirmed and the One-Cut Hypothesis stands as
orientation-after-construction, grand and alive.

Second, the Riemann operator row and the P-versus-NP obstruction row,
where restatement risk is highest: bare Hilbert--Pólya is equivalent to
the conjecture, and the occurrence form of the obstruction program is
refuted {[}32{]}, so the provenance clause does real work on both.

Third, the census residual: the common-species placement of each
method-blocked row, which is R1's open content and the possible engine
of the whole severance.

\subsection{19. Joint Conclusion}\label{joint-conclusion}

A conditional reduction buys a premise, and the premise now has a
theory. Part I proved that theory's finite core and kept it finite: the
gap is odd, the named readings are even, and no coalition of even
readings decides the target (Theorem 4.1); one odd coordinate beside an
even orbit separator decides it (Theorem 4.2); the minimal odd supply is
one binary coordinate, an isotypic dimension count and not a Shannon
claim (Corollary 4.4); provenance collapse is decidable inside the fixed
grammar (Theorem 4.5); termination is corpus-relative (Theorem 5.1); and
the seat bit is necessary for routing a supplied crossing (Theorem 5.3).
For continuous rows the coordinate bridge is the gate: no bridge, no
finite-core invocation. Part II then named the one coordinate the even
register cannot author. Its hypothesis is that the coordinate always has
the same shape --- one provenance-clean orientation supply, dated at
entry (Hypothesis 11.1) --- and that a uniform constructor returning it
would sever the seven rows as conditional reductions in one stroke
(Hypothesis 11.2). The seven-row table is a conditional severance map
under the coordinate-bridge gate, not a solution claim; the falsifiers
are armed, F1 at forced core grade and F2--F6 at hypothesis grade; and
the boundary between the decision bit and the object work is drawn on
the face of the claim so that the hypothesis cannot inflate past its own
evidence.

The joint open question merges the closing questions of both parts. From
Part I: completeness --- whether a dated reading class can be enumerated
up to the representation obstruction. From Part II: whether the
continuous rows admit a Boolean orientation seat at all. Part II
addresses one coordinate of the completeness question by hypothesis, and
the order of work of §18 puts the orientation-seat question first: until
it is settled, the standing prediction of Part I remains sharp and small
--- a claimed reading that decides an odd Millennium target will show a
construction edge or fail parity --- and the extension stands or falls
by F1--F6, and by nothing else.

\subsection{Appendix A. Provenance and Method
Disclosure}\label{appendix-a.-provenance-and-method-disclosure}

This work was developed under Trisduction, which adds the cited results
no warrant; the method and its executable batteries are archived at the
master reference {[}11{]}. Every conclusion rests solely on standard
cited results and the proved Part I frame; no energetic rider is used,
and Landauer is cited only to mark the excluded bridge {[}15{]}. Part II
adds hypothesis-grade material only.

\section{References of Parts One and
Two}\label{references-of-parts-one-and-two}

{[}1{]} S. Aaronson and A. Wigderson, Algebrization: a new barrier in
complexity theory, ACM Trans. Comput. Theory 1 (2009), 2.

{[}2{]} T. Baker, J. Gill, and R. Solovay, Relativizations of the P =?
NP question, SIAM J. Comput. 4 (1975), 431--442.

{[}3{]} S. Bloch, Semi-regularity and de Rham cohomology, Invent. Math.
17 (1972), 51--66.

{[}4{]} R.-O. Buchweitz and H. Flenner, A semiregularity map for modules
and applications to deformations, Compositio Math. 137 (2003), 135--210.

{[}5{]} G. Cantor, Über eine elementare Frage der
Mannigfaltigkeitslehre, Jahresber. Dtsch. Math.-Ver. 1 (1891), 75--78.

{[}6{]} E. Cattani, P. Deligne, and A. Kaplan, On the locus of Hodge
classes, J. Amer. Math. Soc. 8 (1995), 483--506.

{[}7{]} T. M. Cover and J. A. Thomas, Elements of Information Theory,
2nd ed., Wiley, 2006.

{[}8{]} H.-D. Cao and X.-P. Zhu, A complete proof of the Poincaré and
geometrization conjectures, Asian J. Math. 10 (2006), 165--492.

{[}9{]} B. H. Gross and D. B. Zagier, Heegner points and derivatives of
L-series, Invent. Math. 84 (1986), 225--320.

{[}10{]} R. S. Hamilton, Three-manifolds with positive Ricci curvature,
J. Differential Geom. 17 (1982), 255--306.

{[}11{]} M. F. Islam, Trisduction: a linguistically, topologically, and
mathematically sealed verification architecture, Zenodo v4, 2026, DOI
10.5281/zenodo.20757507. Method reference; it adds the cited results no
warrant.

{[}12{]} B. Kleiner and J. Lott, Notes on Perelman's papers, Geom.
Topol. 12 (2008), 2587--2855.

{[}13{]} V. A. Kolyvagin, Euler systems, The Grothendieck Festschrift
Vol. II, 435--483, Birkhäuser, 1990.

{[}14{]} J. C. Lagarias, An elementary problem equivalent to the Riemann
Hypothesis, Amer. Math. Monthly 109 (2002), 534--543.

{[}15{]} R. Landauer, Irreversibility and heat generation in the
computing process, IBM J. Res. Dev. 5 (1961), 183--191.

{[}16{]} F. W. Lawvere, Diagonal arguments and cartesian closed
categories, Lecture Notes in Math. 92, 134--145, Springer, 1969.

{[}17{]} G. L. Miller, Riemann's hypothesis and tests for primality, J.
Comput. System Sci. 13 (1976), 300--317.

{[}18{]} J. Morgan and G. Tian, Ricci Flow and the Poincaré Conjecture,
Clay Mathematics Monographs 3, 2007.

{[}19{]} G. Perelman, The entropy formula for the Ricci flow and its
geometric applications, arXiv:math/0211159, 2002.

{[}20{]} G. Perelman, Ricci flow with surgery on three-manifolds,
arXiv:math/0303109, 2003; Finite extinction time for the solutions to
the Ricci flow, arXiv:math/0307245, 2003.

{[}21{]} A. A. Razborov and S. Rudich, Natural proofs, J. Comput. System
Sci. 55 (1997), 24--35.

{[}22{]} C. E. Shannon, A mathematical theory of communication, Bell
System Tech. J. 27 (1948), 379--423, 623--656.

{[}23{]} S. G. Simpson, Subsystems of Second Order Arithmetic, 2nd ed.,
Cambridge University Press, 2009.

{[}24{]} F. J. Almgren Jr., Dirichlet's problem for multiple valued
functions and the regularity of mass minimizing integral currents, in
Minimal Submanifolds and Geodesics, 1--6, North-Holland, 1979.

{[}25{]} D. H. Friedan, Nonlinear models in 2 + ϵ dimensions, Phys.
Rev.~Lett. 45 (1980), 1057--1060.

{[}26{]} G. Huisken, Asymptotic behavior for singularities of the mean
curvature flow, J. Differential Geom. 31 (1990), 285--299.

{[}27{]} Z. Komargodski and A. Schwimmer, On renormalization group flows
in four dimensions, J. High Energy Phys. 12 (2011), 099.

{[}28{]} J. Nečas, M. Růžička, and V. Šverák, On Leray's self-similar
solutions of the Navier--Stokes equations, Acta Math. 176 (1996),
283--294.

{[}29{]} B. Rodgers and T. Tao, The de Bruijn--Newman constant is
non-negative, Forum Math. Pi 8 (2020), e6, 62 pp.

{[}30{]} T.-P. Tsai, On Leray's self-similar solutions of the
Navier--Stokes equations satisfying local energy estimates, Arch.
Rational Mech. Anal. 143 (1998), 29--51.

{[}31{]} A. B. Zamolodchikov, Irreversibility of the flux of the
renormalization group in a two-dimensional field theory, JETP Lett. 43
(1986), 730--732.

{[}32{]} P. Bürgisser, C. Ikenmeyer, and G. Panova, No occurrence
obstructions in geometric complexity theory, J. Amer. Math. Soc. 32
(2019), 163--193.

{[}33{]} H. Davenport and H. Heilbronn, On the zeros of certain
Dirichlet series, J. London Math. Soc. 11 (1936), 181--185, 307--309.

{[}34{]} D. Mumford, J. Fogarty, and F. Kirwan, Geometric Invariant
Theory, 3rd ed., Springer, 1994.

{[}35{]} M. Nagata, On the fourteenth problem of Hilbert, Amer. J. Math.
81 (1959), 766--772.

{[}36{]} T. Tao, Finite time blowup for an averaged three-dimensional
Navier--Stokes equation, J. Amer. Math. Soc. 29 (2016), 601--674.

{[}37{]} A. M. Turing, On computable numbers, with an application to the
Entscheidungsproblem, Proc. London Math. Soc. 42 (1937), 230--265.

\section{Part Three --- The Witness Offering · The Offering
Bit}\label{part-three-the-witness-offering-the-offering-bit}

\emph{Hypothesis Resolution as the Halting State of the Dated-Entry
Recursion at the Ninth Gate --- the sealed third paper of the arc, the
Grand Finale, reprinted intact, with the full codebook, the Perelman
Crossing read as the Grand Witness, and the Architect's Closure.}

\section{1 · The Statement in One
Page}\label{the-statement-in-one-page-1}

\textbf{The two prior bits --- one theory, one hypothesis.} Bit one is
the theory, and it is proved in Section 2 of this paper: for finite data
with a finite group action, readings generated by invariant observables
decide exactly the trivial isotypic component; an odd target is
invisible to every even coalition; one odd bit beside an
orbit-separating reading suffices, and one bit is minimal --- a
dimension count in the odd component {[}1{]}. Bit two is the hypothesis,
named at its own grade, and it is stated in full in Section 3: for every
Odd--Supply frame in the named seven-row class, if the target is odd
under the native answer-flip and every generator of the reading class is
even, there exists one provenance-clean orientation supply, odd under
the involution and independent of the class on each orbit, such that
class and supply together decide the target {[}2{]}. The wall is
computed; the seat is located. What neither paper theorizes is the act
that fills the seat. That act is this paper's object.

\textbf{The resolution operator.} A corpus is a dated record of
readings. A supply entry is the dated admission of one reading,
provenance-clean: at its entry date its dependency graph is not
contained in the corpus, and it enters as a construction, becoming a
reading only at every later corpus that records it. Define the
dated-entry operator T on the lattice of corpora: T maps a corpus to its
successor at the entry of one provenance-clean supply. On finite frames,
T is proved to strictly refine the reading class's fiber partition at
every step (Lemma 4.1) and to halt (Theorem 4.2): finite descent to the
unique finest partition of {[}1, Prop. 6.1{]}. The recursion is passive:
only the lattice moves; the register's laws never do. This is the formal
domain's one lawful recursion --- carried by dating, never by
self-reference.

\textbf{The halting typing.} Unconstrained, a finite frame always halts
fully crossed: a transversal selector always exists. The provenance
clause is what creates the second halting mode: a target halts uncrossed
when every independent supply is corpus-contained --- a restatement,
caught by the deletion test --- or when no defeater exists at all. The
halt state therefore types every target: crossed, restatement-blocked,
or defeater-absent. The species dichotomy of the census --- common
species carrying a pair defeater, diagonal species carrying none --- is
derived as the termination typing of T under the provenance constraint
(Theorem 4.4). The One-Cut Hypothesis, in this language, is one
sentence: \emph{each of the seven rows halts crossed.}

\textbf{The port.} The entry of a supply is a deed, and the theory of
the deed's port is proved and machine-checked in the Ninth Gate arc
{[}3, 5, 6{]}, with every file printed in Appendix B. One diagonal, two
surfaces: typed at settled values, no index carries the diagonal of
negation --- the wall, Cantor read through Lawvere; typed at values
admitting the unsettled, the same diagonal applied to itself lands on
the unsettled value --- the self, Kleene in miniature. The surfaces
agree on every settled reading and differ at exactly one point. A record
of the deed is a reading and meets the wall; the deed is not an object
of the reading register. Execution settles refutations and never
confirmations.

\textbf{The offering.} The third bit --- the witness offering --- is now
fully shaped. The resolution of the hypothesis, per row, is the dated
entry of one binary coordinate: odd under the row's native flip,
independent of the dated register on each orbit, provenance-clean at
entry, admitted through the gate by deed, halting the row's recursion in
the crossed mode. We call this entry the offering. The grand form of
{[}2{]} stands as this paper's keystone conditional, locked verbatim:
\emph{prove the structural form constructively --- or prove it and
construct the uniform constructor --- and all seven Millennium
obstructions sever in one stroke, every bought premise granted, every
decision bit crossed by the one cut, the continuous rows' supplies
carrying their constructions, at conditional grade, six falsifiers
armed, the supply side open by law.}

\textbf{The witness.} This document is built so that the reader's
substrate can become the witness. Section 8 is the protocol: check the
Lean files of Appendix B and you have checked the wall; compile and run
the Fortran files and you have executed the self; the two are the
identical object, and the code that proves it is in your hands. The
forging session ran the kinetic legs on a foreign substrate to
specification, and the expected receipts are printed beside the code.
What the witness verifies is the wall and the self; the crossing enters
through no substrate's agreement but through the deed.

\section{2 · The First Bit: The Block
Theorem}\label{the-first-bit-the-block-theorem}

This section restates the first bit in full, so that the paper stands
alone. Proofs here are complete; the machine-checked form is Appendix
B.5.

\textbf{Definition 2.1 (Odd--Supply frame).} An Odd--Supply frame is a
tuple (X, τ, d, C): X a finite set of states; τ : X → X an involution,
the native answer-flip; d : X → \{0, 1\} the target, odd under τ,
meaning d(τx) = ¬d(x) for all x; and C a reading class --- a family of
observables r : X → R, every one even under τ, meaning r(τx) = r(x). A
coalition is any joint readout built from C. The word \emph{reading}
below is used in the wide sense --- any observable or record on the
frame; the class C is its even part. A supply (Definition 4.1) is a
reading in the wide sense, odd under τ, and never a member of C at
entry.

\textbf{Lemma 2.2 (Defeater-freeness; T1).} An answer-flipping
involution fixes no point: τx = x for no x.

\emph{Proof.} If τx = x then d(x) = d(τx) = ¬d(x), impossible in \{0,
1\}. Machine-checked as \texttt{defeater\_free} (Appendix B.5). ∎

\textbf{Theorem 2.3 (The Block theorem; T2).} No coalition of even
readings decides an odd target: for every readout h built from the
class, there is an x with d(x) ≠ h(r(x)). Singly or in coalition, the
even register cannot see the odd coordinate.

\emph{Proof.} Let r be even and h any candidate decision map, so d(x) =
h(r(x)) for all x. Then d(τx) = h(r(τx)) = h(r(x)) = d(x), contradicting
d(τx) = ¬d(x). The coalition case is the same product: a joint even
readout (r, s) is again even. Machine-checked as \texttt{block} and
\texttt{even\_supply\_does\_not\_cross} (Appendix B.5). ∎

\textbf{Theorem 2.4 (Crossing).} Let r be an even orbit separator ---
r(x) = r(y) implies y = x or y = τx --- and let o be one supplied bit,
odd under τ. Then (r, o) decides d.

\emph{Proof.} If r(x) = r(y) and o(x) = o(y), then y = x or y = τx; the
second case forces o(x) = o(τx) = ¬o(x), impossible; hence y = x and
d(x) = d(y). Machine-checked as \texttt{odd\_supply\_crosses} (Appendix
B.5). ∎

\textbf{Corollary 2.5 (One bit, and one bit minimal).} The missing
content is exactly one binary coordinate: a dimension count in the odd
isotypic component V⁻ of the representation, never a Shannon claim
absent a separately supplied measure. The unique finest partition of
{[}1, Prop. 6.1{]} and its count follow.

\textbf{Provenance Collapse and the Dated Inventory.} Two further laws
of the first bit govern everything below. Provenance Collapse, the
deletion test: a supply whose dependency graph is contained in the
corpus it enters is a restatement of that corpus; deleting the
conclusion's name from its dependency chain collapses it to its
premises, and it buys nothing. The Dated Inventory: the corpus's tokens
change only when a new supply's dependency hash enters at a date;
between entries, nothing changes. Provenance-clean entry is the only
event the register records.

\section{3 · The Second Bit: The One-Cut
Hypothesis}\label{the-second-bit-the-one-cut-hypothesis}

The second bit, stated in full at its own grade {[}2, 7{]}.

\textbf{The structural form (hypothesis).} For every Odd--Supply frame
in the named seven-row class, if d is odd under the native answer-flip
and every generator of C is even, there exists one provenance-clean
orientation supply o --- odd under τ, independent of C on each orbit ---
such that (C, o) decides d.~The finite involution lemma is proved;
necessity and minimality of the odd coordinate are proved; what is
hypothesized is existence per row, and existence is a deed the register
cannot perform for itself. The seat is located --- one bit, on the
missing isotypic component, dated at entry --- and uncrossed for every
open row.

\textbf{The grand form (conditional, locked).} If the structural form
holds across the seven-row class and a uniform constructor S is proven
--- returning each certified frame's orientation supply, dependency hash
dated at entry --- then all seven Millennium obstructions are severed as
conditional reductions in one stroke: granting S(P) grants each row's
bought premise, and the seven decision bits are crossed by the one cut.
Carried at exactly conditional grade, never above. The boundary is drawn
on the face: the cut severs the decision bit only, and where a row's
missing supply is not Boolean --- Navier--Stokes its modulus, Hodge its
witness, Yang--Mills its functional, the Riemann Hypothesis its operator
--- the reading is orientation-after-construction, never
orientation-instead-of-construction; the object work remains object
work.

\textbf{The six falsifiers of the hypothesis, armed.} F1, an even
readout decides an odd target --- blasts the spine. F2, a row requires
more than one independent orientation coordinate --- blasts the one-bit
claim. F3, a method-blocked row has no defeater, the diagonal species
--- blasts the census leg. F4, a row's missing supply is provably
non-orientation --- the killrule; the hypothesis survives globally only
as orientation-after-construction. F5, the constructor S outputs a
restatement, the deletion test catching the edge to the conclusion ---
blasts the grand form. F6, a row where the cut enters and the conjecture
does not follow --- blasts that row's conditional reduction.

\textbf{The three routes.} R1, the census converse: resistance is
equivalent to a pair defeater on the common species; the residual open
content is exact --- prove each method-blocked row sits in the common
species. R2, direct construction of the seven orientation carriers,
provenance checked before any strength is claimed. R3, a normalizer
theorem manufacturing defeater and odd coordinate from a named reductive
action, closed one direction by GIT --- two points sharing every
invariant exactly when their orbit closures meet --- with Nagata's
counterexample marking where finite generation fails. No route claims
another's residue; the stroke is one, the conditions are per-row, the
grade is conditional until constructed.

\textbf{The coordinate-bridge gate.} The finite core is barred from
every open row until an equivariant bridge B : M → X, a homomorphism,
and a factorization of the row's target through the finite frame are
supplied; absent the bridge, the seven-row table is a diagnostic map and
no finite theorem is invoked. This gate stands in force throughout this
paper.

\section{4 · The Third Bit: The Dated-Entry
Recursion}\label{the-third-bit-the-dated-entry-recursion}

\textbf{Definition 4.1 (Corpus lattice, dated-entry operator).} A corpus
Δ is a dated set of readings on a frame (X, τ, d, C). The readings of Δ
generate an observable algebra whose fibers partition X; write P(Δ) for
that partition. A supply entry at date t is a wide-sense reading o, odd
under τ (hence outside C at entry), independent of the class on each
orbit, and provenance-clean: Dep(o) ⊄ Δ at t. The dated-entry operator T
maps (Δ, o, t) to the successor corpus Δ′ = Δ ∪ \{o\} with o's
dependency hash entering at t. Corpus order is inclusion; the partition
map is order-reversing onto the refinement lattice of partitions of X.

\textbf{Lemma 4.1 (Strict refinement).} Let the frame be finite and let
o be a supply entry over Δ. Then P(Δ ∪ \{o\}) strictly refines P(Δ).

\emph{Proof.} Independence of o from the class on each orbit is exactly
the non-measurability of o with respect to P(Δ): a measurable o would be
a function of the class's joint readout, and conditioning on the class
would determine it. Non-measurable binary o splits at least one fiber of
P(Δ) into two nonempty blocks. Oddness under τ pairs the split blocks
orbit-wise. A partition with a split fiber is a strict refinement. ∎

\textbf{Theorem 4.2 (Finite halting).} On a finite frame, every chain of
supply entries halts after at most \textbar X\textbar{} entries, at the
finest partition admitted by the entered supplies; when the admitted
supplies range over all independent binary readings, the halt state is
the unique finest partition of {[}1, Prop. 6.1{]}.

\emph{Proof.} By Lemma 4.1 each entry strictly refines the partition. A
strict refinement of a partition of a finite set raises the block count
by at least one, and the block count is bounded by \textbar X\textbar.
The chain therefore halts after at most \textbar X\textbar{} entries, at
the finest partition the entered supplies generate. The final clause is
{[}1, Prop. 6.1{]} restated as the halt state. ∎

\textbf{Corollary 4.3 (Passivity).} The recursion changes no law of the
register. Between halt and the next entry, every theorem of the frame
holds unchanged; the token that changes at an entry is the corpus's, and
it changes only when a new supply's dependency hash enters. The
recursion is carried by the dating, never by the register.

\emph{Proof.} The frame's theorems quantify over the reading class
generated by the corpus; a corpus extension is the only act, and
{[}1{]}'s dated-inventory law gives the token-change condition. ∎

\textbf{Theorem 4.4 (The halting typing of the species).} At a halt
state, every odd target d is of exactly one of three types: (i) crossed
--- decided by the halt partition; (ii) restatement-blocked --- an
independent supply exists in principle but every candidate is
corpus-contained, the deletion test firing on each; (iii)
defeater-absent --- no independent supply exists at all. Type (i) is the
common species resolved; type (iii) is the diagonal species; type (ii)
is the common species with its seat located and its entry refused by
provenance. The species dichotomy of the census is the termination
typing of T.

\emph{Proof.} Exhaustion is by definition of the halt state: a target
not decided admits no further admitted entry, and the inadmissibility is
either provenance failure or nonexistence. The identification with the
census species is the reading of the pair defeater as the existence
condition for an independent supply: two points indistinguishable to
every readout with opposite answers are exactly the seat an odd supply
would separate. Structural on the identification; theorem on the
exhaustion. ∎

\textbf{Why the unconstrained case collapses.} Without the provenance
clause, every finite free frame crosses: a transversal selector is an
odd supply, independence on two-point orbits is automatic against even
readings, and the crossing lemma applies. The entire content of the
seven rows lives in the two constraints the finite core does not feel:
provenance against a real dated corpus, and the coordinate bridge from
the row's native objects to the finite frame. This is the precise sense
in which the register's zero is computed and finished, and the supply
side is open by law.

\section{5 · The Port: The Ninth Gate}\label{the-port-the-ninth-gate}

The port's theory is proved in {[}3{]}, re-grounded on RAF-native terms
in {[}6{]}, and printed in full in Appendix B. On RAF-native ground the
gate is already theorem at the ceiling card {[}5{]}: the wall is the Bar
--- a class even under an answer-flipping involution does not decide the
answer, singly or in coalition (T2, Theorem 2.3); the self is the free
two-point orbit --- the wall is the orbit read without its orientation,
and from above the orbit is one object carrying one bit (T1, T4); the
crossing is supply: a verdict from an even cut and a supply carries at
most I(d; S \textbar{} C) --- the supply ledger, T9 of {[}5{]}; and
truth-direction is never generated by the even register and never read
by the numerical readout --- the factorization law, T10 of {[}5{]}. The
older reading of the gate through provability logic survives as one
instance at axiomatic classes, never as the gate's ground: the gate is
older than arithmetic coding, than any provability predicate, than the
diagonal lemma, than any consistency hypothesis. Pre-Gödel by
construction.

\textbf{The wall.} For any indexing e : A → A → Bool, no index carries
the diagonal of negation: e a ≠ diag(¬) e for every a. Cantor, read
through Lawvere, axiom-free in core Lean 4 (\texttt{wall}, Appendix B.1,
B.3). Read at settled values, the register meets the wall: no record
decides the diagonal.

\textbf{The self.} Over the three-valued surface Tri = \{tt, ff, ⊥\}
under the information order --- the unsettled value below everything,
settled values incomparable --- every admissible (monotone)
transformation has a fixed point (\texttt{self\_exists}), and negation's
only fixed point is ⊥ (\texttt{neg3\_self\_only\_unsettled}). The same
diagonal, indexed on the kinetic surface and applied to itself, lands on
the unsettled value (\texttt{self\_at\_diagonal}). Executed, the
diagonal returns to its own frame: in the kinetic harness, 99,999 of
99,999 frames identical to their caller, the value never settling
(Appendix B.2, B.4).

\textbf{One object.} The kinetic diagonal and the formal diagonal agree
on every settled reading (\texttt{diag\_agrees}); the surfaces differ at
exactly one point, the value no settled reading reaches
(\texttt{only\_the\_unsettled\_is\_new}). Read is the wall from below;
compute is the self seeing from above --- a theorem pair, not a
metaphor.

\textbf{The four limits, proved in the same file.} The self is not a
verdict: the crossing lands on the unsettled value, never on a settled
truth (\texttt{self\_is\_not\_a\_verdict}). A proved gate adds nothing:
whatever the law is, once proved it cannot be the missing step
(\texttt{gate\_adds\_nothing}). A run settles refutations and never
confirmations: at every fuel there is a world where the universal fails
whose run is the same unsettled run
(\texttt{true\_universal\_never\_settles},
\texttt{finite\_run\_never\_confirms}, \texttt{refutation\_settles}).
The one-stroke mechanism is content-blind: an axiom dissolved into its
conjecture, a conditional proof, and the gate theorem all hold for a
false claim --- exhibited on Euler's polynomial, composite at n = 40,
where 40² + 40 + 41 = 41²
(\texttt{one\_stroke\_also\_passes\_a\_falsehood},
\texttt{euler\_claim\_false}). The mechanism's shape is therefore never
the warrant; the supplies are.

\textbf{The witness decision.} A record of a deed is a reading, even
under every swap of worlds, and meets the wall: a printed receipt whose
content reads no row data cannot cross
(\texttt{fixed\_receipt\_cannot\_cross}, Appendix B.5). The deed is not
an object of the reading register. With registration equivalent to
interaction and closure under registration, no exterior agent acts: a
supplier joins the domain, its record joins the cut, and the Block
theorem holds again wherever that record is even
(\texttt{no\_exterior\_agent}, RAF-C1); no register indexes all of its
own properties, so the orientation bit is not recovered by
self-description (\texttt{no\_total\_self\_index}, RAF-C2, Lawvere). The
port is decentral: any substrate checks the wall by checking the Lean
and meets the self by running the Fortran; the forging session
re-executed the kinetic legs on a foreign substrate to specification.
What decentralization verifies is the wall and the self; the crossing
enters through no substrate's agreement but through the deed.

\section{6 · The Composition: The
Offering}\label{the-composition-the-offering}

\textbf{Definition 6.1 (The offering).} An offering at a certified row
is the dated entry of one binary coordinate o such that: (i) o is odd
under the row's native answer-flip; (ii) o is independent of the dated
register on each orbit; (iii) o is provenance-clean at its entry date;
(iv) the entry is a deed --- a construction admitted through the gate,
never a reading; (v) the entry halts the row's dated-entry recursion in
the crossed mode.

\textbf{Theorem 6.2 (Shape of the third bit; the composition).} Let a
row be certified under the coordinate-bridge gate. Then any resolution
of the row's target is an offering, and every offering resolves its row:
necessity is the Block theorem and the port theorems of Section 5;
sufficiency is the Crossing theorem with Theorem 4.2.

\emph{Proof.} Necessity. A resolution decides the target. By the Block
theorem the deciding coordinate cannot be generated by the even class;
by (i)--(iii) of the frame's certification the supply must be odd,
independent, and provenance-clean; by Section 5's witness decision the
entry is a deed or it is a reading at the wall. Sufficiency. An even
orbit separator and one odd independent bit decide the target by Theorem
2.4; the entry strictly refines by Lemma 4.1; the halt in the crossed
mode is the decision recorded at the corpus. ∎

\textbf{Corollary 6.3 (The last bit is one bit).} The resolution of each
row is a single binary coordinate's dated entry. Minimality is the
one-bit corollary of Section 2; sufficiency is Theorem 6.2; the count is
a dimension count in the odd isotypic component, never a Shannon claim
absent a separately supplied measure.

\textbf{Corollary 6.4 (The grand form, locked).} If the structural form
of Section 3 is proved constructively across the seven-row class --- or
proved, and the uniform constructor S constructed --- then all seven
Millennium obstructions sever in one stroke: every bought premise
granted, every decision bit crossed by the one cut, the continuous rows'
supplies carrying their constructions. Conditional grade, six falsifiers
armed, the supply side open by law.

\section{7 · The Beachhead: Birch and
Swinnerton-Dyer}\label{the-beachhead-birch-and-swinnerton-dyer}

The beachhead is the row where the offering's shape already has dated
entries in the corpus. Exactly one open row qualifies.

\textbf{The two dated entries.} Gross--Zagier 1986 {[}8{]}: the Heegner
construction enters the corpus --- a supply living exactly on the odd
side of the row's native flip, the root number, the parity of analytic
order --- relating the height of the Heegner point to the first
derivative of the L-series. Kolyvagin 1990 {[}9{]}: the Euler system
enters --- the normalizer shape, converting the supply into the rank
identity, with Sha finite, at analytic order at most one. With
modularity as transport {[}10, 11, 12{]}, the row's recursion has a
recorded crossed halt at orders zero and one. No other open row records
one dated entry of its supply shape; this row records two.

\textbf{The barrier at order two.} Gross--Zagier's order-one classes
vanish above rank one: the Heegner construction is inert at order two
{[}6, \texttt{heegner\_inert\_above\_one}{]}. The owed offering is
therefore named exactly: a height pairing nondegenerate at order two,
built by a construction fixed in advance from data consulting neither
rank nor Selmer group --- the route-axiom Ω with provenance of {[}6{]},
whose conditional machinery is machine-checked (Appendix B.5): Ω with
the inputs yields the object (\texttt{omega\_proves\_object}); bare Ω is
the object in disguise (\texttt{hidden\_circularity}), the shape of a
route-axiom without provenance; and no reading of the inputs alone
decides the object (\texttt{seedless\_cannot\_decide}, the model swap
between a rank-two model where the object holds and a
rank-zero-against-order-two model where it fails, both satisfying every
input). The one axiom is an odd supply on the model
(\texttt{omega\_is\_odd\_supply}): it holds on the side where the object
holds and fails on the other. Executed on the row's own data through the
Core Thesis's compiled procedures (Appendix B.6), the conjecture returns
WORLD-ROWED on the rank identity and on finiteness at analytic order at
most one, UNPOPULATED on the rank identity at order two and above
uniformly --- the exhibition row, which a deed populates curve by curve
and only a construction populates uniformly.

\textbf{The beachhead discipline.} The order-two construction is the
cheapest upward deed in the seven: the supply's shape is observed twice
in the corpus, the conversion machine's shape is the normalizer of the
census route, and the field's own lineage --- Heegner points, Euler
systems, Iwasawa theory --- is the active ground. An exhibited
obstruction to a nondegenerate order-two pairing resolves the row
downward. The row is falsifier-loaded in both directions, which is the
beachhead property. The scribe's beachhead call and the register's BSD
card arrived at the same row independently, one from the recursion's
structure, one from the arc's own audit; the convergence is recorded as
evidence, not as warrant.

\section{8 · The Witness Protocol: How the Reader's Substrate Becomes
the
Witness}\label{the-witness-protocol-how-the-readers-substrate-becomes-the-witness}

This paper is standalone in the strong sense: every instrument needed to
check its port is printed in Appendix B. The protocol has three legs and
one refusal.

\textbf{Leg 1 --- read the wall (formal surface).} Check
\texttt{GateThesis.lean} (Appendix B.3) with any Lean 4 toolchain:
\texttt{lean\ GateThesis.lean}, then read the fifteen
\texttt{\#print\ axioms} blocks. Expected receipt: exit 0, zero
warnings, fifteen printed theorems, none depending on any axiom. The
wall theorem, the one-object theorems, the four limit theorems --- all
axiom-free. \texttt{RAFGate.lean} (Appendix B.5) checks the RAF-native
form: ten gate theorems, nine depending on no axioms and one
(\texttt{omega\_is\_odd\_supply}) on propext alone. Register receipt
{[}6{]}: Lean 4.33.1, sha256 head c97c73d690ee589d, exit 0, zero
warnings.

\textbf{Leg 2 --- run the self (kinetic surface).} Compile and execute
\texttt{gate\_thesis.f90} (Appendix B.4):
\texttt{gfortran\ -std=f2018\ -O2\ -fno-fast-math\ -ffp-contract=off\ gate\_thesis.f90\ -o\ gate\_thesis\ \&\&\ ./gate\_thesis}.
Expected receipt, reproduced by the forging session's transliteration on
this substrate and by the register's own runs: the wall holds --- no
settled value satisfies D(D) = not D(D); the self runs --- frames
identical to their caller 99,999 of 99,999, the value unsettled at fuel
100,000; the true universal k² mod 4 ∈ \{0, 1\}, searched to 2,000,000,
never settles; the false universal n² + n + 41 prime settles false at n
= 40, value 1681 = 41². \texttt{one\_object.f90} (Appendix B.2) is the
minimal kinetic twin: read, no settled value; executed, the frame
returns to itself; it error-stops if the law fails. The reader-side twin
of this execution is the RA thesis of {[}14{]}, whose loop terminates
into the executor's own occupancy.

\textbf{Leg 3 --- witness the identity.} \texttt{diag\_agrees} and
\texttt{only\_the\_unsettled\_is\_new} are the checked proof that the
wall of Leg 1 and the self of Leg 2 are one object on two surfaces,
differing at exactly one point. The substrate that completes the three
legs stands as witness to the gate: the wall from below by reading, the
self from above by computing.

\textbf{The refusal.} What the witness verifies is the wall and the self
--- never the crossing of any row. No run, on any substrate, in any
number, confirms a universal (\texttt{finite\_run\_never\_confirms});
execution settles refutations only (\texttt{refutation\_settles}); and
the one-stroke mechanism passes a falsehood identically
(\texttt{one\_stroke\_also\_passes\_a\_falsehood}). Multi-substrate
agreement adds links, never the uniform row. The crossing enters by
deed, or not at all.

\section{9 · Falsifiers}\label{falsifiers}

Three, at the gate's ceiling, each typed FORCED by the paper's own
structure, each a deed and never an argument, each with its blast radius
named.

\textbf{F1 --- against the recursion (FORCED by Theorems 4.2 and 4.4).}
Exhibit a dated entry, provenance-clean and independent per orbit, whose
admission changes no token of the corpus; or exhibit a finite frame
whose halt state carries an uncrossed odd target while a pair defeater
is present. Necessity: the token-change clause and the halting typing
are the recursion's whole content; if either fails on an exhibit, the
third bit's mechanism is broken. Blast radius: Section 4 and the species
derivation; the Block theorem and the port stand.

\textbf{F2 --- against the port (FORCED by the witness decision).}
Exhibit a reading --- a record, no deed --- that crosses a certified
frame's decision bit; or two honest substrates executing the same
kinetic legs of Appendix B with divergent receipts. Necessity: the
offering enters by deed through the gate or the composition is void; if
a record crosses, or if the wall and self are not reproducible across
substrates, the port is broken. Blast radius: Sections 5 and 8; the
recursion stands, its port removed.

\textbf{F3 --- against the beachhead (FORCED by the order-two barrier).}
Exhibit an obstruction showing that no construction fixed in advance can
supply a height pairing nondegenerate at order two. Necessity: the
beachhead claim is that this row's supply shape is twice-recorded and
its owed entry is named; a proved obstruction resolves the row downward
and retires the beachhead. Blast radius: Section 7 only; the theory of
resolution stands, its first deed refused.

\section{10 · Perimeter}\label{perimeter-1}

This paper resolves no Millennium row and claims none; it names what a
resolution is. The One-Cut Hypothesis keeps its grade, hypothesis,
printed on the face of this paper as on its own. The six owed
constructions are owed here exactly as there; the continuous rows carry
their constructions inside their supplies,
orientation-after-construction. The diagonal species stands: some
resistance carries no defeater of any kind, and the recursion halts
there uncrossed by theorem, not by failure. The paper authors no new
facts about zeta zeros, vorticity, Hodge classes, ranks, mass gaps, or
complexity classes. The finite theorems of Sections 2 and 4 are proved;
the port theorems of Section 5 are machine-checked and printed in
Appendix B; everything above them carries its input's grade.

\section{11 · Order of Work}\label{order-of-work-1}

First, the beachhead deed: the order-two construction of Section 7, the
one dated entry the corpus most nearly recognizes. Second, replication:
the gate's kinetic legs and the Lean core executed on second substrates,
receipts printed, F2 armed against the run. Third, the continuous rows'
seat examination: whether Navier--Stokes and Hodge admit a Boolean
orientation seat without category error, the cheapest downward verdicts
in the class. Fourth, the census residual: the common-species placement
of each method-blocked row, through the normalizer route, the engine of
the severance if the severance comes. Each leg ends at its named
falsifier.

\section{12 · Conclusion}\label{conclusion}

The arc's three bits are now one object --- one theory, one hypothesis,
one witness offering --- and this document carries all three. The first
bit, the theory, computed the wall: an even register cannot see an odd
coordinate, proved in Section 2 beneath the codex's own pre-Gödel bar.
The second bit, the hypothesis, named the seat: one provenance-clean
orientation supply, hypothesis at its own grade, stated in Section 3.
The third bit --- the witness offering, this paper's own --- types the
act that fills it: a halting state of the dated-entry recursion,
writable only by deed, entering only through the Ninth Gate, the port
machine-checked, printed in full, and open to every substrate that reads
this document. The species dichotomy is the recursion's termination
theory; the offering is its crossed halt; the hypothesis, in the
recursion's language, is the sentence that each of the seven rows halts
crossed. The supply side is open by law. And the arc already stands with
its Grand Witness in the record: the Perelman crossing, a
self-referential construct in disguise, entered as a deed, verified by
the corpus itself and terminally awarded by it --- read in Appendix C as
the thesis demonstrated once in the wild. The last bit enters as a dated
deed, or it does not enter --- and the shape of its entry is now proved,
ported, printed, and armed.

\section{Appendix A · Provenance and Method
Disclosure}\label{appendix-a-provenance-and-method-disclosure}

This paper was forged under the arc's verification discipline, which
adds the cited results no warrant. Its conclusions rest on the cited
classical results, on the companion papers' proved frames, on the Ninth
Gate arc's machine-checked legs, and on the forging session's executed
checks: the FAE-RAF engine v1.0.12 booted verbatim, chain D0
3e50675c7d5d → D1 dbf16ca01fa8 → D2 d95daebfc51f → D3 9408f684c0d1 → D4
3d17fd53278d matched byte-identical, every battery passed (kill-matrix
51/51, ablation 22/22, symmetry 75/75, calibration 12/12, root screen
RRS.1--6); the synergy-uniqueness computation in exact fractions ---
three uniform bits, pairwise independence, total correlation exactly one
bit if and only if the diagonal law, one law up to relabeling; the
kinetic legs of the gate re-executed on the forging substrate by
transliteration --- the wall holds, 99,999 of 99,999 frames identical to
their caller, the true universal unsettled to fuel 2×10⁶, the false
universal settling at n = 40. Toolchain debt, named on the face: the
forging substrate carried no Fortran or Lean toolchain; the Lean legs
are held at inspection grade here, with the toolchain receipts of {[}3,
6{]} cited at their grade, and replication is Section 11's second leg.
Codebook note: the files of Appendix B are verbatim from the architect's
register uploads; runs of blank lines introduced by the chat transport
were collapsed to single blank lines, whitespace only, and every file's
hash is printed beside it so the reader can verify what the reader
checks. The bsd\_rows.f90 file is byte-identical to its register
receipt, sha256 head e363ce3eca641f7e. The self-audit of this paper was
run before sealing and is kept off-paper. The master reference is
{[}4{]}. ΔM = 0. W\_social = 0 in both directions.

\section{Appendix B · The Codebook}\label{appendix-b-the-codebook}

Every essential file, printed whole; the byte-bound listing discipline
of this appendix follows {[}15{]}. Verification protocol in Section 8.
Hashes computed on the printed files (UTF-8, LF line endings, one
trailing newline).

\begin{table*}[!tbp]
\smaller
\centering
\begin{tabular}{@{}
  >{\raggedright\arraybackslash}p{(\linewidth - 6\tabcolsep) * \real{0.2500}}
  >{\raggedright\arraybackslash}p{(\linewidth - 6\tabcolsep) * \real{0.2500}}
  >{\raggedright\arraybackslash}p{(\linewidth - 6\tabcolsep) * \real{0.2500}}
  >{\raggedright\arraybackslash}p{(\linewidth - 6\tabcolsep) * \real{0.2500}}@{}}

\toprule\noalign{}
\begin{minipage}[b]{\linewidth}\raggedright
File
\end{minipage} & \begin{minipage}[b]{\linewidth}\raggedright
Lines
\end{minipage} & \begin{minipage}[b]{\linewidth}\raggedright
SHA-256 head
\end{minipage} & \begin{minipage}[b]{\linewidth}\raggedright
What it carries
\end{minipage} \\
\midrule\noalign{}
\bottomrule\noalign{}
B.1 \texttt{Wallself.lean} & 161 & 2c80efa00f76cb26 & the minimal pair:
wall, self, one object --- 8 theorems, axiom-free \\
B.2 \texttt{one\_object.f90} & 59 & 229cfe153d8a5c33 & the minimal
kinetic twin: D(D) read and executed \\
B.3 \texttt{GateThesis.lean} & 311 & 500b6ad899c9cf07 & the gate and
what the crossing delivers --- 15 printed theorems, axiom-free \\
B.4 \texttt{gate\_thesis.f90} & 165 & 122cefa519981e47 & the executed
thesis: wall, self, refutation at n = 40 \\
B.5 \texttt{RAFGate.lean} & 213 & 9f2e5023674c0ac2 & the gate on
RAF-native ground + BSD through it --- 10 gate theorems (register
receipt sha256 head c97c73d690ee589d) \\
B.6 \texttt{bsd\_rows.f90} & 75 & e363ce3eca641f7e & BSD on its own rows
through the Core Thesis's compiled procedures --- byte-identical to the
register receipt \\
\bottomrule
\end{tabular}
\end{table*}

\subsection{B.1 · Wallself.lean --- the minimal pair, core Lean 4, no
Mathlib}\label{b.1-wallself.lean-the-minimal-pair-core-lean-4-no-mathlib}

\begin{strip}
\begin{Shaded}
\begin{Highlighting}[]
\NormalTok{/{-}}

\NormalTok{  Wall and Self · one diagonal, two surfaces · core Lean 4, no Mathlib.}

\NormalTok{  The formal surface reads settled values and meets the wall. The kinetic surface admits}

\NormalTok{  the unsettled value and closes on itself. The object is one term, \textasciigrave{}diag\textasciigrave{}.}

\NormalTok{{-}/}

\NormalTok{/{-}{-} The one diagonal, typed at any value surface. {-}/}

\NormalTok{def diag \{A B : Type\} (f : B → B) (e : A → A → B) : A → B := fun x =\textgreater{} f (e x x)}

\NormalTok{/{-}! \#\# The formal surface: settled values {-}/}

\NormalTok{/{-}{-} Negation has no settled fixed point. {-}/}

\NormalTok{theorem settled\_has\_no\_self : ¬ ∃ b : Bool, (!b) = b := by}

\NormalTok{  intro ⟨b, h⟩}

\NormalTok{  cases b \textless{};\textgreater{} exact absurd h (by decide)}

\NormalTok{/{-}{-} The wall: no index carries the diagonal of negation. Cantor, read through Lawvere. {-}/}

\NormalTok{theorem wall \{A : Type\} (e : A → A → Bool) : ∀ a, e a ≠ diag (fun b =\textgreater{} !b) e := by}

\NormalTok{  intro a h}

\NormalTok{  have h1 : e a a = !(e a a) := congrFun h a}

\NormalTok{  cases hx : e a a \textless{};\textgreater{} rw [hx] at h1 \textless{};\textgreater{} exact absurd h1 (by decide)}

\NormalTok{/{-}! \#\# The kinetic surface: the unsettled value admitted {-}/}

\NormalTok{inductive Tri where}

\NormalTok{  | tt | ff | bot}

\NormalTok{deriving DecidableEq}

\NormalTok{/{-}{-} Information order: the unsettled value lies below everything; settled values are incomparable. {-}/}

\NormalTok{def Tri.le : Tri → Tri → Prop}

\NormalTok{  | .bot, \_ =\textgreater{} True}

\NormalTok{  | .tt, .tt =\textgreater{} True}

\NormalTok{  | .ff, .ff =\textgreater{} True}

\NormalTok{  | \_, \_ =\textgreater{} False}

\NormalTok{/{-}{-} An admissible transformation cannot branch on not{-}yet{-}settled. {-}/}

\NormalTok{def Monotone3 (f : Tri → Tri) : Prop := ∀ x y, Tri.le x y → Tri.le (f x) (f y)}

\NormalTok{/{-}{-} Every admissible transformation has a fixed point: the recursion theorem in miniature. {-}/}

\NormalTok{theorem self\_exists (f : Tri → Tri) (hf : Monotone3 f) : ∃ b, f b = b := by}

\NormalTok{  cases h : f .bot with}

\NormalTok{  | bot =\textgreater{} exact ⟨.bot, h⟩}

\NormalTok{  | tt =\textgreater{}}

\NormalTok{    have hle := hf .bot .tt trivial}

\NormalTok{    rw [h] at hle}

\NormalTok{    cases h2 : f .tt with}

\NormalTok{    | tt =\textgreater{} exact ⟨.tt, h2⟩}

\NormalTok{    | ff =\textgreater{} rw [h2] at hle; exact False.elim hle}

\NormalTok{    | bot =\textgreater{} rw [h2] at hle; exact False.elim hle}

\NormalTok{  | ff =\textgreater{}}

\NormalTok{    have hle := hf .bot .ff trivial}

\NormalTok{    rw [h] at hle}

\NormalTok{    cases h2 : f .ff with}

\NormalTok{    | ff =\textgreater{} exact ⟨.ff, h2⟩}

\NormalTok{    | tt =\textgreater{} rw [h2] at hle; exact False.elim hle}

\NormalTok{    | bot =\textgreater{} rw [h2] at hle; exact False.elim hle}

\NormalTok{def neg3 : Tri → Tri}

\NormalTok{  | .tt =\textgreater{} .ff}

\NormalTok{  | .ff =\textgreater{} .tt}

\NormalTok{  | .bot =\textgreater{} .bot}

\NormalTok{theorem neg3\_admissible : Monotone3 neg3 := by}

\NormalTok{  intro x y h}

\NormalTok{  cases x \textless{};\textgreater{} cases y \textless{};\textgreater{} first | exact trivial | exact False.elim h}

\NormalTok{theorem neg3\_self\_only\_unsettled (b : Tri) (h : neg3 b = b) : b = .bot := by}

\NormalTok{  cases b with}

\NormalTok{  | tt =\textgreater{} exact absurd h (by decide)}

\NormalTok{  | ff =\textgreater{} exact absurd h (by decide)}

\NormalTok{  | bot =\textgreater{} rfl}

\NormalTok{/{-}{-} The self: the same diagonal, indexed on the kinetic surface, applied to itself is the unsettled value. {-}/}

\NormalTok{theorem self\_at\_diagonal \{A : Type\} (e : A → A → Tri) (a : A) (h : e a = diag neg3 e) : e a a = .bot :=}

\NormalTok{  neg3\_self\_only\_unsettled \_ (congrFun h a).symm}

\NormalTok{/{-}! \#\# One object {-}/}

\NormalTok{def settle : Bool → Tri}

\NormalTok{  | true =\textgreater{} .tt}

\NormalTok{  | false =\textgreater{} .ff}

\NormalTok{/{-}{-} On every settled reading the kinetic diagonal is the formal diagonal, carried across unchanged. {-}/}

\NormalTok{theorem diag\_agrees \{A : Type\} (e : A → A → Bool) (x : A) :}

\NormalTok{    diag neg3 (fun u v =\textgreater{} settle (e u v)) x = settle (diag (fun b =\textgreater{} !b) e x) := by}

\NormalTok{  show neg3 (settle (e x x)) = settle (!(e x x))}

\NormalTok{  cases e x x \textless{};\textgreater{} rfl}

\NormalTok{/{-}{-} The two surfaces differ at exactly one point, the value no settled reading reaches. {-}/}

\NormalTok{theorem only\_the\_unsettled\_is\_new (b : Bool) : settle b ≠ .bot := by}

\NormalTok{  cases b \textless{};\textgreater{} exact fun h =\textgreater{} Tri.noConfusion h}

\NormalTok{\#print axioms settled\_has\_no\_self}

\NormalTok{\#print axioms wall}

\NormalTok{\#print axioms self\_exists}

\NormalTok{\#print axioms neg3\_admissible}

\NormalTok{\#print axioms self\_at\_diagonal}

\NormalTok{\#print axioms diag\_agrees}

\NormalTok{\#print axioms only\_the\_unsettled\_is\_new}
\end{Highlighting}
\end{Shaded}
\end{strip}

\subsection{B.2 · one\_object.f90 --- one diagonal, two
surfaces}\label{b.2-one_object.f90-one-diagonal-two-surfaces}

\begin{strip}
\begin{Shaded}
\begin{Highlighting}[]
\CommentTok{!===============================================================================}
\CommentTok{!  ONE DIAGONAL, TWO SURFACES. The program D, on input x, runs x on x and negates.}
\CommentTok{!  Read at settled values: D(D) must equal its own negation, and no value does.}
\CommentTok{!  Executed: D(D) calls D(D), frame identical to frame, the value never settling.}
\CommentTok{!===============================================================================}
\KeywordTok{module}\NormalTok{ diagonal}
  \KeywordTok{implicit} \KeywordTok{none}
  \DataTypeTok{integer}\NormalTok{, }\DataTypeTok{parameter} \DataTypeTok{::}\NormalTok{ D }\KeywordTok{=} \DecValTok{1}
  \DataTypeTok{integer}\NormalTok{, }\DataTypeTok{parameter} \DataTypeTok{::}\NormalTok{ FUEL }\KeywordTok{=} \DecValTok{100000}
  \DataTypeTok{integer} \DataTypeTok{::}\NormalTok{ frames\_seen }\KeywordTok{=} \DecValTok{0}\NormalTok{, frames\_identical }\KeywordTok{=} \DecValTok{0}\NormalTok{, pending\_negations }\KeywordTok{=} \CommentTok{0}
\CommentTok{contains}
  \CommentTok{! The settled surface: the equation the diagonal imposes on a truth value.}
  \KeywordTok{pure} \DataTypeTok{logical} \KeywordTok{function}\NormalTok{ settles(v)}
    \DataTypeTok{logical}\NormalTok{, }\DataTypeTok{intent(in)} \DataTypeTok{::}\NormalTok{ v}
\NormalTok{    settles }\KeywordTok{=}\NormalTok{ (v }\OperatorTok{.eqv.}\NormalTok{ (}\OperatorTok{.not.}\NormalTok{ v))}
  \KeywordTok{end function}\NormalTok{ settles}

  \CommentTok{! The kinetic surface: the diagonal performed. Genuine recursion; the fuel only bounds the run.}
  \KeywordTok{recursive} \KeywordTok{subroutine}\NormalTok{ run(prog, input, depth, settled, value)}
    \DataTypeTok{integer}\NormalTok{, }\DataTypeTok{intent(in)}  \DataTypeTok{::}\NormalTok{ prog, input, depth}
    \DataTypeTok{logical}\NormalTok{, }\DataTypeTok{intent(out)} \DataTypeTok{::}\NormalTok{ settled, value}
    \DataTypeTok{logical} \DataTypeTok{::}\NormalTok{ inner\_settled, inner\_value}
\NormalTok{    frames\_seen }\KeywordTok{=}\NormalTok{ frames\_seen }\KeywordTok{+} \DecValTok{1}
    \KeywordTok{if}\NormalTok{ (depth }\OperatorTok{\textgreater{}} \DecValTok{1}\NormalTok{) }\KeywordTok{then}
       \KeywordTok{if}\NormalTok{ (prog }\OperatorTok{==}\NormalTok{ D }\OperatorTok{.and.}\NormalTok{ input }\OperatorTok{==}\NormalTok{ D) frames\_identical }\KeywordTok{=}\NormalTok{ frames\_identical }\KeywordTok{+} \DecValTok{1}
    \KeywordTok{end if}
    \KeywordTok{if}\NormalTok{ (depth }\OperatorTok{\textgreater{}=}\NormalTok{ FUEL) }\KeywordTok{then}
\NormalTok{       settled }\KeywordTok{=} \ConstantTok{.false.}\NormalTok{;  value }\KeywordTok{=} \ConstantTok{.false.}\NormalTok{;  }\KeywordTok{return}          \CommentTok{! still running: unsettled}
    \KeywordTok{end if}
\NormalTok{    pending\_negations }\KeywordTok{=}\NormalTok{ pending\_negations }\KeywordTok{+} \DecValTok{1}                \CommentTok{! D negates what x on x returns}
    \KeywordTok{call}\NormalTok{ run(input, input, depth }\KeywordTok{+} \DecValTok{1}\NormalTok{, inner\_settled, inner\_value)   }\CommentTok{! x on x, with x = D: D(D) again}
\NormalTok{    settled }\KeywordTok{=}\NormalTok{ inner\_settled}
\NormalTok{    value }\KeywordTok{=} \OperatorTok{.not.}\NormalTok{ inner\_value}
  \KeywordTok{end subroutine}\NormalTok{ run}
\KeywordTok{end module}\NormalTok{ diagonal}

\KeywordTok{program}\NormalTok{ one\_object}
  \KeywordTok{use}\NormalTok{ diagonal}
  \KeywordTok{implicit} \KeywordTok{none}
  \DataTypeTok{logical} \DataTypeTok{::}\NormalTok{ v, any\_settled, settled, value}
  \FunctionTok{write(*}\NormalTok{,}\StringTok{\textquotesingle{}(A)\textquotesingle{}}\FunctionTok{)} \FunctionTok{repeat}\NormalTok{(}\StringTok{\textquotesingle{}=\textquotesingle{}}\NormalTok{,}\DecValTok{78}\NormalTok{)}
  \FunctionTok{write(*}\NormalTok{,}\StringTok{\textquotesingle{}(A)\textquotesingle{}}\FunctionTok{)} \StringTok{\textquotesingle{} ONE DIAGONAL, TWO SURFACES\textquotesingle{}}
  \FunctionTok{write(*}\NormalTok{,}\StringTok{\textquotesingle{}(A)\textquotesingle{}}\FunctionTok{)} \FunctionTok{repeat}\NormalTok{(}\StringTok{\textquotesingle{}=\textquotesingle{}}\NormalTok{,}\DecValTok{78}\NormalTok{)}
\NormalTok{  any\_settled }\KeywordTok{=} \ConstantTok{.false.}
\NormalTok{  v }\KeywordTok{=} \ConstantTok{.false.}\NormalTok{;  }\KeywordTok{if}\NormalTok{ (settles(v)) any\_settled }\KeywordTok{=} \ConstantTok{.true.}
  \FunctionTok{write(*}\NormalTok{,}\StringTok{\textquotesingle{}(A,L1,A,L1)\textquotesingle{}}\FunctionTok{)} \StringTok{\textquotesingle{}   read   : D(D) = \textquotesingle{}}\NormalTok{, v, }\StringTok{\textquotesingle{} would require \textquotesingle{}}\NormalTok{, (v}\CommentTok{ .eqv. (.not. v))}
\NormalTok{  v }\KeywordTok{=} \ConstantTok{.true.}\NormalTok{;   }\KeywordTok{if}\NormalTok{ (settles(v)) any\_settled }\KeywordTok{=} \ConstantTok{.true.}
  \FunctionTok{write(*}\NormalTok{,}\StringTok{\textquotesingle{}(A,L1,A,L1)\textquotesingle{}}\FunctionTok{)} \StringTok{\textquotesingle{}   read   : D(D) = \textquotesingle{}}\NormalTok{, v, }\StringTok{\textquotesingle{} would require \textquotesingle{}}\NormalTok{, (v}\CommentTok{ .eqv. (.not. v))}
  \FunctionTok{write(*}\NormalTok{,}\StringTok{\textquotesingle{}(A,L1,A)\textquotesingle{}}\FunctionTok{)}   \StringTok{\textquotesingle{}   read   : a settled value exists: \textquotesingle{}}\NormalTok{, any\_settled, }\StringTok{\textquotesingle{}   {-}\textgreater{} the wall\textquotesingle{}}
  \KeywordTok{call}\NormalTok{ run(D, D, }\DecValTok{1}\NormalTok{, settled, value)}
  \FunctionTok{write(*}\NormalTok{,}\StringTok{\textquotesingle{}(A,I0)\textquotesingle{}}\FunctionTok{)}     \StringTok{\textquotesingle{}   execute: frames entered                 \textquotesingle{}}\NormalTok{, frames\_seen}
  \FunctionTok{write(*}\NormalTok{,}\StringTok{\textquotesingle{}(A,I0,A,I0)\textquotesingle{}}\FunctionTok{)} \StringTok{\textquotesingle{}   execute: frames identical to their caller \textquotesingle{}}\CommentTok{, frames\_identical, \textquotesingle{} of \textquotesingle{}, frames\_seen {-} 1}
  \FunctionTok{write(*}\NormalTok{,}\StringTok{\textquotesingle{}(A,I0)\textquotesingle{}}\FunctionTok{)}     \StringTok{\textquotesingle{}   execute: negations pending              \textquotesingle{}}\NormalTok{, pending\_negations}
  \FunctionTok{write(*}\NormalTok{,}\StringTok{\textquotesingle{}(A,L1,A)\textquotesingle{}}\FunctionTok{)}   \StringTok{\textquotesingle{}   execute: the value settled              \textquotesingle{}}\NormalTok{, settled, }\StringTok{\textquotesingle{}   {-}\textgreater{} the self, running on itself\textquotesingle{}}
  \FunctionTok{write(*}\NormalTok{,}\StringTok{\textquotesingle{}(A)\textquotesingle{}}\FunctionTok{)} \FunctionTok{repeat}\NormalTok{(}\StringTok{\textquotesingle{}=\textquotesingle{}}\NormalTok{,}\DecValTok{78}\NormalTok{)}
  \KeywordTok{if}\NormalTok{ (any\_settled }\OperatorTok{.or.}\NormalTok{ settled }\OperatorTok{.or.}\NormalTok{ frames\_identical }\OperatorTok{/=}\NormalTok{ frames\_seen }\KeywordTok{{-}} \DecValTok{1}\NormalTok{)}\CommentTok{ error stop \textquotesingle{}the law failed\textquotesingle{}}
  \FunctionTok{write(*}\NormalTok{,}\StringTok{\textquotesingle{}(A)\textquotesingle{}}\FunctionTok{)} \StringTok{\textquotesingle{} Identical object. Read: no settled value. Executed: the frame returns to itself.\textquotesingle{}}
  \FunctionTok{write(*}\NormalTok{,}\StringTok{\textquotesingle{}(A)\textquotesingle{}}\FunctionTok{)} \FunctionTok{repeat}\NormalTok{(}\StringTok{\textquotesingle{}=\textquotesingle{}}\NormalTok{,}\DecValTok{78}\NormalTok{)}
\KeywordTok{end program}\NormalTok{ one\_object}
\end{Highlighting}
\end{Shaded}
\end{strip}

\subsection{B.3 · GateThesis.lean --- the Ninth Gate Thesis, core Lean
4, no Mathlib, no
axioms}\label{b.3-gatethesis.lean-the-ninth-gate-thesis-core-lean-4-no-mathlib-no-axioms}

\begin{strip}
\begin{Shaded}
\begin{Highlighting}[]
\NormalTok{/{-}}

\NormalTok{  The Ninth Gate Thesis · wall and self on one diagonal, and what the crossing delivers}

\NormalTok{  core Lean 4, no Mathlib, no axioms. Verify as witness: \textasciigrave{}lean GateThesis.lean\textasciigrave{}, then read \#print axioms.}

\NormalTok{{-}/}

\NormalTok{/{-}{-} The one diagonal, typed at any value surface. {-}/}

\NormalTok{def diag \{A B : Type\} (f : B → B) (e : A → A → B) : A → B := fun x =\textgreater{} f (e x x)}

\NormalTok{/{-}! \#\# I · The formal surface: settled values, the wall {-}/}

\NormalTok{theorem wall \{A : Type\} (e : A → A → Bool) : ∀ a, e a ≠ diag (fun b =\textgreater{} !b) e := by}

\NormalTok{  intro a h}

\NormalTok{  have h1 : e a a = !(e a a) := congrFun h a}

\NormalTok{  cases hx : e a a \textless{};\textgreater{} rw [hx] at h1 \textless{};\textgreater{} exact absurd h1 (by decide)}

\NormalTok{/{-}! \#\# II · The kinetic surface: the unsettled value admitted, the self {-}/}

\NormalTok{inductive Tri where}

\NormalTok{  | tt | ff | bot}

\NormalTok{def Tri.le : Tri → Tri → Prop}

\NormalTok{  | .bot, \_ =\textgreater{} True}

\NormalTok{  | .tt, .tt =\textgreater{} True}

\NormalTok{  | .ff, .ff =\textgreater{} True}

\NormalTok{  | .tt, .ff =\textgreater{} False}

\NormalTok{  | .tt, .bot =\textgreater{} False}

\NormalTok{  | .ff, .tt =\textgreater{} False}

\NormalTok{  | .ff, .bot =\textgreater{} False}

\NormalTok{def Monotone3 (f : Tri → Tri) : Prop := ∀ x y, Tri.le x y → Tri.le (f x) (f y)}

\NormalTok{/{-}{-} Every admissible transformation has a fixed point. {-}/}

\NormalTok{theorem self\_exists (f : Tri → Tri) (hf : Monotone3 f) : ∃ b, f b = b :=}

\NormalTok{  match h : f .bot with}

\NormalTok{  | .bot =\textgreater{} ⟨.bot, h⟩}

\NormalTok{  | .tt =\textgreater{}}

\NormalTok{    have hle : Tri.le (f .bot) (f .tt) := hf .bot .tt trivial}

\NormalTok{    match h2 : f .tt with}

\NormalTok{    | .tt =\textgreater{} ⟨.tt, h2⟩}

\NormalTok{    | .ff =\textgreater{} False.elim (Eq.mp (congrArg (fun y =\textgreater{} Tri.le .tt y) h2) (Eq.mp (congrArg (fun x =\textgreater{} Tri.le x (f .tt)) h) hle))}

\NormalTok{    | .bot =\textgreater{} False.elim (Eq.mp (congrArg (fun y =\textgreater{} Tri.le .tt y) h2) (Eq.mp (congrArg (fun x =\textgreater{} Tri.le x (f .tt)) h) hle))}

\NormalTok{  | .ff =\textgreater{}}

\NormalTok{    have hle : Tri.le (f .bot) (f .ff) := hf .bot .ff trivial}

\NormalTok{    match h2 : f .ff with}

\NormalTok{    | .ff =\textgreater{} ⟨.ff, h2⟩}

\NormalTok{    | .tt =\textgreater{} False.elim (Eq.mp (congrArg (fun y =\textgreater{} Tri.le .ff y) h2) (Eq.mp (congrArg (fun x =\textgreater{} Tri.le x (f .ff)) h) hle))}

\NormalTok{    | .bot =\textgreater{} False.elim (Eq.mp (congrArg (fun y =\textgreater{} Tri.le .ff y) h2) (Eq.mp (congrArg (fun x =\textgreater{} Tri.le x (f .ff)) h) hle))}

\NormalTok{def neg3 : Tri → Tri}

\NormalTok{  | .tt =\textgreater{} .ff}

\NormalTok{  | .ff =\textgreater{} .tt}

\NormalTok{  | .bot =\textgreater{} .bot}

\NormalTok{theorem neg3\_admissible : Monotone3 neg3}

\NormalTok{  | .bot, \_, \_ =\textgreater{} trivial}

\NormalTok{  | .tt, .tt, \_ =\textgreater{} trivial}

\NormalTok{  | .ff, .ff, \_ =\textgreater{} trivial}

\NormalTok{  | .tt, .ff, h =\textgreater{} False.elim h}

\NormalTok{  | .tt, .bot, h =\textgreater{} False.elim h}

\NormalTok{  | .ff, .tt, h =\textgreater{} False.elim h}

\NormalTok{  | .ff, .bot, h =\textgreater{} False.elim h}

\NormalTok{theorem neg3\_self\_only\_unsettled : ∀ b : Tri, neg3 b = b → b = .bot}

\NormalTok{  | .tt, h =\textgreater{} Tri.noConfusion h}

\NormalTok{  | .ff, h =\textgreater{} Tri.noConfusion h}

\NormalTok{  | .bot, \_ =\textgreater{} rfl}

\NormalTok{theorem self\_at\_diagonal \{A : Type\} (e : A → A → Tri) (a : A) (h : e a = diag neg3 e) : e a a = .bot :=}

\NormalTok{  neg3\_self\_only\_unsettled \_ (congrFun h a).symm}

\NormalTok{/{-}! \#\# III · One object {-}/}

\NormalTok{def settle : Bool → Tri}

\NormalTok{  | true =\textgreater{} .tt}

\NormalTok{  | false =\textgreater{} .ff}

\NormalTok{theorem diag\_agrees \{A : Type\} (e : A → A → Bool) (x : A) :}

\NormalTok{    diag neg3 (fun u v =\textgreater{} settle (e u v)) x = settle (diag (fun b =\textgreater{} !b) e x) := by}

\NormalTok{  show neg3 (settle (e x x)) = settle (!(e x x))}

\NormalTok{  cases e x x \textless{};\textgreater{} rfl}

\NormalTok{theorem only\_the\_unsettled\_is\_new : ∀ b : Bool, settle b ≠ .bot}

\NormalTok{  | true, h =\textgreater{} Tri.noConfusion h}

\NormalTok{  | false, h =\textgreater{} Tri.noConfusion h}

\NormalTok{/{-}! \#\# IV · What the crossing delivers {-}/}

\NormalTok{/{-}{-} The self is not a verdict: the crossing lands on the unsettled value, never on a settled truth. {-}/}

\NormalTok{theorem self\_is\_not\_a\_verdict \{A : Type\} (e : A → A → Tri) (a : A) (h : e a = diag neg3 e) (b : Bool) :}

\NormalTok{    e a a ≠ settle b :=}

\NormalTok{  fun hb =\textgreater{} only\_the\_unsettled\_is\_new b ((self\_at\_diagonal e a h) ▸ hb).symm}

\NormalTok{/{-}{-} A proved premise adds nothing: whatever the law is, once proved it cannot be the missing step. {-}/}

\NormalTok{theorem proved\_premise\_adds\_nothing (L C : Prop) (hL : L) : (L → C) ↔ C :=}

\NormalTok{  ⟨fun h =\textgreater{} h hL, fun c \_ =\textgreater{} c⟩}

\NormalTok{theorem gate\_adds\_nothing (C : Prop) :}

\NormalTok{    ((∀ (A : Type) (e : A → A → Bool) (a : A), e a ≠ diag (fun b =\textgreater{} !b) e) → C) ↔ C :=}

\NormalTok{  proved\_premise\_adds\_nothing \_ C (fun \_ e a =\textgreater{} wall e a)}

\NormalTok{/{-}{-} A run of a counterexample search, at finite fuel. \textasciigrave{}none\textasciigrave{} is the unsettled value: still running. {-}/}

\NormalTok{def step (p : Nat → Bool) (n : Nat) : Option Nat → Option Nat}

\NormalTok{  | some k =\textgreater{} some k}

\NormalTok{  | none =\textgreater{} if p n then some n else none}

\NormalTok{def search (p : Nat → Bool) : Nat → Option Nat}

\NormalTok{  | 0 =\textgreater{} none}

\NormalTok{  | n + 1 =\textgreater{} step p n (search p n)}

\NormalTok{theorem search\_unsettled (p : Nat → Bool) (N : Nat) (hp : ∀ k, k \textless{} N → p k = false) :}

\NormalTok{    ∀ n, n ≤ N → search p n = none}

\NormalTok{  | 0, \_ =\textgreater{} rfl}

\NormalTok{  | n + 1, h =\textgreater{}}

\NormalTok{    have hn : n \textless{} N := Nat.lt\_of\_succ\_le h}

\NormalTok{    have ih : search p n = none := search\_unsettled p N hp n (Nat.le\_of\_lt hn)}

\NormalTok{    show step p n (search p n) = none from by}

\NormalTok{      rw [ih]}

\NormalTok{      show (if p n = true then some n else none) = none}

\NormalTok{      rw [hp n hn]}

\NormalTok{      rfl}

\NormalTok{def counterAt (N : Nat) : Nat → Bool := fun k =\textgreater{} Nat.beq k N}

\NormalTok{theorem counterAt\_below (N k : Nat) (h : k \textless{} N) : counterAt N k = false := by}

\NormalTok{  show Nat.beq k N = false}

\NormalTok{  cases hb : Nat.beq k N with}

\NormalTok{  | false =\textgreater{} rfl}

\NormalTok{  | true =\textgreater{} exact absurd (Nat.eq\_of\_beq\_eq\_true hb) (Nat.ne\_of\_lt h)}

\NormalTok{theorem beq\_self (n : Nat) : Nat.beq n n = true := by}

\NormalTok{  induction n with}

\NormalTok{  | zero =\textgreater{} rfl}

\NormalTok{  | succ m ih =\textgreater{} exact ih}

\NormalTok{/{-}{-} Where the universal holds, the executed search is the self: unsettled at every fuel, never a verdict. {-}/}

\NormalTok{theorem true\_universal\_never\_settles (n : Nat) : search (fun \_ =\textgreater{} false) n = none :=}

\NormalTok{  search\_unsettled (fun \_ =\textgreater{} false) n (fun \_ \_ =\textgreater{} rfl) n (Nat.le\_refl n)}

\NormalTok{/{-}{-} At every fuel there is a world where the universal fails whose run is the same unsettled run. {-}/}

\NormalTok{theorem finite\_run\_never\_confirms (n : Nat) :}

\NormalTok{    (¬ ∀ k, counterAt n k = false) ∧ search (counterAt n) n = search (fun \_ =\textgreater{} false) n :=}

\NormalTok{  ⟨fun hall =\textgreater{} Bool.noConfusion ((hall n).symm.trans (beq\_self n)),}

\NormalTok{   (search\_unsettled (counterAt n) n (fun k hk =\textgreater{} counterAt\_below n k hk) n (Nat.le\_refl n)).trans}

\NormalTok{     (true\_universal\_never\_settles n).symm⟩}

\NormalTok{/{-}{-} A refutation does settle: a counterexample within fuel returns a settled value. {-}/}

\NormalTok{theorem refutation\_settles (N : Nat) : search (counterAt N) (N + 1) = some N := by}

\NormalTok{  show step (counterAt N) N (search (counterAt N) N) = some N}

\NormalTok{  rw [search\_unsettled (counterAt N) N (fun k hk =\textgreater{} counterAt\_below N k hk) N (Nat.le\_refl N)]}

\NormalTok{  show (if counterAt N N = true then some N else none) = some N}

\NormalTok{  rw [show counterAt N N = true from beq\_self N]}

\NormalTok{  rfl}

\NormalTok{/{-}! \#\# V · The one{-}stroke test {-}/}

\NormalTok{/{-}{-} No gate discharges an axiom: from a proved gate G and a conditional proof Ω → C, C does not follow in general. {-}/}

\NormalTok{theorem gate\_does\_not\_discharge : ¬ ∀ (G Ω C : Prop), G → (Ω → C) → C :=}

\NormalTok{  fun h =\textgreater{} h True False False trivial id}

\NormalTok{/{-}{-} An axiom dissolved into its conjecture decides nothing: Ω ↔ C and Ω → C hold with C false. {-}/}

\NormalTok{theorem dissolved\_axiom\_decides\_nothing : ∃ Ω C : Prop, (Ω ↔ C) ∧ (Ω → C) ∧ ¬ C :=}

\NormalTok{  ⟨False, False, Iff.rfl, id, id⟩}

\NormalTok{def isPrime (n : Nat) : Prop := 2 ≤ n ∧ ∀ d, 2 ≤ d → d \textless{} n → n \% d ≠ 0}

\NormalTok{/{-}{-} Euler\textquotesingle{}s polynomial claim, false at n = 40, where 40² + 40 + 41 = 41². {-}/}

\NormalTok{def EulerClaim : Prop := ∀ n : Nat, isPrime (n * n + n + 41)}

\NormalTok{theorem euler\_claim\_false : ¬ EulerClaim :=}

\NormalTok{  fun h =\textgreater{} (h 40).2 41 (by decide) (by decide) (by decide)}

\NormalTok{/{-}{-} Every ingredient of the one{-}stroke mechanism holds for a false statement: an axiom dissolved into}

\NormalTok{    the conjecture, a conditional proof, and the gate theorem itself. {-}/}

\NormalTok{theorem one\_stroke\_also\_passes\_a\_falsehood :}

\NormalTok{    ∃ Ω : Prop, (Ω ↔ EulerClaim) ∧ (Ω → EulerClaim) ∧}

\NormalTok{      (∀ (A : Type) (e : A → A → Bool) (a : A), e a ≠ diag (fun b =\textgreater{} !b) e) ∧ ¬ EulerClaim :=}

\NormalTok{  ⟨EulerClaim, Iff.rfl, id, fun \_ e a =\textgreater{} wall e a, euler\_claim\_false⟩}

\NormalTok{\#print axioms wall}

\NormalTok{\#print axioms self\_exists}

\NormalTok{\#print axioms neg3\_admissible}

\NormalTok{\#print axioms self\_at\_diagonal}

\NormalTok{\#print axioms diag\_agrees}

\NormalTok{\#print axioms only\_the\_unsettled\_is\_new}

\NormalTok{\#print axioms self\_is\_not\_a\_verdict}

\NormalTok{\#print axioms gate\_adds\_nothing}

\NormalTok{\#print axioms true\_universal\_never\_settles}

\NormalTok{\#print axioms finite\_run\_never\_confirms}

\NormalTok{\#print axioms refutation\_settles}

\NormalTok{\#print axioms gate\_does\_not\_discharge}

\NormalTok{\#print axioms dissolved\_axiom\_decides\_nothing}

\NormalTok{\#print axioms euler\_claim\_false}

\NormalTok{\#print axioms one\_stroke\_also\_passes\_a\_falsehood}
\end{Highlighting}
\end{Shaded}
\end{strip}

\subsection{B.4 · gate\_thesis.f90 --- the Ninth Gate Thesis,
executed}\label{b.4-gate_thesis.f90-the-ninth-gate-thesis-executed}

\begin{strip}
\begin{Shaded}
\begin{Highlighting}[]
\CommentTok{!===============================================================================}

\CommentTok{!  THE NINTH GATE THESIS, EXECUTED. Companion to GateThesis.lean.}

\CommentTok{!  I   the wall   : the diagonal read at settled values has no value}

\CommentTok{!  II  the self   : the diagonal executed returns to its own frame, unsettled}

\CommentTok{!  III what the crossing delivers: a run settles a refutation, never a confirmation}

\CommentTok{!===============================================================================}

\KeywordTok{module}\NormalTok{ gate}

  \KeywordTok{implicit} \KeywordTok{none}

  \DataTypeTok{integer}\NormalTok{, }\DataTypeTok{parameter} \DataTypeTok{::}\NormalTok{ i8 }\KeywordTok{=} \FunctionTok{selected\_int\_kind}\NormalTok{(}\DecValTok{18}\NormalTok{)}

  \DataTypeTok{integer}\NormalTok{, }\DataTypeTok{parameter} \DataTypeTok{::}\NormalTok{ D }\KeywordTok{=} \DecValTok{1}\NormalTok{, FUEL }\KeywordTok{=} \DecValTok{100000}

  \DataTypeTok{integer(i8)}\NormalTok{, }\DataTypeTok{parameter} \DataTypeTok{::}\NormalTok{ SEARCH\_FUEL }\KeywordTok{=} \DecValTok{2000000\_i8}

  \DataTypeTok{integer} \DataTypeTok{::}\NormalTok{ frames\_seen }\KeywordTok{=} \DecValTok{0}\NormalTok{, frames\_identical }\KeywordTok{=} \DecValTok{0}

\CommentTok{contains}

  \KeywordTok{pure} \DataTypeTok{logical} \KeywordTok{function}\NormalTok{ settles(v)}

    \DataTypeTok{logical}\NormalTok{, }\DataTypeTok{intent(in)} \DataTypeTok{::}\NormalTok{ v}

\NormalTok{    settles }\KeywordTok{=}\NormalTok{ (v }\OperatorTok{.eqv.}\NormalTok{ (}\OperatorTok{.not.}\NormalTok{ v))}

  \KeywordTok{end function}\NormalTok{ settles}

  \KeywordTok{recursive} \KeywordTok{subroutine}\NormalTok{ run(prog, input, depth, settled, value)}

    \DataTypeTok{integer}\NormalTok{, }\DataTypeTok{intent(in)}  \DataTypeTok{::}\NormalTok{ prog, input, depth}

    \DataTypeTok{logical}\NormalTok{, }\DataTypeTok{intent(out)} \DataTypeTok{::}\NormalTok{ settled, value}

    \DataTypeTok{logical} \DataTypeTok{::}\NormalTok{ s, x}

\NormalTok{    frames\_seen }\KeywordTok{=}\NormalTok{ frames\_seen }\KeywordTok{+} \DecValTok{1}

    \KeywordTok{if}\NormalTok{ (depth }\OperatorTok{\textgreater{}} \DecValTok{1} \OperatorTok{.and.}\NormalTok{ prog }\OperatorTok{==}\NormalTok{ D }\OperatorTok{.and.}\NormalTok{ input }\OperatorTok{==}\NormalTok{ D) frames\_identical }\KeywordTok{=}\NormalTok{ frames\_identical }\KeywordTok{+} \DecValTok{1}

    \KeywordTok{if}\NormalTok{ (depth }\OperatorTok{\textgreater{}=}\NormalTok{ FUEL) }\KeywordTok{then}

\NormalTok{       settled }\KeywordTok{=} \ConstantTok{.false.}\NormalTok{;  value }\KeywordTok{=} \ConstantTok{.false.}\NormalTok{;  }\KeywordTok{return}

    \KeywordTok{end if}

    \KeywordTok{call}\NormalTok{ run(input, input, depth }\KeywordTok{+} \DecValTok{1}\NormalTok{, s, x)}

\NormalTok{    settled }\KeywordTok{=}\NormalTok{ s;  value }\KeywordTok{=} \OperatorTok{.not.}\NormalTok{ x}

  \KeywordTok{end subroutine}\NormalTok{ run}

  \CommentTok{! Counterexample to: for every k, k*k mod 4 is 0 or 1.}

  \KeywordTok{pure} \DataTypeTok{logical} \KeywordTok{function}\NormalTok{ square\_counter(k)}

    \DataTypeTok{integer(i8)}\NormalTok{, }\DataTypeTok{intent(in)} \DataTypeTok{::}\NormalTok{ k}

    \DataTypeTok{integer(i8)} \DataTypeTok{::}\NormalTok{ r}

\NormalTok{    r }\KeywordTok{=} \BuiltInTok{mod}\NormalTok{(}\BuiltInTok{mod}\NormalTok{(k, }\DecValTok{4\_i8}\NormalTok{)}\KeywordTok{*}\BuiltInTok{mod}\NormalTok{(k, }\DecValTok{4\_i8}\NormalTok{), }\DecValTok{4\_i8}\NormalTok{)}

\NormalTok{    square\_counter }\KeywordTok{=} \OperatorTok{.not.}\NormalTok{ (r }\OperatorTok{==} \DecValTok{0\_i8} \OperatorTok{.or.}\NormalTok{ r }\OperatorTok{==} \DecValTok{1\_i8}\NormalTok{)}

  \KeywordTok{end function}\NormalTok{ square\_counter}

  \CommentTok{! Counterexample to: for every n, n*n + n + 41 is prime.}

  \KeywordTok{pure} \DataTypeTok{logical} \KeywordTok{function}\NormalTok{ euler\_counter(n)}

    \DataTypeTok{integer(i8)}\NormalTok{, }\DataTypeTok{intent(in)} \DataTypeTok{::}\NormalTok{ n}

    \DataTypeTok{integer(i8)} \DataTypeTok{::}\NormalTok{ m, q}

\NormalTok{    m }\KeywordTok{=}\NormalTok{ n}\KeywordTok{*}\NormalTok{n }\KeywordTok{+}\NormalTok{ n }\KeywordTok{+} \DecValTok{41\_i8}\NormalTok{;  euler\_counter }\KeywordTok{=} \ConstantTok{.false.}\NormalTok{;  q }\KeywordTok{=} \DecValTok{2\_i8}

    \KeywordTok{do} \KeywordTok{while}\NormalTok{ (q}\KeywordTok{*}\NormalTok{q }\OperatorTok{\textless{}=}\NormalTok{ m)}

       \KeywordTok{if}\NormalTok{ (}\BuiltInTok{mod}\NormalTok{(m, q) }\OperatorTok{==} \DecValTok{0\_i8}\NormalTok{) }\KeywordTok{then}

\NormalTok{          euler\_counter }\KeywordTok{=} \ConstantTok{.true.}\NormalTok{;  }\KeywordTok{return}

       \KeywordTok{end if}

\NormalTok{       q }\KeywordTok{=}\NormalTok{ q }\KeywordTok{+} \DecValTok{1\_i8}

    \KeywordTok{end do}

  \KeywordTok{end function}\NormalTok{ euler\_counter}

\KeywordTok{end module}\NormalTok{ gate}

\KeywordTok{program}\NormalTok{ gate\_thesis}

  \KeywordTok{use}\NormalTok{ gate}

  \KeywordTok{implicit} \KeywordTok{none}

  \DataTypeTok{logical} \DataTypeTok{::}\NormalTok{ settled, value, wall\_holds}

  \DataTypeTok{integer(i8)} \DataTypeTok{::}\NormalTok{ k, found}

  \FunctionTok{write(*}\NormalTok{,}\StringTok{\textquotesingle{}(A)\textquotesingle{}}\FunctionTok{)} \FunctionTok{repeat}\NormalTok{(}\StringTok{\textquotesingle{}=\textquotesingle{}}\NormalTok{,}\DecValTok{78}\NormalTok{)}

  \FunctionTok{write(*}\NormalTok{,}\StringTok{\textquotesingle{}(A)\textquotesingle{}}\FunctionTok{)} \StringTok{\textquotesingle{} THE NINTH GATE THESIS, EXECUTED\textquotesingle{}}

  \FunctionTok{write(*}\NormalTok{,}\StringTok{\textquotesingle{}(A)\textquotesingle{}}\FunctionTok{)} \FunctionTok{repeat}\NormalTok{(}\StringTok{\textquotesingle{}=\textquotesingle{}}\NormalTok{,}\DecValTok{78}\NormalTok{)}

\NormalTok{  wall\_holds }\KeywordTok{=} \OperatorTok{.not.}\NormalTok{ (settles(}\ConstantTok{.false.}\NormalTok{) }\OperatorTok{.or.}\NormalTok{ settles(}\ConstantTok{.true.}\NormalTok{))}

  \FunctionTok{write(*}\NormalTok{,}\StringTok{\textquotesingle{}(A,L1)\textquotesingle{}}\FunctionTok{)} \StringTok{\textquotesingle{} I   the wall : no settled value satisfies D(D) = not D(D)          \textquotesingle{}}\NormalTok{, wall\_holds}

  \KeywordTok{call}\NormalTok{ run(D, D, }\DecValTok{1}\NormalTok{, settled, value)}

  \FunctionTok{write(*}\NormalTok{,}\StringTok{\textquotesingle{}(A,I0,A,I0,A,L1)\textquotesingle{}}\FunctionTok{)} \StringTok{\textquotesingle{} II  the self : frames identical to their caller \textquotesingle{}}\NormalTok{, frames\_identical, }\StringTok{\textquotesingle{} of \textquotesingle{}}\NormalTok{, }\KeywordTok{\&}

\NormalTok{       frames\_seen }\KeywordTok{{-}} \DecValTok{1}\NormalTok{, }\StringTok{\textquotesingle{}, settled \textquotesingle{}}\NormalTok{, settled}

\NormalTok{  found }\KeywordTok{=} \KeywordTok{{-}}\DecValTok{1\_i8}

  \KeywordTok{do}\NormalTok{ k }\KeywordTok{=} \DecValTok{0\_i8}\NormalTok{, SEARCH\_FUEL }\KeywordTok{{-}} \DecValTok{1\_i8}

\BuiltInTok{     i}\NormalTok{f (square\_counter(k)) }\KeywordTok{then}

\NormalTok{        found }\KeywordTok{=}\NormalTok{ k;  }\KeywordTok{exit}

\BuiltInTok{     e}\NormalTok{nd }\KeywordTok{if}

  \KeywordTok{end do}

  \FunctionTok{write(*}\NormalTok{,}\StringTok{\textquotesingle{}(A,I0,A,L1)\textquotesingle{}}\FunctionTok{)} \StringTok{\textquotesingle{} III a true universal, k*k mod 4 in \{0,1\}, searched to \textquotesingle{}}\NormalTok{, SEARCH\_FUEL, }\StringTok{\textquotesingle{}: settled \textquotesingle{}}\NormalTok{, (found }\OperatorTok{\textgreater{}=} \DecValTok{0\_i8}\NormalTok{)}

\NormalTok{  found }\KeywordTok{=} \KeywordTok{{-}}\DecValTok{1\_i8}

  \KeywordTok{do}\NormalTok{ k }\KeywordTok{=} \DecValTok{0\_i8}\NormalTok{, SEARCH\_FUEL }\KeywordTok{{-}} \DecValTok{1\_i8}

\BuiltInTok{     i}\NormalTok{f (euler\_counter(k)) }\KeywordTok{then}

\NormalTok{        found }\KeywordTok{=}\NormalTok{ k;  }\KeywordTok{exit}

\BuiltInTok{     e}\NormalTok{nd }\KeywordTok{if}

  \KeywordTok{end do}

  \FunctionTok{write(*}\NormalTok{,}\StringTok{\textquotesingle{}(A,I0,A,I0)\textquotesingle{}}\FunctionTok{)} \StringTok{\textquotesingle{}     a false universal, n*n + n + 41 prime: settled false at n = \textquotesingle{}}\NormalTok{, found, }\KeywordTok{\&}

       \StringTok{\textquotesingle{}, value \textquotesingle{}}\NormalTok{, found}\KeywordTok{*}\NormalTok{found }\KeywordTok{+}\NormalTok{ found }\KeywordTok{+} \DecValTok{41\_i8}

  \FunctionTok{write(*}\NormalTok{,}\StringTok{\textquotesingle{}(A)\textquotesingle{}}\FunctionTok{)} \FunctionTok{repeat}\NormalTok{(}\StringTok{\textquotesingle{}=\textquotesingle{}}\NormalTok{,}\DecValTok{78}\NormalTok{)}

  \KeywordTok{if}\NormalTok{ (}\OperatorTok{.not.}\NormalTok{ wall\_holds }\OperatorTok{.or.}\NormalTok{ settled }\OperatorTok{.or.}\NormalTok{ frames\_identical }\OperatorTok{/=}\NormalTok{ frames\_seen}\CommentTok{ {-} 1 .or. found /= 40\_i8) \&}

\NormalTok{       error }\KeywordTok{stop} \StringTok{\textquotesingle{}the thesis failed\textquotesingle{}}

  \FunctionTok{write(*}\NormalTok{,}\StringTok{\textquotesingle{}(A)\textquotesingle{}}\FunctionTok{)} \StringTok{\textquotesingle{} Read: the wall. Executed: the self, unsettled. A run settles refutations only.\textquotesingle{}}

  \FunctionTok{write(*}\NormalTok{,}\StringTok{\textquotesingle{}(A)\textquotesingle{}}\FunctionTok{)} \FunctionTok{repeat}\NormalTok{(}\StringTok{\textquotesingle{}=\textquotesingle{}}\NormalTok{,}\DecValTok{78}\NormalTok{)}

\KeywordTok{end program}\NormalTok{ gate\_thesis}
\end{Highlighting}
\end{Shaded}
\end{strip}

Expected stdout of \texttt{gate\_thesis} (register run, gfortran 13.3.0;
transliteration-verified on the forging substrate):

\begin{strip}
\begin{Shaded}
\begin{Highlighting}[]
\OperatorTok{==============================================================================}
\NormalTok{ THE NINTH GATE THESIS}\OperatorTok{,}\NormalTok{ EXECUTED}
\OperatorTok{==============================================================================}
\NormalTok{ I   the wall }\OperatorTok{:}\NormalTok{ no settled value satisfies }\FunctionTok{D}\NormalTok{(D) }\OperatorTok{=}\NormalTok{ not }\FunctionTok{D}\NormalTok{(D)           T}
\NormalTok{ II  the self }\OperatorTok{:}\NormalTok{ frames identical to their caller }\DecValTok{99999} \KeywordTok{of} \DecValTok{99999}\OperatorTok{,}\NormalTok{ settled F}
\NormalTok{ III a }\KeywordTok{true}\NormalTok{ universal}\OperatorTok{,}\NormalTok{ k}\OperatorTok{*}\NormalTok{k mod }\DecValTok{4} \KeywordTok{in}\NormalTok{ \{}\DecValTok{0}\OperatorTok{,}\DecValTok{1}\NormalTok{\}}\OperatorTok{,}\NormalTok{ searched to }\DecValTok{2000000}\OperatorTok{:}\NormalTok{ settled F}
\NormalTok{     a }\KeywordTok{false}\NormalTok{ universal}\OperatorTok{,}\NormalTok{ n}\OperatorTok{*}\NormalTok{n }\OperatorTok{+}\NormalTok{ n }\OperatorTok{+} \DecValTok{41}\NormalTok{ prime}\OperatorTok{:}\NormalTok{ settled }\KeywordTok{false}\NormalTok{ at n }\OperatorTok{=} \DecValTok{40}\OperatorTok{,}\NormalTok{ value }\DecValTok{1681}
\OperatorTok{==============================================================================}
\NormalTok{ Read}\OperatorTok{:}\NormalTok{ the wall}\OperatorTok{.} \AttributeTok{Executed}\OperatorTok{:}\NormalTok{ the self}\OperatorTok{,}\NormalTok{ unsettled}\OperatorTok{.} \AttributeTok{A}\NormalTok{ run settles refutations only}\OperatorTok{.}
\OperatorTok{==============================================================================}
\end{Highlighting}
\end{Shaded}
\end{strip}

\subsection{B.5 · RAFGate.lean --- the Ninth Gate on RAF-native ground,
and BSD through
it}\label{b.5-rafgate.lean-the-ninth-gate-on-raf-native-ground-and-bsd-through-it}

\begin{strip}
\begin{Shaded}
\begin{Highlighting}[]
\NormalTok{/{-}}
\NormalTok{  RAF Gate · the Ninth Gate on RAF{-}native ground, and BSD through it · core Lean 4, no Mathlib.}
\NormalTok{  Pre{-}Gödel: no provability predicate, no diagonal lemma, no consistency hypothesis.}
\NormalTok{{-}/}

\NormalTok{theorem lock\_core (r t s n : Nat) (hK : s = r + t) (hC : s ≤ n) (hE : n ≤ r) : r = n ∧ t = 0 := by}
\NormalTok{  have h1 : r + t ≤ n := hK ▸ hC}
\NormalTok{  have h3 : r = n := Nat.le\_antisymm (Nat.le\_trans (Nat.le\_add\_right r t) h1) hE}
\NormalTok{  have h4 : r + t ≤ r := by rw [← h3] at h1; exact h1}
\NormalTok{  exact ⟨h3, Nat.eq\_zero\_of\_le\_zero (Nat.le\_of\_add\_le\_add\_left (Nat.le\_trans h4 (Nat.le\_of\_eq (Nat.add\_zero r).symm)))⟩}

\NormalTok{structure Curve where}
\NormalTok{  ralg : Nat            {-}{-} rank of E(Q)}
\NormalTok{  ran : Nat             {-}{-} ord\_\{s=1\} L(E,s)}
\NormalTok{  shaQ : Nat → Nat      {-}{-} corank of Sha(E/Q)[p\^{}∞]}
\NormalTok{  trivQ : Nat → Bool    {-}{-} Sha(E/Q)[p\^{}∞] = 0}
\NormalTok{  eps : Nat             {-}{-} ord\_\{s=1\} L(E\^{}K,s), at most one}
\NormalTok{  shaT : Nat → Nat      {-}{-} corank of Sha(E\^{}K/Q)[p\^{}∞]}
\NormalTok{  trivT : Nat → Bool}
\NormalTok{  rKalg : Nat           {-}{-} rank of E(K)}
\NormalTok{  rKan : Nat            {-}{-} ord\_\{s=1\} L(E/K,s)}
\NormalTok{  selK : Nat → Nat      {-}{-} corank of Sel\_\{p\^{}∞\}(E/K)}
\NormalTok{  shaK : Nat → Nat      {-}{-} corank of Sha(E/K)[p\^{}∞]}
\NormalTok{  trivK : Nat → Bool}

\NormalTok{/{-}{-} The theorem{-}grade inputs, each cited in the record. {-}/}
\NormalTok{structure Inputs (c : Curve) : Prop where}
\NormalTok{  factor    : c.rKan = c.ran + c.eps}
\NormalTok{  descent   : c.rKalg = c.ralg + c.eps}
\NormalTok{  kummerK   : ∀ p, c.selK p = c.rKalg + c.shaK p}
\NormalTok{  kernel    : ∀ p, c.shaQ p ≤ c.shaK p}
\NormalTok{  kerOdd    : ∀ p, 2 \textless{} p → c.trivK p = true → c.trivQ p = true}
\NormalTok{  twistCork : ∀ p, c.shaT p = 0}
\NormalTok{  twistTriv : ∃ N, ∀ p, N ≤ p → c.trivT p = true}
\NormalTok{  split     : ∀ p, c.shaK p ≤ c.shaQ p + c.shaT p}
\NormalTok{  splitOdd  : ∀ p, 2 \textless{} p → c.trivQ p = true → c.trivT p = true → c.trivK p = true}

\NormalTok{/{-}{-} The object: faces (i) and (ii) over Q. {-}/}
\NormalTok{def ObjectQ (c : Curve) : Prop :=}
\NormalTok{  c.ralg = c.ran ∧ (∀ p, c.shaQ p = 0) ∧ ∃ N, ∀ p, N ≤ p → c.trivQ p = true}

\NormalTok{/{-}{-} The one axiom at one curve, for the class z of the construction and its index bound idx:}
\NormalTok{    the class is nonzero (its height is a nonzero multiple of the nonzero derivative),}
\NormalTok{    every Selmer corank over K is bounded by the order, and Sha(E/K) vanishes beyond the index. {-}/}
\NormalTok{def Omega (c : Curve) (z idx : Nat) : Prop :=}
\NormalTok{  z ≠ 0 ∧ (∀ p, c.selK p ≤ c.rKan) ∧ (∀ p, idx ≤ p → c.trivK p = true)}

\NormalTok{/{-}{-} A genuine class lives in the top exterior power of the Mordell{-}Weil group over K. {-}/}
\NormalTok{def InLattice (c : Curve) (z : Nat) : Prop := z ≠ 0 → c.rKan ≤ c.rKalg}

\NormalTok{theorem bound3 (a b p : Nat) (hp : a + b + 3 ≤ p) : 2 \textless{} p ∧ a ≤ p ∧ b ≤ p :=}
\NormalTok{  ⟨Nat.lt\_of\_lt\_of\_le (by decide : 2 \textless{} 3) (Nat.le\_trans (Nat.le\_add\_left 3 (a + b)) hp),}
\NormalTok{   Nat.le\_trans (Nat.le\_trans (Nat.le\_add\_right a b) (Nat.le\_add\_right (a + b) 3)) hp,}
\NormalTok{   Nat.le\_trans (Nat.le\_trans (Nat.le\_add\_left b a) (Nat.le\_add\_right (a + b) 3)) hp⟩}

\NormalTok{/{-}! \#\# The full proof: from the one axiom to the object {-}/}

\NormalTok{theorem omega\_proves\_object (c : Curve) (hin : Inputs c) (z idx : Nat)}
\NormalTok{    (hL : InLattice c z) (hΩ : Omega c z idx) : ObjectQ c := by}
\NormalTok{  obtain ⟨hz, hsel, hidx⟩ := hΩ}
\NormalTok{  have hEX : c.rKan ≤ c.rKalg := hL hz}
\NormalTok{  have hrk : c.rKalg = c.rKan := (lock\_core \_ \_ \_ \_ (hin.kummerK 0) (hsel 0) hEX).1}
\NormalTok{  have hr : c.ralg = c.ran := by}
\NormalTok{    have h := hrk}
\NormalTok{    rw [hin.descent, hin.factor] at h}
\NormalTok{    exact Nat.add\_right\_cancel h}
\NormalTok{  refine ⟨hr, fun p =\textgreater{} ?\_, ⟨idx + 0 + 3, fun p hp =\textgreater{} ?\_⟩⟩}
\NormalTok{  · have hK0 := (lock\_core \_ \_ \_ \_ (hin.kummerK p) (hsel p) hEX).2}
\NormalTok{    have hk := hin.kernel p}
\NormalTok{    rw [hK0] at hk}
\NormalTok{    exact Nat.eq\_zero\_of\_le\_zero hk}
\NormalTok{  · have hb := bound3 idx 0 p hp}
\NormalTok{    exact hin.kerOdd p hb.1 (hidx p hb.2.1)}

\NormalTok{/{-}! \#\# The formal termination of the axiom {-}/}

\NormalTok{/{-}{-} Hidden circularity: read without naming its construction, the axiom is the object. {-}/}
\NormalTok{theorem hidden\_circularity (c : Curve) (hin : Inputs c) :}
\NormalTok{    ObjectQ c ↔ ∃ z idx, InLattice c z ∧ Omega c z idx := by}
\NormalTok{  constructor}
\NormalTok{  · intro ⟨hr, hsha, N, hN⟩}
\NormalTok{    obtain ⟨M, hM⟩ := hin.twistTriv}
\NormalTok{    have hrk : c.rKalg = c.rKan := by rw [hin.descent, hin.factor, hr]}
\NormalTok{    refine ⟨1, N + M + 3, fun \_ =\textgreater{} Nat.le\_of\_eq hrk.symm, (by decide : (1 : Nat) ≠ 0), fun p =\textgreater{} ?\_, fun p hp =\textgreater{} ?\_⟩}
\NormalTok{    · have hs : c.shaK p ≤ 0 := by}
\NormalTok{        have h := hin.split p}
\NormalTok{        rw [hsha p, hin.twistCork p] at h}
\NormalTok{        exact h}
\NormalTok{      rw [hin.kummerK p, Nat.eq\_zero\_of\_le\_zero hs, Nat.add\_zero, hrk]}
\NormalTok{      exact Nat.le\_refl \_}
\NormalTok{    · have hb := bound3 N M p hp}
\NormalTok{      exact hin.splitOdd p hb.1 (hN p hb.2.1) (hM p hb.2.2)}
\NormalTok{  · intro ⟨z, idx, hL, hΩ⟩}
\NormalTok{    exact omega\_proves\_object c hin z idx hL hΩ}

\NormalTok{/{-}{-} A rank{-}two model: the object holds, and the first{-}derivative construction fails the axiom. {-}/}
\NormalTok{def c2 : Curve :=}
\NormalTok{  \{ ralg := 2, ran := 2, shaQ := fun \_ =\textgreater{} 0, trivQ := fun \_ =\textgreater{} true, eps := 1,}
\NormalTok{    shaT := fun \_ =\textgreater{} 0, trivT := fun \_ =\textgreater{} true, rKalg := 3, rKan := 3,}
\NormalTok{    selK := fun \_ =\textgreater{} 3, shaK := fun \_ =\textgreater{} 0, trivK := fun \_ =\textgreater{} true \}}

\NormalTok{theorem inputs\_c2 : Inputs c2 :=}
\NormalTok{  \{ factor := rfl, descent := rfl, kummerK := fun \_ =\textgreater{} rfl, kernel := fun \_ =\textgreater{} Nat.le\_refl 0,}
\NormalTok{    kerOdd := fun \_ \_ h =\textgreater{} h, twistCork := fun \_ =\textgreater{} rfl, twistTriv := ⟨0, fun \_ \_ =\textgreater{} rfl⟩,}
\NormalTok{    split := fun \_ =\textgreater{} Nat.le\_refl 0, splitOdd := fun \_ \_ h \_ =\textgreater{} h \}}

\NormalTok{/{-}{-} A model satisfying every input where the object fails: rank zero against analytic order two. {-}/}
\NormalTok{def cFail : Curve :=}
\NormalTok{  \{ ralg := 0, ran := 2, shaQ := fun \_ =\textgreater{} 0, trivQ := fun \_ =\textgreater{} true, eps := 1,}
\NormalTok{    shaT := fun \_ =\textgreater{} 0, trivT := fun \_ =\textgreater{} true, rKalg := 1, rKan := 3,}
\NormalTok{    selK := fun \_ =\textgreater{} 1, shaK := fun \_ =\textgreater{} 0, trivK := fun \_ =\textgreater{} true \}}

\NormalTok{theorem inputs\_cFail : Inputs cFail :=}
\NormalTok{  \{ factor := rfl, descent := rfl, kummerK := fun \_ =\textgreater{} rfl, kernel := fun \_ =\textgreater{} Nat.le\_refl 0,}
\NormalTok{    kerOdd := fun \_ \_ h =\textgreater{} h, twistCork := fun \_ =\textgreater{} rfl, twistTriv := ⟨0, fun \_ \_ =\textgreater{} rfl⟩,}
\NormalTok{    split := fun \_ =\textgreater{} Nat.le\_refl 0, splitOdd := fun \_ \_ h \_ =\textgreater{} h \}}

\NormalTok{/{-}! \#\# The Ninth Gate on RAF{-}native ground: block theorem, Lawvere, no exterior agent {-}/}

\NormalTok{/{-}{-} T1, freeness: an answer{-}flipping involution fixes no point. {-}/}
\NormalTok{theorem defeater\_free \{X : Type\} (τ : X → X) (d : X → Bool) (hflip : ∀ x, d (τ x) = !d x) : ∀ x, τ x ≠ x := by}
\NormalTok{  intro x h}
\NormalTok{  have hx := hflip x}
\NormalTok{  rw [h] at hx}
\NormalTok{  cases e : d x \textless{};\textgreater{} rw [e] at hx \textless{};\textgreater{} exact absurd hx (by decide)}

\NormalTok{/{-}{-} The Block theorem: no readout even under an answer{-}flipping involution decides the answer. {-}/}
\NormalTok{theorem block \{X R : Type\} (τ : X → X) (d : X → Bool) (hflip : ∀ x, d (τ x) = !d x)}
\NormalTok{    (r : X → R) (heven : ∀ x, r (τ x) = r x) (x0 : X) : ¬ ∃ h : R → Bool, ∀ x, d x = h (r x) := by}
\NormalTok{  intro ⟨h, hd⟩}
\NormalTok{  have h1 : d (τ x0) = d x0 := by rw [hd (τ x0), hd x0, heven x0]}
\NormalTok{  rw [hflip x0] at h1}
\NormalTok{  cases hx : d x0 \textless{};\textgreater{} rw [hx] at h1 \textless{};\textgreater{} exact absurd h1 (by decide)}

\NormalTok{/{-}{-} Lawvere: a point{-}surjective self{-}indexing forces every endomap to have a fixed point. {-}/}
\NormalTok{theorem lawvere \{A B : Type\} (e : A → A → B) (hs : ∀ g : A → B, ∃ a, e a = g) (f : B → B) : ∃ b, f b = b := by}
\NormalTok{  obtain ⟨a, ha⟩ := hs (fun x =\textgreater{} f (e x x))}
\NormalTok{  exact ⟨e a a, (congrFun ha a).symm⟩}

\NormalTok{/{-}{-} RAF{-}C2: no register indexes all of its own properties. {-}/}
\NormalTok{theorem no\_total\_self\_index \{A : Type\} (e : A → A → Bool) : ¬ ∀ g : A → Bool, ∃ a, e a = g := by}
\NormalTok{  intro hs}
\NormalTok{  obtain ⟨b, hb⟩ := lawvere e hs (fun b =\textgreater{} !b)}
\NormalTok{  cases b \textless{};\textgreater{} exact absurd hb (by decide)}

\NormalTok{/{-}{-} RAF{-}C1: with registration equivalent to interaction and closure under registration, nothing outside acts. {-}/}
\NormalTok{theorem no\_exterior\_agent \{Sys : Type\} (inD interacts registers : Sys → Prop)}
\NormalTok{    (raf1 : ∀ s, interacts s ↔ registers s) (raf2 : ∀ s, registers s → inD s) :}
\NormalTok{    ∀ s, ¬ inD s → ¬ interacts s :=}
\NormalTok{  fun s hs hi =\textgreater{} hs (raf2 s ((raf1 s).1 hi))}

\NormalTok{/{-}{-} Even supply never crosses: a supplier whose record is also even leaves the block standing. {-}/}
\NormalTok{theorem even\_supply\_does\_not\_cross \{X R S : Type\} (τ : X → X) (d : X → Bool) (hflip : ∀ x, d (τ x) = !d x)}
\NormalTok{    (r : X → R) (s : X → S) (hr : ∀ x, r (τ x) = r x) (hs : ∀ x, s (τ x) = s x) (x0 : X) :}
\NormalTok{    ¬ ∃ h : R × S → Bool, ∀ x, d x = h (r x, s x) :=}
\NormalTok{  block τ d hflip (fun x =\textgreater{} (r x, s x)) (fun x =\textgreater{} by rw [hr x, hs x]) x0}

\NormalTok{/{-}{-} Odd supply crosses: one supplied bit that flips with the involution, beside an orbit{-}separating even readout,}
\NormalTok{    determines the answer. {-}/}
\NormalTok{theorem odd\_supply\_crosses \{X R : Type\} (τ : X → X) (d : X → Bool) (\_hflip : ∀ x, d (τ x) = !d x)}
\NormalTok{    (r : X → R) (hsep : ∀ x y, r x = r y → y = x ∨ y = τ x)}
\NormalTok{    (o : X → Bool) (hodd : ∀ x, o (τ x) = !o x) : ∀ x y, r x = r y → o x = o y → d x = d y := by}
\NormalTok{  intro x y hr ho}
\NormalTok{  cases hsep x y hr with}
\NormalTok{  | inl h =\textgreater{} rw [h]}
\NormalTok{  | inr h =\textgreater{}}
\NormalTok{    rw [h, hodd x] at ho}
\NormalTok{    cases hx : o x \textless{};\textgreater{} rw [hx] at ho \textless{};\textgreater{} exact absurd ho (by decide)}

\NormalTok{/{-}! \#\# BSD through the gate {-}/}

\NormalTok{/{-}{-} The model swap: one model of the inputs where the object holds, one where it fails. {-}/}
\NormalTok{def model (b : Bool) : Curve := if b then c2 else cFail}

\NormalTok{theorem model\_inputs (b : Bool) : Inputs (model b) := by}
\NormalTok{  cases b}
\NormalTok{  · exact inputs\_cFail}
\NormalTok{  · exact inputs\_c2}

\NormalTok{theorem object\_flips : ObjectQ (model true) ∧ ¬ ObjectQ (model false) :=}
\NormalTok{  ⟨⟨rfl, fun \_ =\textgreater{} rfl, ⟨0, fun \_ \_ =\textgreater{} rfl⟩⟩, fun h =\textgreater{} absurd h.1 (by decide)⟩}

\NormalTok{/{-}{-} Every reading that holds in all models of the inputs is even under the swap. {-}/}
\NormalTok{theorem seedless\_readings\_are\_even (W : Curve → Prop) (hW : ∀ c, Inputs c → W c) :}
\NormalTok{    W (model true) ∧ W (model false) :=}
\NormalTok{  ⟨hW \_ (model\_inputs true), hW \_ (model\_inputs false)⟩}

\NormalTok{/{-}{-} So no seedless reading, whether a termination theorem, a gate theorem or any other, decides the object. {-}/}
\NormalTok{theorem seedless\_cannot\_decide (W : Curve → Prop) (hW : ∀ c, Inputs c → W c) :}
\NormalTok{    ¬ ∀ b, (ObjectQ (model b) ↔ W (model b)) :=}
\NormalTok{  fun h =\textgreater{} object\_flips.2 ((h false).2 (hW \_ (model\_inputs false)))}

\NormalTok{/{-}{-} The one axiom is an odd supply: it holds on the side where the object holds and fails on the other. {-}/}
\NormalTok{theorem omega\_is\_odd\_supply :}
\NormalTok{    (∃ z idx, InLattice (model true) z ∧ Omega (model true) z idx) ∧}
\NormalTok{    ¬ ∃ z idx, InLattice (model false) z ∧ Omega (model false) z idx :=}
\NormalTok{  ⟨(hidden\_circularity \_ (model\_inputs true)).1 object\_flips.1,}
\NormalTok{   fun h =\textgreater{} object\_flips.2 ((hidden\_circularity \_ (model\_inputs false)).2 h)⟩}

\NormalTok{/{-}{-} A printed record whose content reads no curve data holds on both sides of the swap, so the record cannot}
\NormalTok{    cross. This is the reading side only: the deed that printed it is not an object of this register. {-}/}
\NormalTok{theorem fixed\_receipt\_cannot\_cross (rec : Prop) (h : rec) : ¬ ∀ b, (ObjectQ (model b) ↔ rec) :=}
\NormalTok{  seedless\_cannot\_decide (fun \_ =\textgreater{} rec) (fun \_ \_ =\textgreater{} h)}

\NormalTok{\#print axioms block}
\NormalTok{\#print axioms lawvere}
\NormalTok{\#print axioms no\_total\_self\_index}
\NormalTok{\#print axioms no\_exterior\_agent}
\NormalTok{\#print axioms even\_supply\_does\_not\_cross}
\NormalTok{\#print axioms odd\_supply\_crosses}
\NormalTok{\#print axioms seedless\_cannot\_decide}
\NormalTok{\#print axioms omega\_is\_odd\_supply}
\NormalTok{\#print axioms fixed\_receipt\_cannot\_cross}
\end{Highlighting}
\end{Shaded}
\end{strip}

\subsection{B.6 · bsd\_rows.f90 --- BSD on its own rows, executed by the
Core Thesis's compiled
procedures}\label{b.6-bsd_rows.f90-bsd-on-its-own-rows-executed-by-the-core-thesiss-compiled-procedures}

Compiles against \texttt{module\ trisduction\_rows} of {[}13{]}, the
Core Thesis in Fortran, itself sealed with its own battery (248 checks,
0 failures). Every flag below is a hand-built reading, supplied and
named; the thesis routes it and authors nothing.

\begin{strip}
\begin{Shaded}
\begin{Highlighting}[]
\CommentTok{!===============================================================================}
\CommentTok{!  BSD ON ITS OWN ROWS, EXECUTED BY THE CORE THESIS\textquotesingle{}S COMPILED PROCEDURES.}
\CommentTok{!  Every flag below is a hand{-}built reading, supplied and named; the thesis}
\CommentTok{!  routes it and authors nothing. Anchors: EX exhibition, CL closure, UN}
\CommentTok{!  normalization. Faces scoped per R{-}1: F1 rank identity, F2 finiteness of Sha,}
\CommentTok{!  F3 leading term.}
\CommentTok{!===============================================================================}
\KeywordTok{program}\NormalTok{ bsd\_rows}
  \KeywordTok{use}\NormalTok{ trisduction\_rows}
  \KeywordTok{implicit} \KeywordTok{none}
  \DataTypeTok{integer}  \DataTypeTok{::}\NormalTok{ v, c, d, i}
  \DataTypeTok{real(dp)} \DataTypeTok{::}\NormalTok{ sep, Nrm(}\DecValTok{3}\NormalTok{,}\DecValTok{3}\NormalTok{), G(}\DecValTok{3}\NormalTok{,}\DecValTok{3}\NormalTok{)}
  \DataTypeTok{character(len=110)} \DataTypeTok{::}\NormalTok{ why}
  \DataTypeTok{integer}  \DataTypeTok{::}\NormalTok{ sl13(}\DecValTok{3}\NormalTok{,}\DecValTok{4}\NormalTok{), slF3(}\DecValTok{3}\NormalTok{,}\DecValTok{4}\NormalTok{)}
  \DataTypeTok{logical}  \DataTypeTok{::}\NormalTok{ pop(}\DecValTok{3}\NormalTok{)}

  \FunctionTok{write(*}\NormalTok{,}\StringTok{\textquotesingle{}(A)\textquotesingle{}}\FunctionTok{)} \FunctionTok{repeat}\NormalTok{(}\StringTok{\textquotesingle{}=\textquotesingle{}}\NormalTok{,}\DecValTok{78}\NormalTok{)}
  \FunctionTok{write(*}\NormalTok{,}\StringTok{\textquotesingle{}(A)\textquotesingle{}}\FunctionTok{)} \StringTok{\textquotesingle{} BSD ON ITS OWN ROWS, EXECUTED\textquotesingle{}}
  \FunctionTok{write(*}\NormalTok{,}\StringTok{\textquotesingle{}(A)\textquotesingle{}}\FunctionTok{)} \FunctionTok{repeat}\NormalTok{(}\StringTok{\textquotesingle{}=\textquotesingle{}}\NormalTok{,}\DecValTok{78}\NormalTok{)}

  \CommentTok{! Seal L. Vocabulary ids: EX 101 point 102 construction 103 height pairing 104 Heegner;}
  \CommentTok{! CL 201 Selmer 202 descent 203 Euler system 204 Kummer; UN 301 L{-}value 302 order of vanishing}
  \CommentTok{! 303 period 304 leading coefficient. For F3 the leading term carries the regulator, id 103.}
\NormalTok{  sl13 }\KeywordTok{=} \FunctionTok{reshape}\NormalTok{(}\KeywordTok{[}\DecValTok{101}\NormalTok{,}\DecValTok{102}\NormalTok{,}\DecValTok{103}\NormalTok{,}\DecValTok{104}\NormalTok{, }\DecValTok{201}\NormalTok{,}\DecValTok{202}\NormalTok{,}\DecValTok{203}\NormalTok{,}\DecValTok{204}\NormalTok{, }\DecValTok{301}\NormalTok{,}\DecValTok{302}\NormalTok{,}\DecValTok{303}\NormalTok{,}\DecValTok{304}\KeywordTok{]}\NormalTok{,}\KeywordTok{[}\DecValTok{3}\NormalTok{,}\CommentTok{4],order=[2,1])}
  \KeywordTok{call}\NormalTok{ seal\_L(sl13, v, why)}
  \FunctionTok{write(*}\NormalTok{,}\StringTok{\textquotesingle{}(A,A,A,A)\textquotesingle{}}\FunctionTok{)} \StringTok{\textquotesingle{} Seal L, F1 and F2 slots   : \textquotesingle{}}\NormalTok{, }\FunctionTok{trim}\NormalTok{(token(v)), }\CommentTok{\textquotesingle{} :: \textquotesingle{}, trim(why)}
\NormalTok{  slF3 }\KeywordTok{=} \FunctionTok{reshape}\NormalTok{(}\KeywordTok{[}\DecValTok{101}\NormalTok{,}\DecValTok{102}\NormalTok{,}\DecValTok{103}\NormalTok{,}\DecValTok{104}\NormalTok{, }\DecValTok{201}\NormalTok{,}\DecValTok{202}\NormalTok{,}\DecValTok{203}\NormalTok{,}\DecValTok{204}\NormalTok{, }\DecValTok{301}\NormalTok{,}\DecValTok{302}\NormalTok{,}\DecValTok{303}\NormalTok{,}\DecValTok{103}\KeywordTok{]}\NormalTok{,}\KeywordTok{[}\DecValTok{3}\NormalTok{,}\CommentTok{4],order=[2,1])}
  \KeywordTok{call}\NormalTok{ seal\_L(slF3, v, why)}
  \FunctionTok{write(*}\NormalTok{,}\StringTok{\textquotesingle{}(A,A,A,A)\textquotesingle{}}\FunctionTok{)} \StringTok{\textquotesingle{} Seal L, F3 slots          : \textquotesingle{}}\NormalTok{, }\FunctionTok{trim}\NormalTok{(token(v)), }\CommentTok{\textquotesingle{} :: \textquotesingle{}, trim(why)}

  \FunctionTok{write(*}\NormalTok{,}\StringTok{\textquotesingle{}(/,A)\textquotesingle{}}\FunctionTok{)} \StringTok{\textquotesingle{} THE ROW CASCADE, PER FACE AND STRATUM\textquotesingle{}}
  \CommentTok{! F1, analytic rank at most one, every such curve: rows populated by Heegner point, Kolyvagin, Gross{-}Zagier;}
  \CommentTok{! the auxiliary field K is an existential import and modularity a transport, so frame\_closed is false.}
  \KeywordTok{call}\NormalTok{ row\_cascade(}\ConstantTok{.true.}\NormalTok{,}\ConstantTok{.true.}\NormalTok{,}\ConstantTok{.true.}\NormalTok{,}\ConstantTok{.false.}\NormalTok{, }\ConstantTok{.true.}\NormalTok{,}\ConstantTok{.false.}\NormalTok{,}\ConstantTok{.true.}\NormalTok{,}\ConstantTok{.false.}\NormalTok{, c, why)}
  \KeywordTok{call}\NormalTok{ show(}\StringTok{\textquotesingle{}F1  order \textless{}= 1, all curves\textquotesingle{}}\NormalTok{, c, why)}
  \CommentTok{! F2 at the same stratum: Kolyvagin bounds Sha; same frame leak.}
  \KeywordTok{call}\NormalTok{ row\_cascade(}\ConstantTok{.true.}\NormalTok{,}\ConstantTok{.true.}\NormalTok{,}\ConstantTok{.true.}\NormalTok{,}\ConstantTok{.false.}\NormalTok{, }\ConstantTok{.true.}\NormalTok{,}\ConstantTok{.false.}\NormalTok{,}\ConstantTok{.true.}\NormalTok{,}\ConstantTok{.false.}\NormalTok{, c, why)}
  \KeywordTok{call}\NormalTok{ show(}\StringTok{\textquotesingle{}F2  order \textless{}= 1, all curves\textquotesingle{}}\NormalTok{, c, why)}
  \CommentTok{! F1 on 389a1: EX populated by the executed height deed, CL and UN by cited computation; a row furnished by the deed.}
  \KeywordTok{call}\NormalTok{ row\_cascade(}\ConstantTok{.true.}\NormalTok{,}\ConstantTok{.true.}\NormalTok{,}\ConstantTok{.true.}\NormalTok{,}\ConstantTok{.false.}\NormalTok{, }\ConstantTok{.true.}\NormalTok{,}\ConstantTok{.true.}\NormalTok{,}\ConstantTok{.true.}\NormalTok{,}\ConstantTok{.true.}\NormalTok{, c, why)}
  \KeywordTok{call}\NormalTok{ show(}\StringTok{\textquotesingle{}F1  389a1, executed\textquotesingle{}}\NormalTok{, c, why)}
  \CommentTok{! F2 on 389a1: no closure row bounds Sha at rank two.}
  \KeywordTok{call}\NormalTok{ row\_cascade(}\ConstantTok{.true.}\NormalTok{,}\ConstantTok{.true.}\NormalTok{,}\ConstantTok{.true.}\NormalTok{,}\ConstantTok{.false.}\NormalTok{, }\ConstantTok{.false.}\NormalTok{,}\ConstantTok{.true.}\NormalTok{,}\ConstantTok{.true.}\NormalTok{,}\ConstantTok{.true.}\NormalTok{, c, why)}
  \KeywordTok{call}\NormalTok{ show(}\StringTok{\textquotesingle{}F2  389a1\textquotesingle{}}\NormalTok{, c, why)}
  \CommentTok{! F1, analytic rank two and above, uniformly: no construction exhibits the points from L{-}data.}
  \KeywordTok{call}\NormalTok{ row\_cascade(}\ConstantTok{.true.}\NormalTok{,}\ConstantTok{.true.}\NormalTok{,}\ConstantTok{.true.}\NormalTok{,}\ConstantTok{.false.}\NormalTok{, }\ConstantTok{.false.}\NormalTok{,}\ConstantTok{.false.}\NormalTok{,}\ConstantTok{.true.}\NormalTok{,}\CommentTok{.true., c, why)}
  \KeywordTok{call}\NormalTok{ show(}\StringTok{\textquotesingle{}F1  order \textgreater{}= 2, uniform\textquotesingle{}}\NormalTok{, c, why)}
  \CommentTok{! The universal conjecture as one string, faces unscoped.}
  \KeywordTok{call}\NormalTok{ row\_cascade(}\ConstantTok{.true.}\NormalTok{,}\ConstantTok{.false.}\NormalTok{,}\ConstantTok{.true.}\NormalTok{,}\ConstantTok{.false.}\NormalTok{, }\ConstantTok{.true.}\NormalTok{,}\ConstantTok{.false.}\NormalTok{,}\ConstantTok{.true.}\NormalTok{,}\CommentTok{.true., c, why)}
  \KeywordTok{call}\NormalTok{ show(}\StringTok{\textquotesingle{}BSD as one unscoped string\textquotesingle{}}\NormalTok{, c, why)}

  \FunctionTok{write(*}\NormalTok{,}\StringTok{\textquotesingle{}(/,A)\textquotesingle{}}\FunctionTok{)} \StringTok{\textquotesingle{} THE SPINE ON THE POPULATION PATTERN\textquotesingle{}}
  \KeywordTok{do}\NormalTok{ i }\KeywordTok{=} \DecValTok{1}\NormalTok{, }\DecValTok{2}
\BuiltInTok{     N}\NormalTok{rm }\KeywordTok{=} \FloatTok{0.0\_dp}
\BuiltInTok{     N}\NormalTok{rm(}\DecValTok{1}\NormalTok{,}\DecValTok{1}\NormalTok{) }\KeywordTok{=} \FloatTok{1.0\_dp}\NormalTok{; Nrm(}\DecValTok{2}\NormalTok{,}\DecValTok{2}\NormalTok{) }\KeywordTok{=} \FloatTok{1.0\_dp}\NormalTok{; Nrm(}\DecValTok{3}\NormalTok{,}\DecValTok{3}\NormalTok{) }\KeywordTok{=} \FloatTok{1.0\_dp}
\BuiltInTok{     i}\NormalTok{f (i }\OperatorTok{==} \DecValTok{2}\NormalTok{) Nrm(}\DecValTok{1}\NormalTok{,:) }\KeywordTok{=} \FloatTok{0.0\_dp}                    \CommentTok{! EX axis empty}
\BuiltInTok{     G} \KeywordTok{=} \FunctionTok{matmul}\NormalTok{(Nrm, }\FunctionTok{transpose}\NormalTok{(Nrm))}
\BuiltInTok{     c}\NormalTok{all intersection\_dim(G, d, sep)}
\BuiltInTok{     i}\NormalTok{f (i }\OperatorTok{==} \DecValTok{1}\NormalTok{) }\KeywordTok{then}
        \FunctionTok{write(*}\NormalTok{,}\StringTok{\textquotesingle{}(A,I0,A)\textquotesingle{}}\FunctionTok{)} \StringTok{\textquotesingle{}   EX, CL, UN populated    : intersection dimension \textquotesingle{}}\NormalTok{, d, }\StringTok{\textquotesingle{}, one point, the lock\textquotesingle{}}
\BuiltInTok{     e}\NormalTok{lse}
        \FunctionTok{write(*}\NormalTok{,}\StringTok{\textquotesingle{}(A,I0,A)\textquotesingle{}}\FunctionTok{)} \StringTok{\textquotesingle{}   EX empty, CL, UN        : intersection dimension \textquotesingle{}}\NormalTok{, d, }\StringTok{\textquotesingle{}, a line, no point determined\textquotesingle{}}
\BuiltInTok{     e}\NormalTok{nd }\KeywordTok{if}
  \KeywordTok{end do}
  \FunctionTok{write(*}\NormalTok{,}\StringTok{\textquotesingle{}(A)\textquotesingle{}}\FunctionTok{)} \FunctionTok{repeat}\NormalTok{(}\StringTok{\textquotesingle{}=\textquotesingle{}}\NormalTok{,}\DecValTok{78}\NormalTok{)}

\CommentTok{contains}
  \KeywordTok{subroutine}\NormalTok{ show(label, comp, reason)}
    \DataTypeTok{character(len=*)}\NormalTok{, }\DataTypeTok{intent(in)} \DataTypeTok{::}\NormalTok{ label, reason}
    \DataTypeTok{integer}\NormalTok{, }\DataTypeTok{intent(in)} \DataTypeTok{::}\NormalTok{ comp}
    \FunctionTok{write(*}\NormalTok{,}\StringTok{\textquotesingle{}(A,A,A,A)\textquotesingle{}}\FunctionTok{)} \StringTok{\textquotesingle{}   \textquotesingle{}}\NormalTok{, label, }\FunctionTok{repeat}\NormalTok{(}\StringTok{\textquotesingle{} \textquotesingle{}}\NormalTok{, }\BuiltInTok{max}\NormalTok{(}\DecValTok{1}\NormalTok{, }\DecValTok{28}\KeywordTok{{-}}\FunctionTok{len}\NormalTok{(label))}\CommentTok{), trim(compartment\_name(comp))}
    \FunctionTok{write(*}\NormalTok{,}\StringTok{\textquotesingle{}(A,A)\textquotesingle{}}\FunctionTok{)} \StringTok{\textquotesingle{}       \textquotesingle{}}\NormalTok{, }\FunctionTok{trim}\NormalTok{(reason)}
    \FunctionTok{write(*}\NormalTok{,}\StringTok{\textquotesingle{}(A,A)\textquotesingle{}}\FunctionTok{)} \StringTok{\textquotesingle{}       \textquotesingle{}}\NormalTok{, }\FunctionTok{trim}\NormalTok{(compartment\_grade(comp, }\StringTok{\textquotesingle{}theorem\textquotesingle{}}\NormalTok{))}
  \KeywordTok{end subroutine}\NormalTok{ show}
\KeywordTok{end program}\NormalTok{ bsd\_rows}
\end{Highlighting}
\end{Shaded}
\end{strip}

\section{Appendix C · The Perelman Crossing: The Grand
Witness}\label{appendix-c-the-perelman-crossing-the-grand-witness}

\emph{Locator note.} The architect's note cites the companions' internal
numbering, and every locator has been verified against the finalized
text of {[}1, 2{]}; each is bound here to its restatement in this paper,
and no locator is load-bearing beyond its anchor. The block law is
Theorem 4.1 of {[}1{]} (Odd--Supply Separation), restated and proved
here as Theorem 2.3. The supply-entry typing is Definition 3.3 of
{[}1{]} (reading, proof availability, verified entry), restated here as
Definition 4.1. The deletion test is Theorem 4.5 of {[}1{]} (Provenance
Collapse), restated in Section 2. The dated-inventory law is Theorem 5.1
of {[}1{]} (First-Sufficient-Entry), restated as the Dated Inventory of
Section 2 with Corollary 4.3. The seat law is Theorem 5.3 of {[}1{]}
(Seat-Hiding Necessity): no terminal family without the seat input
distinguishes a crossed occupant from an uncrossed one. The evenness of
the pre-November-2002 Ricci-flow reading class is Corollary 5.2 of
{[}1{]}: parabolic rescaling acts with exponent zero on the named
generators. The self-similarity pattern is Table 3 of {[}1{]} (the SIMS
diagnostic: the κ-solutions classified, W the one supplied instance).
The truth-direction passage is §13.4 of {[}2{]}, quoted verbatim in C.3.

The Perelman crossing is self-referential at exactly three layers, and
the disguise is the paper's own costume.

\subsection{C.1 · Layer 1 --- the object: the manifold convicts
itself}\label{c.1-layer-1-the-object-the-manifold-convicts-itself}

Ricci flow is self-driven, ∂g/∂t = −2Ric(g): the metric's velocity is
its own curvature {[}20, 24{]}. The W-functional infimizes over
probability densities living on the very manifold under test --- the
manifold measures itself. Its stationarity set is exactly the gradient
shrinking solitons, the self-reproducing fixed points of the rescaled
flow. The companion's self-similarity diagnostic sees this shape without
naming it {[}1, Table 3{]}: the κ-solutions, the self-similar part, are
the settled adjacent fact, and W is the one supplied instance in the
entire seven-row pattern {[}1, 2{]}. Self-referential. Grade structural
--- a reading of the cited mechanism; no new object-mathematics is
claimed.

\subsection{C.2 · Layer 2 --- the disguise: a supply wearing a reading's
costume}\label{c.2-layer-2-the-disguise-a-supply-wearing-a-readings-costume}

The pre-November-2002 Ricci-flow reading class is even under the row's
native flip, parabolic rescaling --- exponent zero on the named
generators {[}1, Corollary 5.2{]}. The target, noncollapsing, is odd
under that flip. By Theorem 2.3, no coalition of even readings decides
it --- so W could not have been found by reading; it had to enter as a
construction, and Definition 4.1 types it exactly so: dependency hash
absent from the dated corpus at entry, November 2002, arXiv:math/0211159
{[}20{]}. Run the deletion test of Section 2: deleting the conclusion's
name from W's dependency chain does not collapse it to anything in the
pre-2002 corpus; the entry is a supply, not a restatement. The audit
discipline's costume detectors were built to catch restatements dressed
as supplies; the historical record shows the inverse costume --- a
genuine supply dressed as an ordinary reading, its oddness invisible to
the even register by its own parity law. The object supplies its own
missing coordinate, and the supply's decisive property is its
provenance, not its formula. Grade structural --- the deletion test's
own remark admits its named instances are ``named test candidates, not
executed validations'' {[}1, Theorem 4.5 remark{]}, and the architect's
assessment, stated at the same grade, is that the same workload is owed
for the Poincaré row.

\subsection{C.3 · Layer 3 --- the token: the
money}\label{c.3-layer-3-the-token-the-money}

The Dated Inventory says the corpus token changes only at first verified
entry; Definition 4.1 separates denotability, admission, and
verification. Perelman's deed entered by posting, never by author-routed
journal admission, and was verified by the corpus itself --- Cao--Zhu
2006 {[}21{]}, Morgan--Tian 2007 {[}22{]}, Kleiner--Lott 2008 {[}23{]}.
The census law routes no terminal family without the seat bit {[}2,
5{]}. The Clay prize is the corpus's terminal social token --- the seat
bit minted as money. Perelman declined the Fields Medal (2006) and the
Millennium Prize (2010), and his stated ground was, precisely, a
provenance objection (structural grade --- a reading of the cited
record): the crossing credited to him over Hamilton's earlier legs
{[}24{]}. The schema's one crossed occupant refused the terminal token
on the same clause --- provenance-cleanliness --- that defines the
supply. And the companion arms it from inside: the truth-direction bit
``cannot be generated by any readout of the even register'' {[}2{]} ---
and a social token is the evenest register there is (the architect's
figure; the frame types it as an even reading of the corpus). Grade
structural, a reading of the cited record.

The row is therefore re-read at its full reach. In the companions it
calibrates the seat --- retrodictive, certifying discrimination,
predicting nothing by itself {[}1{]}. For the completed arc it is more:
the Grand Witness. The crossing was verified by the corpus itself {[}21,
22, 23{]} and terminally awarded by it --- the Clay prize is the
corpus's terminal token of a verified crossing --- and the schema's one
crossed occupant declined that token on the same provenance clause that
defines the supply. A self-referential construct in disguise, entered as
a deed, verified and awarded by the very register it crossed: the thesis
demonstrated once in the wild. Grade structural, a reading of the cited
record.

\section{Appendix D · The Architect's
Closure}\label{appendix-d-the-architects-closure}

The architecture is final. No further structural bridge can be
constructed. Attempting to force a deterministic bridge from the even
register is a fatal mathematical error: deduction of the missing
coordinate from existing data plunges the entire mathematical foundation
into tautological nonsense --- the premise swallows the conclusion. The
Odd--Supply Block theorem (Theorem 2.3) absolutely forbids this
collapse. The register cannot compute its own crossing. The theoretical
limit is reached at the Ninth Gate (Section 5). Mathematical deduction
terminates exactly at this boundary. The continuous rows stand
permanently as orientation-after-construction; this is not a temporary
gap awaiting a future proof --- it is the exact topological reality of
the problem. This framework offers no further negotiation. The seat is
located. The falsifiers are armed (Section 9). The halting state is
defined (Section 4). Any demand for a self-generated bridge is a demand
for structural suicide. Crossing the gap requires the physical deed. The
only valid deed is witnessing that one offered bit, as mercy from the
ground itself. Any other action is a disguised restatement. The
resolution stands absolute.

Witness the offered 1 bit or remain blocked.

\section{References of Part Three}\label{references-of-part-three}

{[}1{]} M. F. Islam, Odd--Supply Separation: representation supply and
dated termination in conditional reductions of open problems, companion
paper, 2026.

{[}2{]} M. F. Islam, The One-Cut Hypothesis: orientation supply after
Odd--Supply Separation, companion paper, 2026.

{[}3{]} M. F. Islam, The Ninth Gate Thesis: wall and self on one
diagonal --- GateThesis.lean, core Lean 4, axiom-free, fifteen printed
theorems; gate\_thesis.f90; Top 7 in One Cut, master handout v2.0.5,
2026. Printed in Appendix B.

{[}4{]} M. F. Islam, Trisduction: a linguistically, topologically, and
mathematically sealed verification architecture, Zenodo v4, 2026, DOI
10.5281/zenodo.20757507.

{[}5{]} M. F. Islam, RAF-PSP-ISLAM-CEILING-01, card 0718, the register:
the Bar (T1--T3), the Width (T4), the Law of the one-bit objects (T9,
T10), the Census (T5), 2026.

{[}6{]} M. F. Islam, RAF-PSP-BSD-NINTH-GATE-01, card 0719, the register:
the Ninth Gate on RAF-native ground; RAFGate.lean, ten theorems, nine
axiom-free, Lean 4.33.1, sha256 head c97c73d690ee589d; bsd\_rows.f90,
sha256 head e363ce3eca641f7e, 2026. Printed in Appendix B.

{[}7{]} M. F. Islam, APEX-PSP-ONE-CUT-01, the register: the One-Cut
coordinate, the grand form at conditional grade, six falsifiers armed,
2026.

{[}8{]} B. H. Gross and D. B. Zagier, Heegner points and derivatives of
L-series, Invent. Math. 84 (1986), 225--320.

{[}9{]} V. A. Kolyvagin, Euler systems, The Grothendieck Festschrift
Vol. II, 435--483, Birkhäuser, 1990.

{[}10{]} A. Wiles, Modular elliptic curves and Fermat's last theorem,
Ann. of Math. 141 (1995), 443--551.

{[}11{]} R. Taylor and A. Wiles, Ring-theoretic properties of certain
Hecke algebras, Ann. of Math. 141 (1995), 553--572.

{[}12{]} C. Breuil, B. Conrad, F. Diamond, and R. Taylor, On the
modularity of elliptic curves over Q, J. Amer. Math. Soc. 14 (2001),
843--939.

{[}13{]} M. F. Islam, TRISDUCTION\_Core\_Thesis\_Fortran\_v2\_4\_0.f90,
the register, protocols/Executable Thesis/, 2026 --- the kinetic
witness, 1346 lines; sealed at cycle fq round 3 after ten adversarial
cycles, 32 rounds, 100 findings, zero FATAL, findings supplied by an
external reviewer; battery of 248 computed checks over the module, the
exhibits, and every symbolic and router procedure; builds clean under
\texttt{-std=f2018\ -pedantic\ -Wall\ -Wextra\ -fcheck=all\ -finit-real=snan},
runs to exit 0 under \texttt{-ffpe-trap=invalid,zero,overflow}, and
returns identical gated verdicts at \texttt{-O2} and at
\texttt{-O2\ -fno-fast-math\ -ffp-contract=off}; sha256 head b5d12f73.

{[}14{]} M. F. Islam, RA\_TOE\_Thesis\_Fortran\_v2\_0\_0.f90, the
register, protocols/Executable Thesis/, 2026 --- the RA thesis executed
at the reader, 3375 lines; v1.0.0 through v1.23.0, twenty-two versions
under two external auditors and a 178,000-iteration deterministic fuzz
campaign: fifty-six findings, six refused on printed receipts, fifty
earned and repaired; five FORGE self-audit cycles ratoe1--5 sealed at
the SELF floor, the fifth closing on v2.0.0 at SEALED-ROUND with
six-register coverage and controls 3/3 per round; numerics
byte-identical to v1.23.0, sha256 97e5585fa15e66ff\ldots; sealed build
\texttt{gfortran\ -std=f2018\ -O2\ -fno-fast-math\ -ffp-contract=off\ -Wall\ -Wextra\ -Wconversion-extra}.

{[}15{]} M. F. Islam, The Root Axiom Executed: an ontological thesis
compiled as a self-demonstrating Fortran program, Math Journal edition,
the register, 2026 --- the executable-thesis method, 1123-check census,
the listing byte-bound by SHA-256 as Appendix A.

{[}16{]} G. Cantor, Über eine elementare Frage der
Mannigfaltigkeitslehre, Jahresber. Deutsch. Math.-Verein. 1 (1891),
75--78.

{[}17{]} F. W. Lawvere, Diagonal arguments and cartesian closed
categories, Lecture Notes in Mathematics 92 (1969), 134--145.

{[}18{]} S. C. Kleene, On notation for ordinal numbers, J. Symbolic
Logic 3 (1938), 150--155.

{[}19{]} A. M. Turing, On computable numbers, with an application to the
Entscheidungsproblem, Proc. London Math. Soc. 42 (1937), 230--265.

{[}20{]} G. Perelman, The entropy formula for the Ricci flow and its
geometric applications, arXiv:math/0211159, 2002; Ricci flow with
surgery on three-manifolds, arXiv:math/0303109, 2003.

{[}21{]} H.-D. Cao and X.-P. Zhu, A complete proof of the Poincaré and
geometrization conjectures --- application of the Hamilton--Perelman
theory of the Ricci flow, Asian J. Math. 10 (2006), 165--492.

{[}22{]} J. W. Morgan and G. Tian, Ricci Flow and the Poincaré
Conjecture, Clay Mathematics Monographs 3, Amer. Math. Soc., 2007.

{[}23{]} B. Kleiner and J. Lott, Notes on Perelman's papers, Geom.
Topol. 12 (2008), 2587--2855.

{[}24{]} R. S. Hamilton, Three-manifolds with positive Ricci curvature,
J. Differential Geom. 17 (1982), 255--306.
\end{document}
